% !TeX program = lualatex
\documentclass[12pt,oneside]{book}
\usepackage[utf8]{inputenc}
\usepackage[T1]{fontenc}
\usepackage{geometry}
\geometry{margin=1in}
\usepackage{amsmath,amssymb}
\usepackage{graphicx}
\usepackage{xcolor}
\usepackage{tcolorbox}
\tcbuselibrary{skins,breakable}

% Suit symbols – used only in imported generator files, kept for compatibility
\newcommand{\spadesuit}{♠}
\newcommand{\heartsuit}{♥}
\newcommand{\clubsuit}{♣}
\newcommand{\diamondsuit}{♦}

% Almanac environment
\newenvironment{almanac}
  {\begin{quote}\small\itshape}
  {\end{quote}}

\begin{document}

% ----------------------------------------------------------------------
% FRONT MATTER
% ----------------------------------------------------------------------
\frontmatter

\title{\Huge Fate's Edge\\[0.3\baselineskip]
\Large The Way of Silk}
\author{Sarnai, Knot‑Tyer’s Hand \\ \small (apprentice to Caravan Judge Sarnai of the Nine Knots)}
\date{\today}
\maketitle

\tableofcontents

% ----------------------------------------------------------------------
% PREFACE – Framing device (Tulkani apprentice)
% ----------------------------------------------------------------------
\chapter*{Preface: The Knot‑Tyer’s Ledger}
\addcontentsline{toc}{chapter}{Preface}

\noindent
{\Large
\textit{My grandmother tied a knot for every caravan she led. When she died, her cord had nine hundred knots. I inherited the cord and the debt. She had promised to write this book. Now I must.}
}

The road is not a line on a map. It is a cord, and every traveller ties a knot. The Fhara tie knots in water – they call them wells. The Sidhi tie knots in light – they call them lighthouses. The Pereshi tie knots in fire – they call them coals. The Kuvani tie knots in grass – they call them remount stations. And we, the Tulkani, tie knots in memory – we call them stories.

I am Sarnai, Knot‑Tyer’s Hand. My grandmother was Caravan Judge of the Nine Knots. She travelled the Way of Silk from the Brass Coast to the eastern isles, from the Violet Steppe to the Crimson Valley. She promised to write a guide for those who would follow – a ledger of customs, strings, and tolls. She died before she could finish.

This book is the interest on her promise.

You will find no card‑drawing oracles here. The road does not speak in suits; it speaks in etiquette, in water shared, in lamps lit at dusk, in the ninth cup left unpoured. I have replaced the cartomancer’s compass with a caravan‑master’s toolkit: season wheels, mount tags, score types, and the cold arithmetic of standing and feud.

The peoples of the Way are not exotic. They are neighbours. They have their own laws, their own debts, and their own ways of remembering. The Fhara weigh water before silver. The Sidhi keep green flames burning for the dead. The Pereshi erase names they will not speak. The Kuvani read the sky in horse‑hoof echoes. And the Tulkani – we tie knots. When you cross a threshold, knock eight times. The ninth is for the Hollow.

I have tied a knot for every story in this book. The cord is heavy. The debt is not yet paid.

\vspace{1em}
\noindent\textit{— Sarnai, at a caravanserai whose name she is still learning to pronounce}

% ----------------------------------------------------------------------
% HOW TO USE THIS BOOK
% ----------------------------------------------------------------------
\chapter*{How to Use This Book}
\addcontentsline{toc}{chapter}{How to Use This Book}

This book is a companion to the \textit{Fate’s Edge Core Amaranthine Worldbook} and the \textit{System Reference Document (SRD)}. It provides cultural and procedural content for the eastern caravan routes – the Way of Silk – using a **travel‑and‑trade engine** instead of the card‑based Cartomancy chapter.

\subsection*{What You Need}
\begin{itemize}
  \item The \textit{Fate’s Edge SRD or Player/GM Guides} for core rules (attributes, skills, actions, Position/Effect, Clocks, Boons, Story Beats).
  \item A **Caravan/Tribe Sheet** (see Appendix) to track Standing, Feud, Exposure, Oath Ledger, and Notables.
  \item **Mount/Camel Tags** for desert, steppe, and highland travel.
  \item The **Season Wheel** to guide the campaign’s rhythm.
\end{itemize}

\subsection*{Regional Chapters}
Each chapter covers one of the five peoples of the Way:
\begin{itemize}
  \item An introductory paragraph (in Sarnai’s voice, biased but affectionate).
  \item A local quote.
  ítem An \emph{almanac} – a quick Q\&A about daily life.
  \item \textbf{Etiquette Hooks} – cultural actions that grant Position or Effect.
  \item \textbf{Strings} – tokens that act as leverage (water‑share tablets, freedman’s seals, fire‑coals, etc.).
  \item \textbf{Venue Tags} – modifiers for common locations (well‑courts, lighthouse courts, fire‑temples, remount posts, wagon circles).
  \item Then the full generator tables (Spades/Hearts/Clubs/Diamonds) imported from separate files.
\end{itemize}

\subsection*{Appendix: Caravan \& Tribe Rules}
This appendix distils the procedures from the \textit{Sands of Moon and Brass} expansion: Caravan/Tribe Sheets, Mount Tags, Season Wheel, Score Types (Desert Caravan, Oasis Siege, Council Moot, etc.), Weather/Heat Matrix, and a selection of generators (d66 targets, d12 hazards, council cases, verses & boons).

\vspace{1em}
\noindent\textit{“The road is a cord. Every knot is a promise. Untie it, and you pay the debt.”} – Sarnai’s grandmother, the first Caravan Judge

% ----------------------------------------------------------------------
% REGIONAL CHAPTERS
% ----------------------------------------------------------------------
\mainmatter

\chapter{Gateways \& Travel}
\emph{The road is a cord. Every knot is a promise. Untie it, and you pay the toll.}

\section*{Travel Modes of the Way of Silk}

\subsection*{Caravan Routes}
The Way of Silk is a network of seasonal caravan tracks connecting Fhara, Sidhi, Pereshi, Kuvani, and Tulkani territories. Travel is governed by the \textbf{Season Wheel}: spring for caravans, summer for heat survival, autumn for trade, winter for councils.

\subsection*{Desert Legs (Fhara)}
Oasis courts and well‑spirits. A \textbf{Water‑Share Tablet} or a \textbf{Drowned Coin} buys water. Never refuse a date offered by a well‑keeper – it is an oath of hospitality.

\subsection*{Coast Roads (Sidhi)}
Lighthouse courts and freedman’s seals. A \textbf{Lighthouse Charter} grants safe harbour. Always honour the \textbf{Bread \& Salt} ritual before negotiation.

\subsection*{Highland Roads (Pereshi)}
Fire temples and name‑erasure. A \textbf{Fire‑Shrine Token} or a \textbf{Rose from the Ruined Garden} pacifies the mountain spirits. Never step on a fallen book – written words carry the weight of law.

\subsection*{Steppe Legs (Kuvani)}
Remount stations and sky‑road songs. A \textbf{Remount Chain Mark} or a \textbf{Grazing Compact} buys fresh horses. Never ride a Kuvani horse without asking its name – it will throw you.

\subsection*{Wagon Circles (Tulkani)}
Itinerant courts and knot‑memory. A \textbf{Memory‑Bead} or a \textbf{Safe‑Haven Oath} guarantees sanctuary. Never sleep outside the circle at night – the Hollow walks.

\section*{Gateways to Other Expansions}

\begin{itemize}
  \item \textbf{Amaranthine Core (Kahfagia, Ecktoria, etc.)}: From Fhara, the \textbf{Salt Road} leads north to Oshiira. From Kuvani, the \textbf{Eastern Steppe Road} connects to Ykrul.
  \item \textbf{Southern Realms (Oshiira, Ashaan, Ngomebe)}: From Sidhi ports, sail to the Brass Coast (Dhahara) or to Oshiira. From Pereshi, the high passes lead to Ashaan.
  \item \textbf{Eastern Realms (Sihai, Nihori, Ayokha)}: From Kuvani, the \textbf{Silk Road} east leads to Sihai. From Tulkani, the \textbf{Processional Roads} connect to Nihori’s shrine circuits.
  \item \textbf{Dhahara (Brass Coast, Nine Kingdoms, Himadri)}: From Sidhi, the \textbf{Brass Coast} is a direct sea route. From Kuvani, the \textbf{Himadri passes} lead to the iron mines.
\end{itemize}
\newpage
% -------------------- FHARA --------------------
\chapter{Fhara}
\emph{The Oasis Courts.}

\noindent
Sarnai says: \textit{“The Fhara measure water before they measure silver. Their wells have ears. Their coffee has three cups – the fourth is for the Hollow. I tied a knot for every well I drank from. My cord is heavy.”}

\noindent
The Fhara dwell on a sun‑baked peninsula where scattered oases and dry riverbeds connect the coast to the interior. Their cities are built around wells and date groves; their caravans are the arteries of the spice trade. Water is law – the well‑keepers are judges, and the courthouse is the shade of a palm. Coffee is not a drink; it is a ritual. Three cups: welcome, negotiation, contract. Refuse the third, and you have broken faith.

\vspace{0.5em}
\noindent\textit{“Never refuse a date offered by a well‑keeper. The fruit is sacred to their well‑spirits. Refusing it ends the conversation – and your access to water.”} \\
— Well‑Judge Fatima al‑Wadi

\vspace{1em}
\begin{almanac}
\textbf{What do people eat for breakfast?} Dates and thin coffee – the dates from the well‑keepers’ groves, the coffee brewed with cardamom and camel milk. The milk is sour. The sourness is for endurance.

\textbf{What does a child learn before they can walk?} To count cups. The eight cups of negotiation. The ninth cup is for the well‑spirit. The well‑spirit is always thirsty.

\textbf{What is the first thing you notice about a stranger?} Their water‑skin. Full means not travelled far. Dry means travelled too far. No water‑skin means either a well‑keeper or a fool.

\textbf{What is the most common crime?} Water theft – a declaration of feud. The well‑spirit records every theft in the aquifer. The verdict is thirst.

\textbf{What does no one talk about but everyone knows?} The wells are not natural. The first Fhara carved them and bargained with the sand: water for memory. The Fhara remember everything. They cannot forget. That is the curse.

\textbf{Where do people go when they die?} To the oasis. Buried in date groves, marked with a single stone. If the name is remembered, the dead rest. If forgotten, they walk.

\textbf{How do you apologize for a serious wrong?} You offer the ninth cup, poured from your own water‑skin. Placed at the well’s edge. If the spirit drinks, you are forgiven. If the water evaporates, you are not.

\textbf{What is the most important festival?} The Filling of the Wells (spring). Every well cleaned and blessed. The water‑clerks recite the names of the year’s debtors. If their reflection wavers, they are forgiven. If it holds, they must pay.

\textbf{What is the secret that could destroy a powerful person?} The water‑clerks’ ledger is a forgery. The original was lost in a sandstorm three generations ago. The clerks have been inventing water rights ever since. The wells remember the original. They are waiting.
\end{almanac}

\subsection*{Etiquette Hooks}
\begin{itemize}
  \item \textbf{Water First}: Offer water before any business → +1 Position on Petition or Broker in Fhara courts.
  \item \textbf{Three Cups, One Contract}: Complete the three‑cup ritual (welcome, negotiation, contract) before a deal → +1 Effect on the final bargain.
  \item \textbf{Date Offering}: Accept a date from a well‑keeper → gain a temporary String (Well‑Spirit’s Favour) usable once in that oasis.
  \item \textbf{Faux Pas (SB)}: Refusing a date lets the GM start \textit{Thirst Clock} [4].
\end{itemize}

\subsection*{Strings (sample)}
\begin{itemize}
  \item \textbf{Water‑Share Tablet} – proof of contribution; reroll a failed Survival test when finding water.
  \item \textbf{Coffee Cup (Royal Seal)} – +1 Position in Fhara courts.
  \item \textbf{Drowned Coin of Qalat al‑Madina} – safe passage from any well‑spirit once.
\end{itemize}

\subsection*{Venue Tags}
\begin{itemize}
  \item \textbf{Well‑Court}: Neutral ground; Lamp‑Oaths bind as durable Strings; breaking ticks Exposure +1.
  \item \textbf{Coffee Bazaar}: Market starts at +1 Position; on a 1, \textit{Lamplighter Gossip} flips an Audience tag.
\end{itemize}

\newpage
\section{Fhara --- ``The Oasis Courts''}
\label{chap:fhara}

\subsection*{Quick Reference}
\textbf{Tagline:} \textit{Water is a gift from the ancestors. The well remembers every cup.}

\begin{quote}
  \textbf{Well-Judge (Elite):}  
  ``Water is not a right. It is a gift from the ancestors, passed down through the sand. The oasis remembers every drop drawn, every cup poured, every promise made at its edge. I do not judge the living. I witness the balance.''
\end{quote}

\begin{quote}
  \textbf{Water-Clerk (Commoner):}  
  ``My grandfather dug this well with a shovel and a prayer. My mother sealed the walls with clay from the dry river. I sit in the shade and measure every bucket. The sand does not forgive a broken trust. But it does forgive a forgetful mind – if you return to share.''
\end{quote}

\textbf{Genre / Mood:} Desert law, coffee ritual, water-as-community.

\textbf{Starting Location:} The shaded stone bench of a well-court, where the air is cool and the water is sweet. A Well-Judge dips a clay cup into the basin: ``You have come for water, or for trade, or for a story. First, drink. Then tell me which promise you intend to keep today.''

\vspace{2mm}
\begin{tcolorbox}[colback=black!3,colframe=blue!60!black,title={Theme \& Atmosphere}]
The Fhara dwell on a great, sun-baked peninsula where scattered oases and dry riverbeds connect the coast to the interior. Their cities are built around wells and date groves, and their caravans are the arteries of the spice trade. Water is law – the well-keepers are judges, and the courthouse is the shade of a palm. The Fhara are merchants, poets, and negotiators; they will bargain for an hour over a cup of coffee, and a day over a camel. Their honour is measured in generosity, and their memory is measured in water-shares. Coffee is not a drink; it is a ritual. Three cups: the first is welcome, the second is negotiation, the third is a contract. Refuse the third, and you have broken faith.

\vspace{2mm}
\textbf{What you notice first:} The smell of roasting coffee and overripe dates. The way every well has a stone bench for the judge. The silence when the sun passes directly overhead – no one moves, no one speaks.

\textbf{The one rule that kills newcomers:} Never refuse a date offered by a well-keeper. The fruit is sacred to their well-spirits. Refusing it ends the conversation – and your access to water.
\end{tcolorbox}

\subsection*{Regional Mood: The Weight of Water}

Power here is measured in well-depth, date groves, and the length of a coffee ritual. Each oasis is a sovereign court, ruled by a Sheikh and a Well-Judge. The Fhara have no single king; their confederacy is a web of marriage alliances, water-sharing treaties, and caravan tariffs. The desert is not empty – it is a record of dry wells, abandoned caravanserais, and the bones of those who could not share. A Fhara merchant can cross the waste with only a waterskin and a coffee pot, because every well has a keeper, every keeper a price, and every price can be negotiated – but the negotiation is a ritual of reciprocity, not of obligation.

\subsection*{The Antagonists: The Nerethi}
\label{subsec:nerethi}

From the deep desert comes the \textbf{Nerethi}. They are acentric, psionic, and distrusted by all. They do not believe in debt; to them, the concept is alien. They do not owe, they do not remember, and they do not forgive. They walk the dunes like shadows, their minds linked in a hive that rejects the Fhara's water-ledger.

\begin{itemize}
  \item \textbf{Philosophy:} Memory is a chain. Debt is a shackle. The Nerethi burn both.
  \item \textbf{Mechanic:} Nerethi agents can erase a \textbf{Bond} or \textbf{Obligation} from a character sheet, but the character loses the memory associated with it permanently.
  \item \textbf{Conflict:} The Fhara see them as thieves of history. The Nerethi see the Fhara as slaves to the past.
\end{itemize}

\subsection*{The Lost Capital: Qalat al-Madina (The Sunken City)}
\label{subsec:qalat}

Before the Kuvani and the Eastern Ykrul swept from the steppe, the Fhara had a single great city – Qalat al-Madina, the Mother of Wells. It stood where the peninsula's largest aquifer surfaced, a ring of gardens and fountains, its walls white as salt. Merchants from Kahfagia and Thepyrgos paid the city's Water-Tithe to cross the waste. Poets sang of its nine hundred wells, each with a different taste.

But Qalat al-Madina's true secret was older than the city itself. Beneath its deepest well, a faction of scholars unearthed a vault of \textbf{forbidden knowledge} – fragments of a civilisation that had turned to ash twenty thousand years before the first Utaran legion marched. The knowledge was not magic; it was a \textit{reckoning}. It spoke of a fire that could scour the world clean.

The three Omens – \textbf{Inaea, Ikasha, and Isoka} – appeared to the Eastern Ykrul and the Kuvani as whispers in the grass and cracks in the sky. They gave a single command: \textit{Bury it. Seal it. If it wakes, the world burns.} The horse-lords rode against Qalat al-Madina, not for plunder, but to honour the Omens' warning. They poisoned the wells, drowned the city in its own aquifer, and salted the earth. The curse they carved into the central well was not cruelty – it was a lock: \textit{``Let no Fhara ever drink here without paying a death.''} The death was the knowledge, buried alive.

Now Qalat al-Madina is a sinkhole field, half-submerged in brackish water, its towers leaning into the mud. The ghosts of the drowned well-keepers rise at dusk, still trying to measure water that no longer exists – and to guard the vault that no one has found. The Kuvani and Ykrul have long since left, but the \textbf{Omen-Whisper} remains: in certain wells, at certain tides, a traveller may hear a voice speaking a single word: \textit{``Remember.''}

\paragraph{What remains:}
\begin{itemize}
  \item \textbf{The Sunken Archive} – a flooded chamber where clay tablets still float, their judgments still legally binding on the living.
  \item \textbf{The Nine Hundred Wells} – now sinkholes, some connected to the underworld.
  \item \textbf{The Cursed Spring} – a pool that rises every new moon; those who drink from it forget their names for a year and a day.
  \item \textbf{The Drowned Coin} – a silver piece minted in the city's last year; it is said to buy safe passage from any Fhara well-spirit, because it still bears the Mother of Well's seal – and the memory of the Omens' warning.
\end{itemize}

\paragraph{Adventure Hooks:}
\begin{itemize}
  \item A Fhara sheikh wants to recover a water-tablet from the Sunken Archive. The tablet proves his clan never ceded the oasis – but the archive is haunted by the drowned judges who whisper of the forbidden vault.
  \item The Kuvani are sending an expedition to the ruins to \textit{renew} the curse. The Fhara need someone to sabotage them – or to learn why the Omens' seal is cracking.
  \item A ghost from Qalat al-Madina appears at a well-court, demanding justice for a murder committed three centuries ago. The accused is a modern Fhara noble, but the ghost's testimony is legally binding under ancient water law. The ghost also whispers a fragment of the forbidden knowledge: a single number.
  \item \textbf{Nerethi Incursion:} A Nerethi psion offers to erase the curse on Qalat al-Madina. The price: the party must forget why they came to the desert in the first place.
\end{itemize}

\begin{tcolorbox}[colback=black!2,colframe=red!70!black,title={The Last Scribe of Qalat al-Madina}]
  \textit{``I wrote what I should not have. Then I ran.''}
  
  The night the Kuvani and Eastern Ykrul rode against Qalat al-Madina, a young water-clerk named \textbf{Samir al-Qalat} sat in the deep archive, copying ancient tablets by lamplight. He was not a hero. He was the grandson of a well-keeper, and he had a bad habit of reading the tablets he was supposed to copy.
  
  That evening, he unrolled a scroll sealed with a symbol he did not recognise – three interlocking rings, each a different shade of green. The ink was still wet. The scroll had been written that morning, by a hand that no longer existed. It was the \textbf{Overseer's fragment}, smuggled out of a vault that should have stayed sealed.
  
  Samir read eight lines. Then he stopped. He knew that if he read the ninth, he would never stop.
  
  When the first torches appeared at the city's edge, he wrapped the scroll in oilcloth, stuffed it into a clay water-jar, and ran. Behind him, the wells were poisoned. Ahead, the desert. He walked for nine days, drinking only the moisture he could lick from the clay.
  
  On the ninth day, he stumbled into a Tulkani caravan circle. The Caravan Judge, an old woman with a cord that dragged on the ground, looked at the jar and said, \textit{``You have brought the ninth cup. We do not pour it. But we can carry it.''}
  
  Samir gave the scroll to the Tulkani. He did not ask for payment. He asked only that they keep it moving – that no single camp hold it for more than nine nights, no single wagon keep it sealed for more than a season. \textit{``The record is warm,''} he said. \textit{``The ink is still wet. The Hollow reads over my shoulder.''}
  
  The Tulkani honoured his request. The fragment passed from wagon to wagon, from story to story, until it reached the Archivolt of Thepyrgos, where it was sealed in a lead-lined box, encased in \textit{karadasarythm}, with a lock that requires eight witnesses to open. The ninth witness, as the \textit{Words of the Ninth} say, is always the Hollow – and the Hollow is always available.
  
  Samir al-Qalat died in the desert three years later, alone except for a single date-palm that grew from a seed he had carried from his mother's grove. The tree still stands. Travellers who rest in its shade sometimes hear a whisper: \textit{``Count to eight. Then stop.''}
  
  \vspace{2mm}
  \textbf{GM Hook:} The Tulkani still carry the memory of Samir's scroll. A Caravan Judge may offer the party a single glimpse of a copy – but only if they promise to carry a message to the Sunken City: \textit{``The ink is still wet. Has the Overseer finished writing?''}
  \end{tcolorbox}

\subsection*{Prominent NPCs of the Fhara}
\label{sec:fhara-npcs}

\begin{longtable}{|p{0.2\textwidth}|p{0.25\textwidth}|p{0.3\textwidth}|p{0.15\textwidth}|}
\hline
\textbf{Name} & \textbf{Role/Title} & \textbf{Motivation} & \textbf{Rumoured Quirk} \\
\hline
Sheikh Mansur & Oasis lord of the Three Wells & Keep his well's water rights against a rival clan & He has not tasted water in ten years – only coffee. \\
\hline
Fatima al-Wadi & Well-Judge of the central route & Settle disputes without bloodshed & She carries a dipper made from a date-palm leaf; it rots after each case. \\
\hline
Rashid & Caravan master of the Spice Road & Find a new route past a dried well & He has three wives, each from a different oasis. The wives do not speak. \\
\hline
Layla & Coffee-house poet & Sell verses for coin and secrets & Her poems are in code; the code is the price of coffee. \\
\hline
The Drowned Well-Keeper & A ghost of a well that went dry & Warn travellers of poisoned springs & She appears as a column of steam at dusk. She points at safe wells. \\
\hline
Harith al-Qalat & Last living survivor of Qalat al-Madina (immortal? cursed?) & Lift the curse on his city – or die trying & He drinks only rainwater. He has not aged in three centuries. He speaks Kuvani with a Fhara accent. \\
\hline
\textbf{Nerethi Whisperer} & \textbf{Psionic Agent of the Deep} & \textbf{Erase all debts and memories} & \textbf{Has no face; wears the skin of those they have erased.} \\
\hline
\end{longtable}

\subsection*{Names of the Fhara}
\label{sec:fhara-names}

Inspired by the desert peninsula: soft consonants, vowel-heavy, with clan prefixes (bin, bint, al-). Surnames often denote a well, a grove, or a trade.

\begin{longtable}{|p{0.25\textwidth}|p{0.25\textwidth}|p{0.2\textwidth}|p{0.2\textwidth}|}
\hline
\textbf{First Names (Male)} & \textbf{First Names (Female)} & \textbf{Surnames} & \textbf{Epithets} \\
\hline
Mansur & Fatima & al-Wadi & \textit{the Well-Born} \\
\hline
Rashid & Layla & bin Hassan & \textit{Coffee-Tongue} \\
\hline
Tariq & Amira & al-Sahel & \textit{the Ninth Cup} \\
\hline
Jamal & Nadia & Qasimi & \textit{Date-Father} \\
\hline
Hakim & Samira & al-Nur & \textit{the Dry-Handed} \\
\hline
Harith & Zubayda & al-Madina & \textit{the Drowned} \\
\hline
\end{longtable}

\paragraph*{(Well/Court/Bazaar)} A well-court with a stone bench and a judge's dipper; a coffee bazaar under canvas, roasting beans; a dried riverbed where old caravans are buried.

\section*{Spades — Places}
\label{sec:fhara-places}
\begin{enumerate}
\setcounter{enumi}{1}
\item \textbf{Well-Court} – A stone bench under a palm; a judge with a clay dipper.
  \hfill \textit{The dipper has a chip. The chip is from a sword that broke a well – and a trust.} \texttt{[JUSTICE]}
\item \textbf{Coffee Bazaar} – Stalls under canvas; the smell of roasting beans.
  \hfill \textit{The best coffee is served in a chipped cup. The chip is a signature.} \texttt{[TRADE]}
\item \textbf{Date Grove} – Rows of palms; the ground soft with fallen fruit.
  \hfill \textit{The fruit ferments at night. The drunk monkeys are the guards. They accept bribes of honey.} \texttt{[COMMUNITY]}
\item \textbf{Dry Riverbed (Old Caravan Route)} – A wide, sandy channel with bleached bones.
  \hfill \textit{The bones are not human. They are camels that died of thirst. Their ghosts still groan a warning.} \texttt{[TRAVEL]}
\item \textbf{Tent Court (Mobile)} – A carpet under a canvas roof; a judge with a scale.
  \hfill \textit{The scale weighs water, not silver. One side is always empty – for the next dispute.} \texttt{[JUSTICE]}
\item \textbf{Caravanserai of the Three Wells} – A walled courtyard with three different springs.
  \hfill \textit{The sweet well is for guests. The brackish well is for animals. The bitter well is for prisoners.} \texttt{[TRAVEL]}
\item \textbf{Salt Pan (Western Edge)} – A white flat where the caravans load salt.
  \hfill \textit{The salt is mixed with ancient tears. It tastes of old grief. Cooks use it sparingly.} \texttt{[TRADE]}
\item \textbf{Lighthouse of the Drowned} – A stone tower on the coast, never lit.
  \hfill \textit{The keeper is blind. He rings a bell when he hears ships. The ships are ghosts.} \texttt{[DEATH]}
\item \textbf{Oasis of the Ninth Palm} – A legend: a grove with nine palms, but the ninth is invisible.
  \hfill \textit{Drink from the ninth well, and you will forget your name for a day. Your friends will remind you.} \texttt{[LEGEND]}
\item[J] \textbf{Abandoned Well-Keeper's Hut} – A collapsed stone shelter with a dry basin.
  \hfill \textit{The basin is carved with a name. The name is the keeper's. He died of thirst – but the well is full.} \texttt{[DEATH]}
\item[Q] \textbf{Coffee Warehouse} – A cool, dark building with sacks of beans from a dozen oases.
  \hfill \textit{Each sack is stamped with a well's mark. The marks are family crests. The beans gossip.} \texttt{[TRADE]}
\item[K] \textbf{Well-Judge's Archive} – A cave filled with clay tablets. Each tablet is a water judgment.
  \hfill \textit{The tablets are baked hard. They are also porous. They weep when a judgment is reversed.} \texttt{[JUSTICE]}
\item[A] \textbf{The First Well} – A mythic spring at the centre of the peninsula, said to never dry.
  \hfill \textit{No one has found it. Those who claim they have are always lying – or dead. The dead are not lying.} \texttt{[LEGEND]}
\item[A+] \textbf{Qalat al-Madina (Sunken City)} – A field of sinkholes and half-submerged towers.
  \hfill \textit{The water is brackish. The ghosts are polite. They offer you a cup. The cup is empty. Drink it anyway.} \texttt{[RUIN]} \texttt{[DEATH]}
\end{enumerate}

\paragraph*{(Judge/Poet/Master)} Well-Judge with a clay dipper and a rotted palm-leaf; coffee-house poet with a stringed instrument; caravan master with a rope of camel hair.

\section*{Hearts — People \& Factions}
\label{sec:fhara-people}
\begin{enumerate}
\setcounter{enumi}{1}
\item \textbf{Well-Judge} – Grey-bearded, blind in one eye. He tastes the water before each ruling.
  \hfill \textit{``Salt means guilt,'' he says. ``Sweet means innocence. Bitter means appeal – and a cup of coffee.''} \texttt{[JUSTICE]}
\item \textbf{Coffee-House Poet} – A woman with hennaed hands and a silver bell on her ankle.
  \hfill \textit{She recites a verse. The verse is a price. The price is never in coin – it is a favor, a memory, a silence.} \texttt{[SOCIAL]}
\item \textbf{Caravan Master} – Leathery skin, a brass whistle on a thong.
  \hfill \textit{He can find water by the colour of the sand. He has never been wrong. He charges for the lesson.} \texttt{[TRAVEL]}
\item \textbf{Water-Clerk} – A young woman with a clay tablet and a reed pen.
  \hfill \textit{She records every bucket drawn. Her ink is water mixed with soot. A dry tablet means a dry well.} \texttt{[COMMUNITY]}
\item \textbf{Date Merchant} – Fat, cheerful, missing two fingers to a knife fight.
  \hfill \textit{He offers a date to every customer. The date is always from his own grove. He knows its story.} \texttt{[TRADE]}
\item \textbf{Oasis Lord (Sheikh)} – Robed in white, a bronze key on his belt.
  \hfill \textit{The key opens the sluice gate. He never lets it leave his sight. He sleeps with it under his tongue.} \texttt{[AUTHORITY]}
\item \textbf{Well-Spirit (Kulub)} – A shimmer of heat above the water at noon.
  \hfill \textit{It asks: ``Who drew without sharing?'' Answer truthfully, or the well will dry. The sharing is a story.} \texttt{[SPIRIT]}
\item \textbf{Camel Breeder} – A tall, silent man with a staff of acacia wood.
  \hfill \textit{He names each camel after a dead wife. The camels answer to the names. They never forget a slight.} \texttt{[COMMUNITY]}
\item \textbf{Salt Carrier} – A slave (illegal, but present in the deep desert). He wears a collar of braided rope.
  \hfill \textit{He is saving to buy his freedom. The price is a thousand measures of salt. His guild saves with him.} \texttt{[LABOR]}
\item[J] \textbf{The Drowned Well-Keeper's Ghost} – A woman in wet robes, standing at a dry well.
  \hfill \textit{She offers a cup of water. The cup is empty. Drink it anyway. The taste is memory.} \texttt{[DEATH]}
\item[Q] \textbf{Coffee Roaster} – A man with blackened hands and a brass drum.
  \hfill \textit{He roasts beans while reciting the genealogy of each grove. The beans remember. They also judge.} \texttt{[TRADE]}
\item[K] \textbf{The Ninth Palm's Guardian} – A hooded figure who appears at the legendary oasis.
  \hfill \textit{He asks: ``Why do you seek the invisible well?'' The answer determines if you see it. There is no wrong answer – only a price.} \texttt{[MYSTERY]}
\item[A] \textbf{The First Well-Keeper (Myth)} – A skeleton sitting on a stone bench, still holding a dipper.
  \hfill \textit{If you take the dipper, you become the new well-keeper. You cannot leave. The water is sweet. The silence is sweeter.} \texttt{[LEGEND]}
\item[A+] \textbf{Harith al-Qalat} – The immortal survivor of the Sunken City.
  \hfill \textit{He drinks only rainwater. He has not aged in three centuries. He knows where the Drowned Coin lies. He will not say.} \texttt{[LEGEND]} \texttt{[CURSED]}
\item[N] \textbf{Nerethi Mind-Thief} – A faceless figure wrapped in sand-coloured silk.
  \hfill \textit{They offer to take your pain. The price is the memory of who hurt you.} \texttt{[PSIONIC]}
\end{enumerate}

\paragraph*{(Dry/Coffee/Dispute)} A well is found dry – the keeper is dead; a coffee ritual is interrupted – the third cup is not poured; a water dispute erupts between two wells – the judge demands a bribe.

\section*{Clubs — Complications/Threats}
\label{sec:fhara-complications}
\begin{enumerate}
\setcounter{enumi}{1}
\item \textbf{Dry Well} – The oasis keeper is dead. Lose 1 Supply.
  \hfill \textit{The keeper's body is curled around the basin. He died of thirst. The well is full – but the spirit is angry.} \texttt{[DEATH]} \texttt{[SURVIVAL]}
\item \textbf{Coffee Ritual Interrupted} – A stranger refuses the third cup. The host draws a knife.
  \hfill \textit{The stranger is a rival clan member. The insult is a declaration of feud. The coffee is still hot.} \texttt{[SOCIAL]} \texttt{[COMBAT]}
\item \textbf{Water Dispute} – Two well-keepers claim the same spring. The judge demands a story.
  \hfill \textit{The story must be true. The judge will know if you lie. She has heard them all.} \texttt{[JUSTICE]} \texttt{[SOCIAL]}
\item \textbf{Date Blight} – A fungus attacks the groves. Prices triple. Starvation looms.
  \hfill \textit{The blight was spread by a merchant from the coast. The merchant's ship flew no flag. The dates weep.} \texttt{[DISASTER]} \texttt{[TRADE]}
\item \textbf{Sandstorm (Khamsin)} – Red dust obscures the path. Navigation DV +1, mark Fatigue.
  \hfill \textit{The dust carries whispers. The whispers are the names of dead caravaners. They are asking for directions.} \texttt{[WEATHER]}
\item \textbf{Customs Inspection (Forged Seal)} – A local governor demands papers. Yours are false.
  \hfill \textit{The seal is a copy. The forger is in the party. He is sweating. The governor is counting the drops.} \texttt{[AUTHORITY]} \texttt{[TRAP]}
\item \textbf{Elephant Stampede (Rare)} – A herd of stray elephants thunders through camp.
  \hfill \textit{The elephants are from a Dhaharan stable. They are escaped. They are looking for their mahout. He is hiding in your tent.} \texttt{[DISASTER]} \texttt{[COMBAT]}
\item \textbf{Mirage (False Well)} – A shimmer on the horizon. Lose 1 Supply chasing it.
  \hfill \textit{The mirage is a memory. The memory is of a well that dried a century ago. The ghost of the well drinks your water.} \texttt{[WEATHER]} \texttt{[TRAP]}
\item \textbf{Obligation Collector (Tulkani Memory-Cord)} – Someone recognises a knot you carry. They demand a story.
  \hfill \textit{The story must be about a broken promise. The knot tightens as you speak.} \texttt{[SOCIAL]} \texttt{[OBLIGATION]}
\item[J] \textbf{Well-Spirit Offended} – A traveller urinated in the basin. The well refuses water.
  \hfill \textit{The spirit demands a story. The story must be true. The well will remember.} \texttt{[SPIRIT]} \texttt{[SOCIAL]}
\item[Q] \textbf{Coffee House Fire} – A roasting drum explodes. The bazaar is burning.
  \hfill \textit{The fire was set by a rival oasis. The rival wants to monopolise the coffee trade. The smoke smells of cinnamon and spite.} \texttt{[DISASTER]} \texttt{[TRADE]}
\item[K] \textbf{Well-Judge's Verdict Forged} – A clay tablet is found altered. The judge is blamed.
  \hfill \textit{The forger is the judge's own clerk. The clerk owes a favour to the Salt Syndicate. The tablet weeps – but the tears are ink.} \texttt{[JUSTICE]} \texttt{[TRAP]}
\item[A] \textbf{The First Well Appears} – The mythic spring manifests. All who drink from it forget their broken promises.
  \hfill \textit{The forgetting is permanent. The sand shifts. The Hollow is pleased. The creditors are not.} \texttt{[LEGEND]} \texttt{[MAGIC]}
\item[A+] \textbf{Qalat al-Madina's Curse} – The Sunken City's ghosts rise and follow the party.
  \hfill \textit{They are polite. They offer water. The water is brackish. They will not leave until someone pays the Drowned Coin – or speaks the word the Omens whispered.} \texttt{[DEATH]} \texttt{[CURSE]}
\item[N] \textbf{Nerethi Memory Erasure} – A party member wakes up unable to recall a Bond or Obligation.
  \hfill \textit{The gap in their mind hums. The Nerethi are watching.} \texttt{[PSIONIC]} \texttt{[CURSE]}
\end{enumerate}

\paragraph*{(Water-share/Court/Dipper)} Water-share tablet (right to draw from a named well for a season); caravan court writ (neutral arbitration for one dispute); well-judge's dipper (one reroll on a negotiation involving water).

\section*{Diamonds — Rewards/Leverage}
\label{sec:fhara-rewards}
\begin{enumerate}
\setcounter{enumi}{1}
\item \textbf{Water-Share Tablet (Stewardship)} – A clay tile. Right to draw from a named well for one season.
  \hfill \textit{The tablet is a seal of trust. Break it, and the well's water becomes free for all – but the breaker becomes the new steward.} \texttt{[COMMUNITY]}
\item \textbf{Caravan Court Writ} – A palm-leaf document. Demands neutral arbitration for any dispute.
  \hfill \textit{The writ is sealed with coffee grounds. The seal is edible. Eat it to void the writ – and the dispute.} \texttt{[JUSTICE]}
\item \textbf{Well-Judge's Dipper} – A clay cup on a stick. Once, reroll a failed negotiation involving water.
  \hfill \textit{The dipper has a chip. The chip is from a sword that broke a well – and a trust. It knows a lie by taste.} \texttt{[SOCIAL]}
\item \textbf{Coffee Cup (Royal Seal)} – A chipped cup used by a Sheikh. Gain +1 Position in any Fhara court.
  \hfill \textit{The chip is a signature. No one will question it. The cup is still warm.} \texttt{[SOCIAL]}
\item \textbf{Date-Palm Seed} – A single seed from a sacred grove. Plant it to claim water rights.
  \hfill \textit{The seed grows into a palm within a year. The well-spirit will recognise you. The dates are a tithe.} \texttt{[COMMUNITY]}
\item \textbf{Oasis Lord's Key} – A bronze key to a sluice gate. Open any well's flow for one hour.
  \hfill \textit{The key is warm. It glows when a well is being poisoned. It burns the hand of a liar.} \texttt{[AUTHORITY]}
\item \textbf{Coffee Merchant's Record} – A record of every coffee sale for a decade. Contains names of spies.
  \hfill \textit{The record is written in disappearing ink. It reappears when heated. The spies are listed in the margins.} \texttt{[TRADE]} \texttt{[SECRET]}
\item \textbf{Salt Carrier's Rope Collar} – A braided rope that once held a slave. Present it to free a captive.
  \hfill \textit{The rope is old. It smells of sweat and freedom. One use. The freed slave will owe you a story.} \texttt{[COMMUNITY]}
\item \textbf{Well-Spirit's Boon} – A vial of water from a haunted well. Drink it to see the nearest spring.
  \hfill \textit{The water is cold. It tastes of copper. The vision lasts one hour. The spring may be guarded.} \texttt{[MAGIC]}
\item[J] \textbf{The Drowned Keeper's Cup} – An empty clay cup. Offer it to a dry well, and the dead keeper will fill it once.
  \hfill \textit{The cup weeps water. The water is salty. Drink it anyway. The dead keeper will ask for a story in return.} \texttt{[DEATH]}
\item[Q] \textbf{Coffee Roaster's Drum} – A brass drum. Roast beans in it, and the smoke forms words.
  \hfill \textit{The words are the future price of coffee. The price is always a warning. The smoke smells of regret.} \texttt{[MAGIC]} \texttt{[OMEN]}
\item[K] \textbf{The Ninth Palm's Date} – A single fruit from the invisible tree. Eat it to forget a single broken promise.
  \hfill \textit{The date is sweet. The forgetting is permanent. The well-spirit will still remember – but it will not collect.} \texttt{[MAGIC]}
\item[A] \textbf{The First Well's Water} – A flask from the mythic spring. Drink it to clear all obligations to Fhara patrons.
  \hfill \textit{The water tastes of nothing. The spirits will not forgive you for drinking it. But they will forget you owed.} \texttt{[LEGEND]} \texttt{[MAGIC]}
\item[A+] \textbf{Drowned Coin of Qalat al-Madina} – A silver piece minted in the lost capital.
  \hfill \textit{Present it to any Fhara well-spirit, and it will grant safe passage once. The coin was paid to a murderer. It still smells of blood – and of the Omens' whisper.} \texttt{[LEGEND]} \texttt{[OBLIGATION]}
\item[N] \textbf{Nerethi Silence} – A psionic token that blocks all Obligation tracking for one session.
  \hfill \textit{You feel light. You feel empty. You do not remember why you came here.} \texttt{[PSIONIC]} \texttt{[LEVERAGE]}
\end{enumerate}

\section*{Quick use notes}
\label{sec:fhara-quick-use}
\begin{itemize}
\item Draw until you have all four suits: Spade = place, Heart = actor, Club = pressure, Diamond = leverage. Highest rank sets the main Timer (2--5 \(\rightarrow\) 4, 6--10 \(\rightarrow\) 6, J/Q/K \(\rightarrow\) 8, A \(\rightarrow\) 10).
\item Diamonds are codified outcomes (tablets, writs, dippers) that change position rather than call for a roll.
\item If any A appears, echo \textbf{water-omens} – a dry well that weeps, coffee that brews itself, a date that bleeds.
\item If the party visits Qalat al-Madina, treat it as a Class IV Anchor site (see Saikou Ira's compendium). The ghosts are not hostile, but they are very, very patient – and they remember the Omens' warning.
\end{itemize}

\subsection*{Additional Features (Cultural Mechanics)}
\label{sec:fhara-mechanics}

\begin{itemize}
\item \textbf{Three Cups, One Contract:} Fhara negotiations follow a rigid ritual. First cup: welcome and genealogy. Second cup: bargaining and threats disguised as compliments. Third cup: sealing the deal. Refusing the third cup voids the negotiation and insults the host. Accepting the third cup without intending to honour the deal is a capital crime – the liar is tied to a date palm until the sun sets.
\item \textbf{Water-Sharing Tablets:} Water rights are recorded on clay tablets, stored in the Well-Judge's archive. A tablet can be split, traded, or shared. The tablet is the well's memory; break it, and the water becomes unclaimed – but the breaker becomes the new steward.
\item \textbf{Well-Spirits (Kulub):} Each well has a spirit, usually benign but watchful. Offerings of dates, coffee, or a true story keep it calm. A well that has witnessed a murder may turn bitter; a well that has been poisoned may turn sentient. A sentient well is a disaster – it follows those who wronged it.
\item \textbf{The Drowned Coin and the Omen-Whisper:} Silver pieces from Qalat al-Madina still circulate. A Drowned Coin can be used to invoke the Water-Right of the Lost Capital – one free, unquestioned draw from any well in Fhara territory. The well-spirit will remember, and it will ask for a story in return. But if you listen carefully while holding the coin, you may hear the Omens' whisper: \textit{``Remember.''}
\item \textbf{Local Patrons (Syncretic Names):}
  \begin{itemize}
    \item \textbf{The Thirst} – Khemesh. \textit{The hunger beneath the sand; the pressure that drinks wells dry.}
    \item \textbf{The Cup Bearer} – The Witness. \textit{The third cup that never lies; the contract sealed in coffee grounds.}
    \item \textbf{The Oasis Father} – The Traveler. \textit{The road that leads to water; the guide who never drinks.}
    \item \textbf{The Date Mother} – Inaea. \textit{The web of groves; the shared harvest that binds families.}
    \item \textbf{The Sunken King} – Varnek Karn. \textit{The ghost of Qalat al-Madina; the negotiator of secrets that outlast cities.}
  \end{itemize}
\end{itemize}

\subsection*{Lore Echoes (Cross-Expansion Hooks)}
\label{sec:fhara-lore}

\begin{itemize}
  \item \textbf{The Omen-Whisper (Inaea, Ikasha, Isoka):} In the deepest wells of the Sunken City, a traveller may hear a voice speaking a single word: \textit{``Remember.''} The voice is not a ghost – it is a fragment of the three Omens' warning, left behind when they commanded the Ykrul and Kuvani to bury the ancient knowledge. A character who drinks from the Cursed Spring and survives may hear the full whisper – a single sentence in a language older than the empire. The GM may reveal a clue about the forbidden vault.
  \item \textbf{Coffee and the Covenant (Mykkiel):} Fhara coffee houses are neutral ground under Covenant law. A well-judge's verdict cannot be appealed outside the oasis. Some Covenant judges are trained in the three-cup ritual – they can tell a liar by the way they hold the cup.
  \item \textbf{Water Tablets and the Aeler Ledger (Ngomebe):} Aeler engineers have been hired to map the underground aquifers. Their breath-books record water pressure and flow. The Fhara see this as an intrusion – or an opportunity for better wells. A shared aquifer is a shared treaty.
  \item \textbf{Date Groves and the Ukwe (Sekogo):} The Ukwe's roots do not reach the dry peninsula, but the Fhara have their own Speaking Palms – trees that remember every date harvested. A date from a sacred grove can witness an oath as surely as a Lethai-ar cord. The tree will not lie.
  \item \textbf{The First Well and the Hollow (Eastern Realms):} The mythic spring that never dries is said to be the Hollow's original well – before it learned to thirst. Aelaerem hedge-keepers sometimes leave offerings at dried wells, hoping to appease the Hollow with water. The Hollow does not drink. It counts.
  \item \textbf{The Sunken City and the Kuvani (Eastern Ykrul):} The Kuvani and Eastern Ykrul never forgot Qalat al-Madina. Some say they still patrol its ruins to prevent the Fhara from reclaiming it – and to ensure the Omens' seal remains unbroken. Others say the curse was their doing, and they hold the only key to lifting it: a bronze paiza stamped with the city's original water treaty, given to them by Inaea herself.
  \item \textbf{The Nerethi and the Void (Psionics):} The Nerethi claim the Hollow is a lie told by creditors. They seek to erase the Ninth Taboo entirely. If they succeed, the Sunken City will rise, and the debt of the world will be wiped clean – along with its history.
\end{itemize}

\subsection*{Fhara Superstitions (burn for +1 die once)}
\begin{itemize}
  \item \textbf{Bread to the Basin:} Crumble a crust into the well. Ignore the first \emph{Dry Well} penalty. \texttt{[SUPERSTITION]}
  \item \textbf{Knot the Dipper:} Tie a single date-palm fibre around the dipper's handle. Cancel \emph{Coffee Ritual Interrupted} or take --1 die on your next negotiation (your choice). \texttt{[SUPERSTITION]}
  \item \textbf{Name to the Well:} Whisper a dead well-keeper's name into the water. Gain +1 Effect when haggling over water rights this scene, but mark Obligation to the Well-Spirit. \texttt{[SUPERSTITION]}
  \item \textbf{Coin to the Ghost:} Toss a Drowned Coin into a dry well. The ghost of Qalat al-Madina will guide you to water once – and whisper the first syllable of the Omens' warning. \texttt{[SUPERSTITION]}
\end{itemize}

\subsection*{Thematic SB Spend Table}
\label{sec:fhara-sb}

\paragraph{Minor Complications (1 SB)}
\begin{itemize}
\item \textbf{Exposure:} Your actions draw unwanted attention from \textbf{Well-Judges or Water-Clerks}.
\item \textbf{Noise:} Sounds of your actions alert nearby \textbf{coffee-house guards or date merchants}.
\item \textbf{Trace:} Evidence of your passage marks your route for \textbf{caravan masters or well-spirits}.
\item \textbf{Delay:} A brief but meaningful setback costs you \textbf{time or favourable water access}.
\item \textbf{Supply Strain:} Mark +1 segment on a relevant \textbf{resource timer}.
\end{itemize}

\paragraph{Moderate Setbacks (2 SB)}
\begin{itemize}
\item \textbf{Alarm Raised:} \textbf{Local Well-Judge or Oasis Lord} becomes aware and begins responding.
\item \textbf{Position Lost:} You lose advantageous ground/cover/stealth due to \textbf{sandstorm or mirage}.
\item \textbf{Foe Appears:} A \textbf{rival oasis guard or obligation collector} arrives on scene.
\item \textbf{Gear Trouble:} A piece of equipment becomes \textbf{Compromised/Neglected}.
\item \textbf{Lock/Barrier:} A simple obstacle now requires a test to overcome.
\end{itemize}

\paragraph{Serious Trouble (3 SB)}
\begin{itemize}
\item \textbf{Reinforcements:} Additional \textbf{well-guards, coffee-house enforcers, or camel riders} arrive.
\item \textbf{Key Gear Breaks:} A crucial tool/weapon becomes temporarily unusable.
\item \textbf{Major Twist:} The situation fundamentally changes – \textbf{well dries / coffee warehouse burns / water tablet breaks}.
\item \textbf{Rail Tick:} Advance a relevant campaign/front timer by 1 segment.
\item \textbf{Condition Applied:} Mark \textbf{Fatigue 1/Harm 1/Condition} appropriate to fiction.
\end{itemize}

\paragraph{Major Turns (4+ SB)}
\begin{itemize}
\item \textbf{Trap Springs:} A prepared danger activates with full effect.
\item \textbf{Authority Arrival:} \textbf{High Well-Judge, Oasis Sheikh, or First Well-Keeper's ghost} intervenes.
\item \textbf{Scene Shift:} The environment changes dramatically – \textbf{well collapses / oasis floods / date grove burns}.
\item \textbf{Patron Omen:} Divine/arcane forces take notice – \textbf{omen appears / blessing lost / curse manifests}.
\item \textbf{Narrative Pivot:} The story takes an unexpected turn that reframes objectives.
\end{itemize}

\paragraph{Region-Specific SB Options}
\begin{itemize}
\item \textbf{Fhara (Water Law):} Well-spirits ripple the surface, clay tablets glow, a dipper's chip grows larger.
\item \textbf{Fhara (Coffee Ritual):} Beans roast themselves, cups overflow without water, the third cup pours black smoke.
\item \textbf{Fhara (Desert):} Sand shifts to reveal old bones, mirages speak, a dry well whispers a name.
\item \textbf{Fhara (Sunken City):} A drowned tablet floats up from a sinkhole, still legible. Its judgment is from a dead judge. It is still binding – and it mentions the Omens.
\item \textbf{Fhara (Nerethi):} The air hums. A party member forgets why they drew their weapon. A Bond vanishes from the sheet.
\end{itemize}

\subsection*{Pursuit Ladder (Near / Mid / Far)}
\begin{itemize}
\item \textbf{Advance/Withdraw (Wits + Survival):} On a hit, shift one band. On a strong hit, also impose \emph{Sand Blind}.
\item \textbf{Cut the Bargain (Presence + Sway):} On a hit, freeze range for one exchange; if a water-share tablet is in your possession, +1 die.
\item \textbf{Weather the Khamsin (Resolve + Endurance):} On a hit in \emph{Sandstorm}, you pick who stays on their feet; on a miss, the GM does.
\end{itemize}

\subsection*{Boss Seeds (Water \& Sand)}
\begin{itemize}
\item \textbf{The Drought-Maker (Nine Dry Wells)} – \emph{Phases}: Well Poisonings \(\rightarrow\) Orchestrated Famine \(\rightarrow\) Oasis Coup. \emph{Cracks}: find a hidden spring, expose a false tablet, survive a season without water.
\item \textbf{The Coffee-Price (Third Cup Cursed)} – \emph{Phases}: Ritual Interruptions \(\rightarrow\) Procession of Empty Cups \(\rightarrow\) Bargain Baptism of a Merchant. \emph{Cracks}: relic coffee bean, survivor testimony, the third cup turning sweet again.
\item \textbf{The Sunken King (Ghost of Qalat al-Madina)} – \emph{Phases}: Drowned Archive \(\rightarrow\) Curse of the Ninth Well \(\rightarrow\) The Coin's Reckoning. \emph{Cracks}: return the Drowned Coin to his skeletal hand, recite the lost water-treaty from a preserved tablet, or convince a living Fhara sheikh to publicly forgive the Kuvani – and to speak the Omen-Whisper aloud.
\item \textbf{The Nerethi Hive-Mind} – \emph{Phases}: Memory Theft \(\rightarrow\) Psionic Intrusion \(\rightarrow\) Total Erasure. \emph{Cracks}: protect a Bond, recover a lost memory, force the Hive to owe a debt.
\end{itemize}

\begin{tcolorbox}[colback=black!3,colframe=black!40!white,title={Patronage \& Power}]
Among the Fhara, power flows through water, coffee, and the weight of a shared bargain. The Well-Judge's word is law, the Oasis Lord controls the sluice, and the Caravan Master commands the desert. But no one rules alone – every decision is a negotiation, every negotiation a ritual, every ritual a shared trust.

\textbf{For the GM:}  
Power in Fhara society is measured in well-depth, coffee cups, and clay tablets. Patronage often takes the form of water-share rights, court writs, or a judge's dipper – each tied to a specific oasis. To emphasise community and reciprocity rather than obligation:
\begin{itemize}
\item Tie rewards to physical tokens (tablets, dippers, keys) that represent shared access rather than obligation. Breaking a token is breaking a trust, not a contract.
\item Let rival oases issue conflicting water claims – resolving them requires a community feast or a joint coffee ritual, not a payment.
\item Use the well-courts, coffee bazaars, and date groves as arenas for social contests, where a cup of coffee is a shared moment and a dry well is a shared tragedy.
\item Introduce the Omen-Whisper as a recurring mystery. The party may find fragments of the Omens' warning carved into ancient well-stones – each fragment a clue to the forbidden knowledge buried beneath Qalat al-Madina.
\item \textbf{Nerethi Contrast:} Use the Nerethi to challenge the party's attachment to their Bonds and Obligations. Is a debt worth remembering if it causes pain?
\end{itemize}
In Fhara society, the sand forgets nothing – and the well remembers every cup shared.
\end{tcolorbox}

% Index entries
\index{Fhara|see{Oasis Courts}}
\index{Theme generator}
\index{Oasis Courts generator}
\index{Places!Fhara}
\index{People!Fhara}
\index{Complications!Fhara}
\index{Rewards!Fhara}
\index{Well-Judge}
\index{Water-Clerk}
\index{Qalat al-Madina}
\index{Sunken City}
\index{Omen-Whisper}
\index{Inaea}
\index{Ikasha}
\index{Isoka}
\index{Eastern Ykrul}
\index{Kuvani}
\index{Local Patrons!Fhara}
\index{Nerethi}
\subsection*{The Narethi — The Unburied of the Red Waste}
\label{subsec:narethi}
\index{Narethi}\index{Fhara!Narethi}

\emph{``The sun remembers what the sand forgets. We remember both. That is our curse. That is our strength.''}

Deep in the Red Waste, where the Fharan oases give way to cracked salt flats and the bones of failed cities, dwell the Narethi. They are not Fhara — they were never Fhara — but they have lived on the edge of Fhara's desert for so long that the water-clerks have learned to leave offerings at the boundary stones and ask no questions.

The Narethi are felid — cat-people, in the crude speech of lowland merchants — with tufted ears, slit-pupiled eyes that narrow to threads in the sun, and a proud, slow way of moving that suggests they are never quite where you expect them to be. Their fur ranges from the pale cream of dried salt to the deep rust of the Red Waste's iron-rich sands. Their voices carry a harmonic resonance, two pitches at once, as if they are always speaking in duet with a shadow of themselves.

What you notice first: The way they never blink in unison. The scent of myrrh and old copper. The silence when they stop speaking — it is not empty; something is listening.
The one rule that kills newcomers: Never ask a Narethi about the ninth syllable. They will not answer. They will not forget that you asked.

\paragraph{Where They Live}
The Narethi have no cities, no wells, no date groves. They dwell in the \textbf{Sunken Courts} — a network of subterranean chambers carved from ancient salt and basalt, hidden beneath the cracked flats of the eastern Red Waste. The entrance is a sinkhole that smells of ozone and old copper. Visitors who are not Narethi are blindfolded and led through nine switchbacks before the tunnels open into the first Court.

The Courts are warm, lit by phosphorescent crystals that the Narethi claim are the teeth of a god who died before the first Fhara well was dug. The air is still and dry, scented with myrrh, old blood, and the faint, sweet smell of mummification spices.

\paragraph{Origins (What the Narethi Believe)}
The Narethi do not speak of their origins to outsiders. When asked, they say: \textit{``We were the ones who did not burn. The world turned to glass. We did not turn with it. We waited. We are still waiting.''}

Their folklore tells of a \textbf{First Scouring} — a war in the sky before the first human king drew breath. They do not name the combatants. They say only that \textit{three stars argued, one dreafellmed, and one of the fallen's children tried to swallow the sun}. The Narethi chose the wrong side — or the right side, depending on which elder you ask — and were exiled to the wastes as punishment. Their task: to guard a \textbf{sleeping thing} buried beneath the salt, to sing it lullabies, and to ensure it never wakes.

They do not remember the war clearly. Their memories are echoes, passed down in harmonic chants that lose one syllable each generation. The oldest chant is now only eight syllables long. The ninth syllable was the name of the thing they guard.

\paragraph{The Unburied — The Mummified One}
At the heart of the Sunken Courts, in a chamber cooled by no draft and lit by no crystal, lies the \textbf{Unburied}. It is a Narethi — or was, once — wrapped in linen and salt, its fur preserved, its eyes sewn shut with silver wire, its mouth open in a silent scream that has not stopped for millennia.

The Narethi do not worship the Unburied. They attend to it. They say it was the first of them, the one who refused to burn, and that its mind still dreams beneath the salt. Its dreams are psionic — they leak into the Courts, shaping the crystals, influencing the Narethi's thoughts, and granting them their strange mental gifts.

The Unburied does not speak. It \textit{resonates}. Narethi who meditate in its chamber report hearing a single word, repeated endlessly, in a voice that sounds like their own: \textit{``Remember. Remember. Remember.''}

\paragraph{The Narethi's Psionic Gifts}
All Narethi are psions, though they do not use that word. They call their abilities \textit{the Resonance} — an echo of the Unburied's eternal dream. A Narethi's psionics are instinctive, shaped by meditation and the harmonic chants they learn as children.

\textbf{Narethi Psionic Traits (Optional Player Character)}
\begin{itemize}
    \item \textbf{Natural Telepathy:} You may communicate wordlessly with any willing creature within Near range. This is not language — it is emotion, image, and intent.
    \item \textbf{Resonance Sense:} Once per session, you may touch a surface (stone, salt, metal) and ask the GM one yes/no question about an event that occurred within Near range in the past day. The resonance lingers.
    \item \textbf{Disciplined Mind:} You gain +1 die to resist psychic assault or mental intrusion. Your thoughts are your own — and the Unburied's.
    \item \textbf{Night-Eyes:} Narethi eyes reflect light like a cat's. In darkness, you see as if in dim light. This is natural, not magical.
\end{itemize}

\paragraph{The Resonance Price}
The Narethi's psionic gifts are not free. Each time a Narethi uses a significant psionic ability (a Telepathy deep read, a Clairvoyance scry, a Psychic Assault), they mark a segment on their \textbf{Resonance Leash [4]}. When the leash fills, the Unburied's dream bleeds through — the Narethi experiences a waking vision of the First Scouring. The GM describes a fragment: a city of black glass, a sky on fire, a scream that has no end. The Narethi marks 1 Fatigue and the leash resets.

\paragraph{Society and Structure}
The Narethi are ruled by a \textbf{Father of Resonance} — an elder who has spent the most time in meditation with the Unburied. She does not command; she interprets the dream. Her word is law because the dream has foreseen the consequence of disobedience. Below her are \textbf{Chanters} (priests and lore-keepers), \textbf{Watchers} (scouts and warriors who patrol the boundaries of the Red Waste), and \textbf{Silent Ones} (Narethi who have taken a vow of non-speech to better hear the Resonance).

Their society is rigid, hierarchical, and secretive. Outsiders are not trusted. A Narethi who leaves the Courts is considered \textit{unbound} — they may return, but they will be questioned about everything they saw and heard. Some unbound Narethi never return. The Father says the dream does not speak of them.

\textbf{Factions Within the Courts:}
\begin{itemize}
    \item \textbf{The Listeners} — Believe the Unburied is trying to wake, and that the Narethi should help it. They study ancient chants, seeking the lost ninth syllable.
    \item \textbf{The Mutes} — Believe the Unburied should never wake. They have taken vows of absolute silence, hoping that if no one speaks, the dream will fade.
    \item \textbf{The Unbound} — Narethi who have left the Courts and lived among surface-dwellers. They are viewed with suspicion, but they are also the Narethi's only source of news about the outside world.
\end{itemize}

\paragraph{Relationship with the Fhara}
The Fhara know of the Narethi, but they do not speak of them. At the boundary stones that mark the edge of the Red Waste, well-judges leave offerings of salt and dried dates twice a year. The Narethi take them. No one meets.

When a Fhara caravan strays too close to the Sunken Courts, it may return with no memory of the detour, its water-skins refilled, its camels calm. The Narethi do not harm — they erase. A traveler who has been touched by the Resonance may wake with a single new knot in their memory-cord, tied in a pattern no Tulkani recognizes.

The Fhara call the Narethi \textit{the Unremembered}. It is not an insult. It is a warning.

\paragraph{The Secret the Mother Keeps}
The Narethi do not know why they guard the Unburied. They only know that if it wakes, the Red Waste will spread. The sand will turn to glass. The sun will not set for nine days.

What the Father of Resonance knows — and does not tell even the Chanters — is that the Unburied is not a prisoner. It is a \textit{lock}. Beneath it, buried deeper than the Courts, deeper than the salt, deeper than the bones of the First Scouring, something else sleeps. The Unburied's dream is the only thing keeping it asleep. If the Unburied ever stops dreaming, the thing beneath will remember what it was.

The Father does not know what that thing is. She only knows that the dream's single word — \textit{``Remember''} — is not a command to the Narethi. It is a command to the thing below. \textit{Remember what you are. Remember what you did. Remember why you must never wake.}

\clearpage
\subsubsection*{Narethi d6 Mini-Generators}

\paragraph{Narethi Places (Spades) — d6}
\begin{enumerate}
    \item \textbf{The Sunken Courts (Seat of Resonance)} — A subterranean complex of salt-crusted chambers lit by phosphorescent crystals. The crystals hum when the Unburied dreams deeply. Visitors are blindfolded through nine switchbacks. [PSIONIC] [SECRET]
    \item \textbf{The Chamber of the Unburied (Heart)} — A sealed room cooled by no draft, lit by no crystal. The mummified Narethi lies on a bed of salt and linen, eyes sewn shut with silver wire. The air tastes of copper and old grief. [DEATH] [PSIONIC]
    \item \textbf{The Hall of Lost Syllables (Archive)} — Walls carved with harmonic chants. The oldest chant has eight syllables; the ninth space is empty, polished smooth by generations of fingers. The stone is warm. [MEMORY] [RIDDLE]
    \item \textbf{The Resonant Well (Psionic Pool)} — A basin of still, dark water. Touch it, and you hear echoes of the Unburied's dream. Narethi meditate here to attune their Resonance. The water never ripples. [PSIONIC] [OMEN]
    \item \textbf{The Watchers' Perch (Boundary)} — A salt-crusted rock outcropping overlooking the Red Waste. Watchers sit here in silence, scanning for intruders. The perch is marked with nine carved eye symbols — the ninth is half-erased. [WATCH] [STEALTH]
    \item \textbf{The Salt Pan of Forgotten Names (Exile Ground)} — A cracked white flat where unbound Narethi who cannot return are said to walk. Their footprints appear at dawn and vanish by noon. Do not follow them. The salt tastes of tears. [DEATH] [MYSTERY]
\end{enumerate}

\paragraph{Narethi People \& Factions (Hearts) — d6}
\begin{enumerate}
    \item \textbf{Father of Resonance (Ruler)} — An elder with fur gone white, eyes that reflect no light. She speaks only to interpret the dream. Her voice is a whisper that carries through the Courts. She has not left the Chamber in twenty years. [AUTHORITY] [SECRET]
    \item \textbf{Listener Chanter (Priest)} — A Narethi with a shaved head, harmonic voice, and a cord of nine knots. She seeks the lost ninth syllable. She will trade a resonance secret for a story about a dream you cannot forget. [PSIONIC] [LORE]
    \item \textbf{Mute Sentinel (Silent One)} — A warrior who has taken a vow of silence. She communicates with hand signs and a small clay slate. Her silence is so complete that footsteps make no sound. She has not spoken in nine years. [STEALTH] [COMBAT]
    \item \textbf{Unbound Exile (Outsider)} — A Narethi who left the Courts and now works as a caravan scout. She wears a hood to hide her eyes. She will guide travelers past the Red Waste — but she asks in return for news of the Courts. She cannot return. [TRAVEL] [TRAGEDY]
    \item \textbf{The Dream-Echo (Haunted)} — A Narethi who spent too long in the Resonant Well. The Unburied's dream now bleeds into his waking hours. He speaks in two voices, one his own, one a whisper of ancient glass. He is looking for someone to cut the resonance. [PSIONIC] [CURSE]
    \item \textbf{The Keeper of Teeth (Crystal-Minder)} — An old, blind Narethi who tends the phosphorescent crystals. She claims the crystals are growing — that they were not here when she was young. She is afraid of what they will become. [CRAFT] [OMEN]
\end{enumerate}

\paragraph{Narethi Complications (Clubs) — d6}
\begin{enumerate}
    \item \textbf{Resonance Bleed} — The Unburied's dream intensifies. All Narethi within the Courts mark 1 Fatigue. Visitors hear whispers in a language they almost understand. The crystals glow brighter. [PSIONIC] [FEAR]
    \item \textbf{The Ninth Syllable Spoken} — Someone speaks the lost syllable aloud — by accident, or by design. The Unburied's mouth twitches. The ground trembles. Start a Waking Dream [6] timer. [DEATH] [APOCALYPSE]
    \item \textbf{Unbound Hunted} — A Narethi exile is being pursued by Watchers. The exile begs the party for passage across the border. Helping them earns a Narethi's gratitude — and the Mother's attention. [STEALTH] [POLITICS]
    \item \textbf{The Crystals Sing} — The phosphorescent crystals emit a harmonic tone that induces waking visions. Anyone who hears it must test Spirit+Resolve (DV 4) or experience a fragment of the First Scouring. Mark 1 Fatigue on failure. [PSIONIC] [OMEN]
    \item \textbf{Mute's Vow Broken} — A Silent One speaks for the first time in years. She speaks a warning: something is waking beneath the salt. The other Narethi do not believe her. She will not speak again. [TRUTH] [TRAGEDY]
    \item \textbf{The Salt Pan Walks} — Footprints in the Salt Pan of Forgotten Names begin moving toward the Courts. The dead unbound are coming home. They are not hostile — they are curious. That is worse. [DEATH] [FEAR]
\end{enumerate}

\paragraph{Narethi Rewards \& Leverage (Diamonds) — d6}
\begin{enumerate}
    \item \textbf{Resonance Shard (Crystal Fragment)} — A piece of phosphorescent crystal from the Courts. Once per session, touch it to gain +1 die on a psionic roll. The shard then dims for a scene. The Unburied knows you used it. [PSIONIC] [LEVERAGE]
    \item \textbf{Mute's Clay Slate (Witness Token)} — A small clay tablet carried by a Silent One. Present it to any Narethi to demand a single truthful answer. The slate then cracks. The Mute who gave it will not speak of it. [SOCIAL] [TRUTH]
    \item \textbf{Listener's Cord (Nine Knots)} — A knotted cord with eight tight knots and one loose. Untie a knot to ask the GM one yes/no question about the Unburied's dream. The knot reties itself after use — but the loose knot tightens slightly. [MEMORY] [RIDDLE]
    \item \textbf{Watcher's Salt Pouch (Boundary Offering)} — A leather pouch of salt blessed by the Watchers. Scatter it to create a temporary threshold that spirits and psionic projections cannot cross. Lasts one scene. The salt tastes of iron. [WARD] [PROTECTION]
    \item \textbf{Unbound's Map (Hidden Route)} — A scrap of hide marked with charcoal lines. Shows a safe route through the Red Waste that avoids both Fhara patrols and Narethi Watchers. The map was drawn by an exile who never returned. [TRAVEL] [STEALTH]
    \item \textbf{The Dream-Echo's Confession (Whisper)} — A single word whispered by a Dream-Echo before he went silent. The word is a fragment of the ninth syllable. Speaking it aloud to a Narethi grants +1 Position in any negotiation with them — but the Unburied will remember your voice. [LEGEND] [DANGER]
\end{enumerate}

\paragraph{Adventure Hooks (d6 Quick Starts)}
\begin{enumerate}
    \item \textbf{The Lost Unbound.} A Narethi who left the Courts a decade ago has resurfaced — as a mercenary captain in Ecktoria. The Father of Resonance sends the party to retrieve them, or to learn why they refuse to return. The Unbound has a fragment of the ninth syllable tattooed on their fur.
    \item \textbf{The Resonance Bleed.} A Fhara well has begun to taste of copper and ozone. Those who drink from it dream of black glass and a screaming sky. The Narethi say the Unburied is restless. Something has disturbed its slumber — and the thing beneath is stirring.
    \item \textbf{The Silent One's Plea.} A mute Narethi appears at a caravanserai, carrying a clay tablet carved with a single symbol. The symbol is not Narethi — it is older. The party must deliver the tablet to the Archivolt in Thepyrgos, where a Lethai-thora keeper recognizes it as a fragment of the First Scouring. The tablet begins to warm.
    \item \textbf{The Ninth Syllable.} A Tulkani story-singer has begun to hum a chant that is not Tulkani — a chant with nine syllables, not eight. The Narethi Listeners want to hire the party to capture the singer. The Mutes want the singer silenced — permanently. The singer does not know where the chant came from.
    \item \textbf{The Salt Pan Pilgrims.} A group of unbound Narethi exiles is walking toward the Courts from the Salt Pan. They have been dead for years. They are not hostile — they are singing a lullaby. The Father wants to know why they have returned. The party must follow them without being seen.
    \item \textbf{The Crystal's Hunger.} The phosphorescent crystals in the Courts have begun to grow toward the surface. A Watcher reports that they pulse in time with the Unburied's dream. The Keeper of Teeth says they are feeding on something — and that something is almost gone.
\end{itemize}

\paragraph{Player Takeaway}
The Narethi are not a character class or a mechanical package. They are a story — a people shaped by a burden they no longer remember, guarding a secret they cannot name. A Narethi player character is \textit{unbound}: they have left the Courts, for reasons of their own. They may be a spy, an exile, a pilgrim, or simply someone who could not bear the silence.

Work with your GM to determine:
\begin{itemize}
    \item Why did you leave the Sunken Courts?
    \item Do you still feel the Resonance? (If yes, you track a Resonance Leash. If no, you have lost some of your psionic gifts — but you are truly free.)
    \item What do you remember of the Unburied's dream?
    \item Who in the Courts wants you to return — and who wants you dead?
\end{itemize}

The Narethi are few. The Courts are quiet. The dream continues.

\noindent \textit{``The ninth syllable is not a sound. It is the space between heartbeats. The Unburied has been counting for a very long time.''}
\newpage

% -------------------- SIDHI --------------------
\chapter{Sidhi}
\emph{The Diaspora of the Sealed Coast.}

\noindent
Sarnai says: \textit{“The Sidhi carry their homeland in a brass seal. They light green flames to show the drowned the way home. I asked one freedman if he would ever return to Galanina. He touched his collar – a scar, not a seal – and said, ‘I never left.’”}

\noindent
The Sidhi are a people of the coast and the diaspora. Their ancestral homeland is Galanina, a rugged shore of limestone cliffs and ancient ports. The Ashaani magocracy enslaved them for generations, scattering families across the Amaranthine. Those who survived the Long Trail – a forced migration spoken of only in whispers and lullabies – carried nothing but the memory of salt and the names of the dead. Now they run ropeways, spice auctions, and port guilds. Their symbol is the lighthouse lamp: a green flame that shows the faces of the drowned, so that the dead may witness the living.

\vspace{0.5em}
\noindent\textit{“Never ask about ancestry. The Ashaani slavery raids are a fresh wound. Asking which clan someone belongs to is asking how their grandparents were captured, sold, or died.”} \\
— Lighthouse Judge of Galanina

\vspace{1em}
\begin{almanac}
\textbf{What do people eat for breakfast?} Salted fish and flatbread – the fish dried on the docks, the bread baked in communal ovens fuelled by driftwood. The driftwood carries the names of lost ships.

\textbf{What does a child learn before they can walk?} To read the lighthouse codes. The eight lamp‑flashes: Safe, Danger, Storm, Pirate, Customs, Quarantine, Home, and the ninth – the flash that signals a freedman’s return. The ninth flash is never used. It is a promise.

\textbf{What is the first thing you notice about a stranger?} Their seal. Every Sidhi freedman carries a brass seal – the mark of manumission. Worn smooth means free for generations. Sharp means just freed. No seal means not a Sidhi. Not a Sidhi is not a person.

\textbf{What is the most common crime?} Seal‑forging – not a crime, a necessity. The Admiralty demands proof of freedom. The forgers are Sidhi. The Sidhi do not prosecute forgers. They hire them.

\textbf{What does no one talk about but everyone knows?} The brass seals are not brass. They are made from melted chains – the chains that bound the first freedmen. The metal remembers the weight of slavery. The seals are warm. The warmth is the memory of wrists.

\textbf{Where do people go when they die?} To the lighthouse. Ashes scattered from the tower’s top. The wind carries them to the sea. The sea carries them to Galanina. The ancestors are waiting.

\textbf{How do you apologize for a serious wrong?} You offer your seal – your own mark of freedom. It is melted down and reforged into a chain. The chain is worn by the wronged party. When the wrong is forgiven, the chain is melted again. One day, there will be nothing left.

\textbf{What is the most important festival?} The Lighting of the Green Flame (winter solstice). The lighthouse lamps are cleaned and refilled. The green flame shows the faces of the drowned. The ancestors watch. They do not judge.

\textbf{What is the secret that could destroy a powerful person?} The brass seals are not all equal. Some were forged by the Admiralty itself – trap‑seals that mark the wearer as a spy. The spies do not know they are spies. Their reports go to the Admiralty. The Admiralty knows everything.
\end{almanac}

\subsection*{Etiquette Hooks}
\begin{itemize}
  \item \textbf{Bread \& Salt}: Share a simple plate before negotiation → +1 Position in Sidhi homes or guild halls.
  \item \textbf{Name the Line}: Give your family or mentor’s lineage when introduced → Audience: Cool → Warm with civic officials.
  \item \textbf{Guest‑Right Token}: Present a stamped guest‑right from a guild → DV –1 on the first Petition in that venue.
  \item \textbf{Faux Pas (SB)}: Cut Line – ignoring lineage or guest‑right protocol imposes Position –1 for the scene.
\end{itemize}

\subsection*{Strings (sample)}
\begin{itemize}
  \item \textbf{Lighthouse Charter} – right to anchor in any Sidhi port without inspection.
  \item \textbf{Freedman’s Seal} – proof of emancipation; glows when a slaver ship is near.
  \item \textbf{Archive Key} – opens one door in the Sealed Archive of Galanina.
\end{itemize}

\subsection*{Venue Tags}
\begin{itemize}
  \item \textbf{Pilots’ Quay}: navigation disputes settled by local charts; Clue +1 on routes with a token.
  \item \textbf{Lighthouse Court}: tolls and wreck law; Board & Brace DV –1 when acting as sworn rescue.
\end{itemize}

\newpage
\section{Sidhi --- ``The Diaspora of the Sealed Coast''}
\label{chap:sidhi}

\begin{quote}
  \textbf{Lighthouse Judge (Elite):}  
  ``Galanina is not a place on a map. It is a promise that our ancestors made to the sea: we will remember, and we will return. The Covenant gave us courts, but the sea gave us memory. The Admiralty gives us tariffs. We give each other salt and stories.''
\end{quote}

\begin{quote}
  \textbf{Freedman Sailor (Commoner):}  
  ``My grandmother wore a brass collar. My mother wore a freedman's seal. I wear a captain's ring. The slavers took our names, but they could not take the tide. The tide always comes back. So do we.''
\end{quote}

\begin{tcolorbox}[colback=black!3,colframe=blue!60!black,title={Theme \& Atmosphere}]
The Sidhi are a people of the coast and the diaspora. Their ancestral homeland is Galanina, a rugged shore of limestone cliffs and ancient ports far to the west. The Ashaani magocracy enslaved them for generations, but Galanina broke free under the Covenant of Mykkiel, whose judges and notaries stamped the liberation with wax and witness. Now the Sidhi are scattered: some returned to Galanina's lighthouse courts, some trade in Oshiiran ports, and many have found refuge on the Brass Coast of Dhahara, where they run ropeways, spice auctions, and Sidhi-run port guilds. They are fierce, proud, and never forget a debt. Their speech is a creole of Ashaani and coastal tongues, their law is maritime, and their symbol is the lighthouse lamp — a green flame that shows the faces of the drowned, so that the dead may witness the living.

\vspace{2mm}
\textbf{What you notice first:} The smell of salt and old parchment. The way every Sidhi touches a small brass key at their belt — the key to a house they may never see again. The green glow of lighthouse lamps along the shore.

\textbf{The one rule that kills newcomers:} Never ask about ancestry. The Ashaani slavery raids are a fresh wound. Asking which clan someone belongs to is asking how their grandparents were captured, sold, or killed.
\end{tcolorbox}

\subsection*{Regional Mood: The Weight of Memory}

Power here is measured in lighthouse charters, freedman's seals, and the weight of witness. The Sidhi have no single king; their authority rests in the Covenant courts of Galanina and the port-masters' guilds of the diaspora. A Sidhi merchant’s word is her bond, because every bond is witnessed by a lighthouse lamp and recorded in a water-resistant ledger. The Admiralty of Kahfagia administers the Brass Coast, but the Sidhi run the real trade — and they remember every raid, every broken treaty, every ancestor who did not return.

\textbf{Starting Location:} The lamp room of a Sidhi lighthouse, where a green flame burns behind a brass mirror. A Lighthouse Judge turns from the flame: ``You have come for safe harbour, or for a charter, or to ask about the old wounds. First, light a lamp. Then tell me which debt you are here to pay — or to collect.''

\subsection*{Prominent NPCs of the Sidhi}
\label{sec:sidhi-npcs}

\begin{longtable}{|p{0.2\textwidth}|p{0.25\textwidth}|p{0.3\textwidth}|p{0.15\textwidth}|}
\hline
\textbf{Name} & \textbf{Role/Title} & \textbf{Motivation} & \textbf{Rumoured Quirk} \\
\hline
Magistrate Jora & Covenant Judge of the Salt Circuit (Galanina) & Enforce maritime law and protect freedman status & She carries a water-clock that runs backwards when she lies. \\
\hline
Captain Idris & Sidhi lighthouse commander (Brass Coast) & Prevent Ashaani slavers from landing & His left eye is a brass lens; he can read the stars even in fog. \\
\hline
Envoys Naveed & Breath-less agent (hidden) & Recruit exiles to resist the Veiled Sisters & He has three names; each is for a different jurisdiction. \\
\hline
Shipwright Elara & Builder of freedmen's boats & Launch one new vessel each year for escaped slaves & She names each boat after a dead family member. The names are never painted — only carved. \\
\hline
The Brass-Fingered & A smuggler with prosthetic fingers of copper & Run goods past the Admiralty's tariffs & He lost his fingers to a customs trap. He made new ones from the trap's own springs. \\
\hline
\end{longtable}

\subsection*{Names of the Sidhi}
\label{sec:sidhi-names}

Inspired by coastal and diaspora traditions: names with soft consonants, often ending in -el, -on, -ia. Surnames denote a port, a ship, or a freed ancestor.

\begin{longtable}{|p{0.25\textwidth}|p{0.25\textwidth}|p{0.2\textwidth}|p{0.2\textwidth}|}
\hline
\textbf{First Names (Male)} & \textbf{First Names (Female)} & \textbf{Surnames} & \textbf{Epithets} \\
\hline
Idris & Jora & el-Galanina & \textit{the Lamp-Keeper} \\
\hline
Naveed & Elara & Pattinam & \textit{Brass-Fingered} \\
\hline
Samir & Leila & Chettiar & \textit{the Unbound} \\
\hline
Rami & Yasmin & al-Bahr & \textit{Ninth Tide} \\
\hline
Tariq & Dalia & Sidhi-el-Kebir & \textit{the Witness} \\
\hline
\end{longtable}

\paragraph*{(Lighthouse/Court/Harbour)} A lighthouse with a green flame and a blind keeper; a Covenant court on stilts above the tide line; a hidden cove where freedmen's boats are built.

\section*{Spades --- Places}
\label{sec:sidhi-places}
\begin{enumerate}
\setcounter{enumi}{1}
\item \textbf{Lighthouse of the Unfaithful Dead} --- A stone tower with a brass mirror and a green flame.
  \hfill \textit{The light shows the faces of those who died in Ashaani chains. They do not blink.}
\item \textbf{Covenant Court of the Salt Marsh} --- A wooden platform on stilts; the tide rises during trials.
  \hfill \textit{If the water touches the judge's feet, the verdict is ``drowned.'' The accused is not present.}
\item \textbf{Sidhi Hidden Cove (Boatyard)} --- A beach accessible only through a sea-cave.
  \hfill \textit{The cave walls are covered in tally marks — each mark is a slave freed.}
\item \textbf{The Unforgotten Ledger's Vault (Galanina)} --- A subterranean room of sealed jars.
  \hfill \textit{Each jar contains a broken oath. The jars whisper if you listen.}
\item \textbf{Brass Coast Ropeway} --- A cable that carries spice baskets from ship to warehouse.
  \hfill \textit{The rope is braided with Sidhi hair. The dead pull the baskets.}
\item \textbf{Freedman's Square (Port Shed)} --- A marketplace where freed slaves trade without tariffs.
  \hfill \textit{The square has a fountain that flows with salt water. Drinking it reminds you of the crossing.}
\item \textbf{Archive of Sealed Doors (Galanina)} --- A fortress of ledgers, each door locked with a judge's brass ring.
  \hfill \textit{The rings are different sizes. The archivist has a key ring with nine rings. The ninth opens nothing.}
\item \textbf{Lighthouse of the Brass Flame (Dhahara)} --- A copper mirror that turns sunlight into a beam.
  \hfill \textit{The keeper can signal a message to any Sidhi ship within a day's sail. The code is a song.}
\item \textbf{Shipbreaker's Yard (Coast)} --- A beach where old vessels are dismantled.
  \hfill \textit{The workers are convicts. They wear chains that drag in the sand. The chains are melted down for anchors.}
\item[J] \textbf{The Admiral's Palace (Occupied)} --- A former Sidhi fort, now a Kahfagian headquarters.
  \hfill \textit{The flag is flown upside down. The Admiral does not know the difference. The Sidhi do.}
\item[Q] \textbf{The Ninth Well (Galanina)} --- A fresh spring that tastes of iron.
  \hfill \textit{The first Covenant judge drowned a liar there. The water still remembers his name.}
\item[K] \textbf{The Drowned Captain's Ship} --- A spectral dhow seen at the edge of storms.
  \hfill \textit{It offers passage to the dead. The living who board it become its crew.}
\item[A] \textbf{The Sealed Coast (Mythic Galanina)} --- A stretch of shore that appears only at low tide.
  \hfill \textit{Those who walk it see the faces of ancestors who did not escape. The ancestors ask for their names back.}
\end{enumerate}

\paragraph*{(Judge/Keeper/Agent)} Covenant magistrate with a water-clock; Sidhi lighthouse keeper, blind in one eye; Breath-less agent dressed as a fisher, hands uncalloused.

\section*{Hearts --- People \& Factions}
\label{sec:sidhi-people}
\begin{enumerate}
\setcounter{enumi}{1}
\item \textbf{Lighthouse Judge} --- Robed in grey, a water-clock on her belt. She drinks only rainwater.
  \hfill \textit{``The sea does not lie,'' she says. ``Neither should you.''}
\item \textbf{Lighthouse Keeper} --- Blind in one eye, the other sees only green. He lights the lamp with a piece of the sun — or so he claims.
  \hfill \textit{He never leaves the tower. The lamp is his family.}
\item \textbf{Breath-less Agent} --- A fisher with soft hands. He speaks in whispers; his words appear on paper without ink.
  \hfill \textit{He carries a hollow ring. Inside is a single word: ``witness.''}
\item \textbf{Shipwright of the Freedmen} --- A woman with a hammer and a jar of brass nails.
  \hfill \textit{She builds boats for escaped slaves. Each boat is named after a dead family member. The names are never painted — only carved.}
\item \textbf{Spice Warehouse Factor} --- A thin man with a brass seal on a chain. He weighs pepper with a scale that has one arm longer than the other.
  \hfill \textit{The long arm is for bribes. The short arm is for justice. He never remembers which is which.}
\item \textbf{Debt Collector of the Anvil (Aeler)} --- A dwarf authorised to collect on maritime loans. He carries a hammer.
  \hfill \textit{He knocks once for a warning, twice for a summons, thrice for a contract. The fourth knock is a coffin.}
\item \textbf{Freedman's Advocate} --- A woman with a freedman's seal branded on her forearm. She argues cases before the Admiralty.
  \hfill \textit{She has never lost a case. Her clients always confess — after she speaks to them privately.}
\item \textbf{Salt Farmer (Galanina)} --- A woman who harvests salt from coastal pans. Her hands are cracked.
  \hfill \textit{She pays her taxes in salt. The Covenant accepts it as ``liquid testimony.''}
\item \textbf{Smuggler (Brass Coast)} --- A man with brass prosthetic fingers. He runs goods past the Admiralty's tariffs.
  \hfill \textit{His fingers pick any lock. He lost the originals to a customs trap. He made new ones from the trap's springs.}
\item[J] \textbf{The Drowned Priest} --- A ghost in wet robes who walks the shore at low tide.
  \hfill \textit{He asks for a coin to pay the ferryman. Give it, and he will not follow. Refuse, and he will ask for your name instead.}
\item[Q] \textbf{The Black Bishop (Covenant)} --- A judge who rules on maritime disputes. He wears a ring made from a ship's nail.
  \hfill \textit{The ship was the first slaver he freed. The nail still has rust on it.}
\item[K] \textbf{The Admiral's Sidhi Clerk} --- A young woman who writes the Admiral's dispatches. She copies them in secret.
  \hfill \textit{She sends the copies to the Breath-less. Her left hand is missing two fingers — a warning she ignored.}
\item[A] \textbf{The First Freedman (Legend)} --- A slave who bought his freedom with a story. He now sits at the bottom of the Ninth Well.
  \hfill \textit{If you tell him a true story, he will tell you the location of a hidden lighthouse.}
\end{enumerate}

\paragraph*{(Raid/Seal/Leak)} Ashaani slave ships are sighted — the lighthouse keeper asks you to hold the light; a freedman's seal is discovered to be a forgery — the party is accused; a Covenant audit demands to see your cargo — the cargo is escaped slaves.

\section*{Clubs --- Complications/Threats}
\label{sec:sidhi-complications}
\begin{enumerate}
\setcounter{enumi}{1}
\item \textbf{Ashaani Raid (Slaver Ships)} --- Black sails on the horizon. The lighthouse keeper asks you to hold the light.
  \hfill \textit{Hold it too long, and they will see you. Hold it too short, and the freedmen will not see the signal.}
\item \textbf{Forged Freedman's Seal} --- A document you carry is discovered to be a copy. The magistrate calls for your arrest.
  \hfill \textit{The forger is a Breath-less agent. The seal was a test. You failed.}
\item \textbf{Covenant Audit (Hidden Cargo)} --- A judge demands to inspect your holds. The lower deck is full of escaped slaves.
  \hfill \textit{The slaves are from Ashaan. One of them is the judge's own niece. He does not know.}
\item \textbf{Lighthouse Out of Phase} --- The keeper signals a warning — but the warning is for the ship behind you.
  \hfill \textit{You are now the decoy. The reef is ahead. The slavers are behind. Choose.}
\item \textbf{Tariff Hike (Admiralty)} --- The port-master doubles duties on all Sidhi goods. The hike is a bribe.
  \hfill \textit{The bribe is to the Black Bishop. The Bishop has already accepted. Your protest is treason.}
\item \textbf{Shipbreaker's Curse} --- A vessel in the yard breaks its chains. It wants a new victim — a ship's captain.
  \hfill \textit{The captain is you. The ship was a slaver. It remembers your face from a raid you did not commit.}
\item \textbf{Breath-less Betrayal} --- An agent sells you to the Sisters. The agent is a slave in the Sisters' household.
  \hfill \textit{The betrayal was a test. The agent wants to see if you can escape. If you do, you are worthy of trust.}
\item \textbf{Green Lantern Lit (Drowned Priest)} --- Someone lights the forbidden lamp at the tavern. The dead come.
  \hfill \textit{They are looking for the person who killed them. That person is in your party.}
\item \textbf{Archive Fire (Galanina)} --- A lamp falls in the vault of sealed jars. Ledgers burn.
  \hfill \textit{The smoke forms names. Those names are the ones the Covenant wants erased. One is yours.}
\item[J] \textbf{The Drowned Captain's Ship Appears} --- A spectral dhow emerges from the fog. It offers passage.
  \hfill \textit{The price is a memory. The memory will be forgotten forever. The destination is the Hollow.}
\item[Q] \textbf{Admiralty Press Gang} --- Marines grab bystanders for the fleet. The marines are drunk and do not recognise rank.
  \hfill \textit{One of the bystanders is a Covenant judge. She is about to be impressed. She asks for your help.}
\item[K] \textbf{The Ninth Well Runs Dry} --- The sacred spring stops flowing. Galanina's courts close for lack of witness water.
  \hfill \textit{The well was poisoned by an Ashaani agent. The poison is a memory. Drinking it makes you forget your name.}
\item[A] \textbf{The Sealed Coast Opens} --- The mythical shore appears. Ancestors walk the tide line, asking for their names.
  \hfill \textit{If you give them a name, they will give you a secret. The secret is always a death.}
\end{enumerate}

\paragraph*{(Charter/Seal/Key)} Lighthouse charter (right to anchor in any Sidhi port); freedman's seal (proof of emancipation — glows when a slave ship is near); archive key (opens one door in the Sealed Archive — the key is a finger bone).

\section*{Diamonds --- Rewards/Leverage}
\label{sec:sidhi-rewards}
\begin{enumerate}
\setcounter{enumi}{1}
\item \textbf{Lighthouse Charter} --- A brass key. The right to anchor in any Sidhi port without inspection.
  \hfill \textit{The key fits a lock on a sea-chest, not a harbour gate. The chest contains a map.}
\item \textbf{Freedman's Seal} --- A brand on a leather thong. Proof of emancipation.
  \hfill \textit{The seal glows when a slaver ship is within a day's sail. It also glows when a Black Hand agent is near.}
\item \textbf{Archive Key (Finger Bone)} --- A bone key. Opens one door in the Sealed Archive of Galanina.
  \hfill \textit{The Archive keeps the corresponding finger. Touch the key to the door, and the finger twitches.}
\item \textbf{Shipwright's Boat} --- A small, fast vessel built for escape. No name, no papers.
  \hfill \textit{The boat cannot be tracked by magic. The wood is from a freed slave's house.}
\item \textbf{Covenant Seal Ring (Brass)} --- Press it into wax to witness any maritime contract.
  \hfill \textit{The ring is warm. It cools if the contract is false. The cooling is immediate.}
\item \textbf{Lighthouse Keeper's Lens} --- A brass tube. Look through it to see the nearest slaver ship.
  \hfill \textit{The lens also shows the faces of the slaves on board. They look back.}
\item \textbf{Salt Farmer's Tithe (Bag)} --- A bag of bitter salt. Legal tender in any Covenant court.
  \hfill \textit{The salt tastes of tears. It also tastes of freedom — if you know what freedom tastes like.}
\item \textbf{Breath-less Cipher Ring (Silver)} --- Twist it to reveal a hidden message. The message is a name.
  \hfill \textit{The name is a debt. The debt is transferable. Sign it over to avoid payment.}
\item \textbf{Smuggler's Brass Fingers} --- A set of prosthetic fingers that pick any lock (once per lock).
  \hfill \textit{The fingers were made from a customs trap. The trap remembers. It may trigger again.}
\item[J] \textbf{Black Bishop's Ring (Ship's Nail)} --- Wear it to be immune to Admiralty arrest in Sidhi ports.
  \hfill \textit{The nail is from a freed slaver. The slaves' ghosts protect you. They also watch.}
\item[Q] \textbf{The Drowned Priest's Coin} --- A copper coin from the ghost's collection. Give it to the ferryman.
  \hfill \textit{Once, you may cross from the living shore to the Hollow and back. The crossing takes one night.}
\item[K] \textbf{First Freedman's Story} --- A true tale told at the Ninth Well. It reveals a hidden lighthouse.
  \hfill \textit{The lighthouse marks a passage through the Brass Coast's tariff zone. No one will tax you there.}
\item[A] \textbf{The Sealed Coast's Water} --- A flask from the mythic shore. Drink it to remember a name you forgot.
  \hfill \textit{The name is always an ancestor's. The ancestor will then ask you to avenge them.}
\end{enumerate}

\section*{Quick use notes}
\label{sec:sidhi-quick-use}
\begin{itemize}
\item Draw until you have all four suits: Spade = place, Heart = actor, Club = pressure, Diamond = leverage. Highest rank sets the main Clock (2–5 → 4, 6–10 → 6, J/Q/K → 8, A → 10).
\item Diamonds are codified outcomes (charters, seals, keys) that change position rather than call for a roll.
\item If any A appears, echo \textbf{sea-omens}—green flames that burn blue, salt that tastes sweet, a drowned face in a still puddle.
\end{itemize}

\subsection*{Additional Features}
\begin{itemize}
\item \textbf{Lighthouse Courts:} Sidhi legal disputes are settled in the lamp room of a lighthouse, with the green flame as witness. The judge is the keeper. The verdict is written in salt on a slate, then washed away by the tide. Only the flame remembers.
\item \textbf{Freedman's Seal:} A brand or pendant that proves emancipation under the Covenant. A freedman cannot be re-enslaved in any Covenant signatory port. The seal is also a key — it opens the door to the Breath-less network.
\item \textbf{The Ninth Well:} A sacred spring in Galanina where the first Covenant judge drowned a liar. The water is said to remember every false oath. Sidhi who need to prove their innocence drink from the well; liars vomit brine.
\end{itemize}

\subsection*{Lore Echoes (Cross-Expansion Hooks)}
\begin{itemize}
  \item \textbf{Galanina and the Veiled Sisters (Ashaan):} The Sidhi diaspora exists because of Ashaani slavery. A Breath-less agent in the Brass Coast may be the same person who escaped Galanina as a child. Revenge is a slow debt.
  \item \textbf{Lighthouse Charters and the Fhara (Way of Silk):} Fhara water-clerks respect Sidhi lighthouse charters as proof of credit. A Sidhi merchant can borrow water rights against a charter; default means the lighthouse lamp goes out.
  \item \textbf{Freedman's Seals and the Covenant (Mykkiel):} The Covenant recognises only one seal higher than a freedman's — the Seal of Mykkiel himself. A party carrying a freedman's seal can demand a hearing before any Covenant judge, anywhere.
  \item \textbf{The Brass Coast and Dhaharan Independence:} The Sidhi of the Brass Coast are a major faction in Dhaharan politics. They pay tariffs to the Admiralty but run their own courts, their own ships, and their own spies. A Sidhi port-master may be the key to breaking the Admiralty's hold.
\end{itemize}

\subsection*{Sidhi Superstitions (burn for +1 die once)}
\begin{itemize}
  \item \textbf{Bread to the Tide:} Crumble a crust into the harbour at dusk. Ignore the first \emph{Ashaani Raid} penalty.
  \item \textbf{Knot the Anchor Chain:} Tie a single sailor's knot in a mooring rope. Cancel \emph{Forged Freedman's Seal} or take –1 die on your next legal roll (your choice).
  \item \textbf{Name to the Lighthouse:} Whisper a dead captain's name into the green flame. Gain +1 Effect when smuggling this scene, but mark \emph{Obligation} to the Drowned Priest.
\end{itemize}

\begin{tcolorbox}[colback=black!3,colframe=black!40!white,title={Patronage \& Power}]
Among the Sidhi, power flows through witness, salt, and the green flame of the lighthouse. The Covenant courts provide law, but the true authority is memory — of the raids, of the escape, of the ancestors who did not return. A Sidhi elder's word is law because she remembers when law was a whip.

\textbf{For the GM:}  
Power in Sidhi society is measured in charters, seals, and the trust of the lighthouse keeper. Patronage often takes the form of lighthouse rights, freedman's proofs, or archive keys — each tied to a specific port, a specific ancestor, a specific debt. To emphasize this:
\begin{itemize}
\item Tie rewards to physical tokens (keys, seals, lenses) that can be lost, stolen, or voided by the Admiralty.
\item Let rival Sidhi port-masters issue conflicting charters, forcing players to choose whose lamp to trust.
\item Use the lighthouse courts, hidden coves, and freedman's squares as arenas for social contests, where a green flame is a witness and a bag of salt is a verdict.
\end{itemize}
In Sidhi society, the sea forgets nothing — and the lighthouse remembers every face.
\end{tcolorbox}

% Index entries
\index{Sidhi|see{Diaspora of the Sealed Coast}}
\index{Theme generator}
\index{Diaspora generator}
\index{Places!Sidhi}
\index{People!Sidhi}
\index{Complications!Sidhi}
\index{Rewards!Sidhi}
\index{Lighthouse courts}
\index{Freedman's seal}
\index{Galanina}
\index{Brass Coast}
\index{Covenant of Mykkiel}
\index{Breath-less}
\index{Ashaani slavery}
\index{Ninth Well}
\index{Sealed Coast}
\index{Saint Adar!Sidhi}
\index{Saint Hervor!Sidhi}
\index{Tide-Washed Covenant!Sidhi}

\section*{Thematic SB Spend Table}
\label{sec:sidhi-sb}

\subsection*{Minor Complications (1 SB)}
\begin{itemize}
\item \textbf{Exposure:} Your actions draw unwanted attention from \textbf{Lighthouse judges or Admiralty customs}.
\item \textbf{Noise:} Sounds of your actions alert nearby \textbf{port-masters or freedmen's advocates}.
\item \textbf{Trace:} Evidence of your passage marks your route for \textbf{Ashaani slavers or Covenant auditors}.
\item \textbf{Delay:} A brief but meaningful setback costs you \textbf{time or favourable tide}.
\item \textbf{Supply Strain:} Mark +1 segment on a relevant \textbf{resource clock}.
\end{itemize}

\subsection*{Moderate Setbacks (2 SB)}
\begin{itemize}
\item \textbf{Alarm Raised:} \textbf{Local lighthouse keeper or Covenant judge} becomes aware and begins responding.
\item \textbf{Position Lost:} You lose advantageous ground/cover/stealth due to \textbf{green flash or rising tide}.
\item \textbf{Foe Appears:} An \textbf{Ashaani slaver or Admiralty press gang} arrives on scene.
\item \textbf{Gear Trouble:} A piece of equipment becomes \textbf{Compromised/Neglected}.
\item \textbf{Lock/Barrier:} A simple obstacle now requires a test to overcome.
\end{itemize}

\subsection*{Serious Trouble (3 SB)}
\begin{itemize}
\item \textbf{Reinforcements:} Additional \textbf{Admiralty marines, Covenant bailiffs, or slaver ships} arrive.
\item \textbf{Key Gear Breaks:} A crucial tool/weapon becomes temporarily unusable.
\item \textbf{Major Twist:} The situation fundamentally changes—\textbf{lighthouse goes dark/charter voided/seal breaks}.
\item \textbf{Rail Tick:} Advance a relevant campaign/front clock by 1 segment.
\item \textbf{Condition Applied:} Mark \textbf{Fatigue 1/Harm 1/Condition} appropriate to fiction.
\end{itemize}

\subsection*{Major Turns (4+ SB)}
\begin{itemize}
\item \textbf{Trap Springs:} A prepared danger activates with full effect.
\item \textbf{Authority Arrival:} \textbf{High Lighthouse Judge, Black Bishop, or Drowned Priest} intervenes.
\item \textbf{Scene Shift:} The environment changes dramatically—\textbf{lighthouse collapses/tide swallows the court/slaver ship rams the pier}.
\item \textbf{Patron Omen:} Divine/arcane forces take notice—\textbf{omen appears/blessing lost/curse manifests}.
\item \textbf{Narrative Pivot:} The story takes an unexpected turn that reframes objectives.
\end{itemize}

\subsection*{Region-Specific SB Options}
\begin{itemize}
\item \textbf{Sidhi (Lighthouse Law):} Green flame flickers, salt crystals form in the air, the judge's water-clock runs backwards.
\item \textbf{Sidhi (Diaspora):} A freedman's seal grows warm, a distant lighthouse answers a question you did not ask, a slave ship's bell rings without being struck.
\item \textbf{Sidhi (Brass Coast):} Ropeways hum with whispers, spice dust forms the shape of Galanina's cliffs, the Admiralty's flag flies at half-mast for no reason.
\end{itemize}

\subsection*{Pursuit Ladder (Near / Mid / Far)}
\begin{itemize}
\item \textbf{Advance/Withdraw (Wits + Sail):} On a hit, shift one band. On a strong hit, also impose \emph{Green Flash}.
\item \textbf{Cut the Charter (Presence + Sway):} On a hit, freeze range for one exchange; if a freedman's seal is in your possession, +1 die.
\item \textbf{Weather the Raid (Resolve + Spirit):} On a hit in \emph{Ashaani Raid}, you pick who holds the light; on a miss, the GM does.
\end{itemize}

\subsection*{Boss Seeds (Salt \& Flame)}
\begin{itemize}
\item \textbf{The Seal-Breaker (Nine Forgeries)} — \emph{Phases}: Freedman Seals Counterfeited $\rightarrow$ Orchestrated Re-enslavement $\rightarrow$ Lighthouse Coup. \emph{Cracks}: recover a true seal, expose a forger, survive a night in the Ninth Well.
\item \textbf{The Drowned Admiral (Ghost of the Brass Coast)} — \emph{Phases}: Lost Ships Called Home $\rightarrow$ Procession of Empty Lighthouses $\rightarrow$ Tide-Baptism of a Judge. \emph{Cracks}: relic lighthouse lens, survivor testimony, the green flame turning white.
\end{itemize}

\newpage

% -------------------- PERESHI --------------------
\chapter{Pereshi}
\emph{The Roof of the Way.}

\noindent
Sarnai says: \textit{“The Pereshi keep coals that never cool. They speak of a prince whose name they erased. I asked a fire‑keeper why they remember him if they erased him. She said, ‘We do not remember him. We remember the debt.’”}

\noindent
The Pereshi highlands are a vast, windswept plateau, cut by deep river valleys and crowned with snow‑capped peaks. Ancient fire‑shrines stand beside caravanserais; terraced vineyards climb the hillsides. The people are scholars, shepherds, and traders – proud, suspicious, and bound by a code of hospitality that demands three days of feasting for any guest, and three generations of vigilance for any insult. The fire in the temple is not just a flame. It is a witness. It remembers every oath, every treaty, every lie told in its light.

\vspace{0.5em}
\noindent\textit{“Never step on a fallen book. Written words carry the weight of law and the weight of witness. To step on a book is to trample a contract, a genealogy, a verdict. The spirits of the words will haunt you.”} \\
— Flame‑Keeper Zarifa

\vspace{1em}
\begin{almanac}
\textbf{What do people eat for breakfast?} Flatbread with honey and clotted cream – the honey from highland bees, the cream from goats that graze the temple steps. The goats are sacred. The goats are also delicious.

\textbf{What does a child learn before they can walk?} To read a contract. The eight clauses of the Covenant. The ninth clause is never written – it is the silence where the unspoken terms live. The ninth clause is the one that kills you.

\textbf{What is the first thing you notice about a stranger?} Their book. Every Pereshi carries a personal ledger of debts owed and collected. A thick book means many debts. A thin book means many creditors. No book means a liar. Liars are not fed.

\textbf{What is the most common crime?} Contract forgery – but forgery is an art. The judges are forgers themselves. The law is a negotiation that never ends.

\textbf{What does no one talk about but everyone knows?} The Forgotten Prince did not vanish. He is buried beneath the largest fire temple, in a lead‑lined coffin sealed with nine locks. The ninth lock has no key – the key is a memory no one wants to have.

\textbf{Where do people go when they die?} To the fire. Cremated on the temple pyres, ashes mixed with clay and pressed into bricks. The bricks build new temples. The dead become walls. The walls witness.

\textbf{How do you apologize for a serious wrong?} You offer a coal – a burning ember from the temple fire. Placed in the wronged party’s hearth. If it stays lit, you are forgiven. If it goes out, you are cursed. The curse is the fire’s silence.

\textbf{What is the most important festival?} The Reading of the Ledgers (autumn). Every household brings its ledger to the temple. The debts are tallied. The dead are named. The fire listens. This year, the ninth coal has not stopped flickering.

\textbf{What is the secret that could destroy a powerful person?} The ninth coal is not coal. It is a piece of the Forgotten Prince’s heart, petrified. It is still warm. It is still beating.
\end{almanac}

\subsection*{Etiquette Hooks}
\begin{itemize}
  \item \textbf{Fruit \& Tea}: Offer fruit and a warm cup on arrival → +1 Position in household negotiations.
  \item \textbf{Shoes at Threshold}: Remove shoes indoors → negate a Position –1 penalty from prior brusqueness.
  \item \textbf{Couplet of Courtesy}: A brief couplet acknowledging hosts’ craft → Audience: Cool → Warm with artisans/scribes.
  \item \textbf{Faux Pas (SB)}: Dust on the Mat – barging in tracks Exposure +1 among neighbours.
\end{itemize}

\subsection*{Strings (sample)}
\begin{itemize}
  \item \textbf{Fire‑Shrine Token} – +1 Position in any Pereshi court or negotiation.
  \item \textbf{Rose from the Ruined Garden} – burn to ask a yes/no question of the fire.
  \item \textbf{Tablet of a Severed Name} – contains an erased name; speak it to restore a person’s memory.
\end{itemize}

\subsection*{Venue Tags}
\begin{itemize}
  \item \textbf{Roof‑Garden Pact}: settlements on rooftops; a hosted pact counts as a String across partnered houses.
  \item \textbf{Archive Court}: citations from the Fire‑Temple archive lower DV by 1 once/scene.
\end{itemize}

\newpage
\section{Pereshi — ``The Roof of the Way''}
\label{chap:pereshi}

\subsection*{Quick Reference}
\textbf{Tagline:} \textit{The fire remembers. The archive endures. The Fhara fear us; the Tulkani remember us; the lowlands do not see us. That is how we like it.}

\begin{quote}
  \textbf{Flame-Keeper (Elite):}  
  ``The fire does not forget. It has witnessed every oath sworn in this valley, every treaty sealed with a dropped coal, every lie told in its light. I do not tend the flame. I listen to what it remembers. And it remembers everything – including the Prince who thought he could burn us. We took his name. We cut his shadow. But shadows grow back, and the ninth coal is warm.''
\end{quote}

\begin{quote}
  \textbf{Caravan Scribe (Commoner):}  
  ``My grandfather copied a tax ledger that had not been opened in three centuries. He found a mistake – a single miscounted bale of silk. That mistake started a war. We do not write for the living. We write for the dead who cannot correct their errors. And we write to remind ourselves: the empire of ash will not rise again. But the Prince's shadow is long, and the ninth coal has started to hum. The Fhara send envoys who cannot meet our eyes. The Tulkani leave knots at our thresholds. We do not explain. We burn.''
\end{quote}

\textbf{Genre / Mood:} High alpine bureaucracy, firelit memory, postcolonial witness, cursed return.

\textbf{Starting Location:} The scriptorium of a mountain caravanserai, where lamps burn at noon and scribes copy tax ledgers by candlelight. A Flame-Keeper looks up from a brazier: ``You have come for a tariff receipt, or for a verse to cross the pass, or to consult a book that should have been burned. First, tell me a truth – something the Ashaani Prince tried to erase. Then I will decide which fire shows you the way. And be quick: the ninth coal is warmer than it should be. It has been humming in my sleep.''

\vspace{2mm}
\begin{tcolorbox}[colback=black!3,colframe=blue!60!black,title={Theme \& Atmosphere}]
The Pereshi highlands are a vast, windswept plateau, cut by deep river valleys and crowned with snow-capped peaks. Ancient fire temples stand beside caravanserais; terraced vineyards climb the hillsides. The people are scholars, shepherds, and traders – proud, suspicious, and bound by a code of hospitality that demands three days of feasting for any guest, and three generations of vigilance for any insult.

But the Pereshi are not a nation. They are a memory of one. Long ago, a highland empire stretched from the western passes to the eastern salt flats – a confederation of fire temples, fortress monasteries, and clan valleys that answered to a council of Flame‑Keepers. That empire fell when the Ashaani Prince marched east, burning libraries and breaking the hearths. What remains is a scatter of independent holdings – cities like Mount Khvarena, temple towns like Adur‑Gah, and countless villages that recognise no king, only the fire and the ninth coal in every hearth.

The Pereshi share a culture: the same fire rituals, the same oral epic (the \textit{Book of Ash}), the same fear of naming the dead. But they do not unite. They have no crown, no single council. The Flame‑Keepers meet once every nine years to renew the seal on the Prince’s heart, but they do not govern. They witness. The rest of the time, each valley minds its own thresholds. This fragmentation is a shield – it has kept conquerors from finding a single throat to cut. It has also kept the Pereshi from reclaiming what they lost.

\vspace{2mm}
\textbf{What you notice first:} The smell of woodsmoke and old parchment. The way every valley has a different accent. The constant call to prayer from fire‑temple towers – but the prayers are in a language no one speaks outside the highlands, a language that predates the Prince's rule. And the way every Flame‑Keeper touches the ninth coal before sleep, checking if it has grown hotter. Lately, they check twice.

\textbf{The one rule that kills newcomers:} Never step on a fallen book. Written words carry the weight of law – and the weight of witness. To step on a book is to trample a contract, a genealogy, a verdict, a memory of the Prince's cruelty. The spirits of the words will haunt you. Also, never speak the Prince's true name aloud. The scroll remembers. The shadow will hear. And the ninth coal will pulse.
\end{tcolorbox}

\subsection*{Regional Mood: The Weight of the Archive and the Warmth of the Coal}

Power here is measured in fire‑temple charters, caravan tariffs, and the length of a family's library. The petty lords of the highlands compete not through armies but through scribes who can produce the oldest charter, the most complete genealogy, the most beautifully illuminated tax receipt. The Covenant has a strong presence in the sacred cities, but the Pereshi follow their own law: the word of a Flame‑Keeper outweighs a judge's gavel. Travelers are expected to recite a verse when crossing a bridge – the spirits of the water demand it – and to present a red rose before petitioning a lord. The high passes are measured in breath‑counts, not miles, and the air is thin enough to make a foreigner's head spin. But the thinnest air of all is in the sealed archive where the Prince's name is buried. No one has opened it in nine generations. The lock has no key. The key is a memory no one wants to have – and lately, the lock has been weeping rust. The rust tastes of iron and old smoke.

\subsection*{The Unfired Crown: How the Pereshi Rule Themselves}
\label{subsec:pereshi-unfired-crown}

The Pereshi do not kneel to a king. They kneel to fire, to ash, and to a name they erased. Each valley or fire‑temple city governs itself through a council of elders, a Flame‑Keeper, and a war captain (if needed). There is no central treasury, no standing army, no code of laws beyond custom and the witness of the ninth coal. What unites them is not a crown but a memory of conquest: every Pereshi child learns the story of the Ashaani Prince, the burning of the great library at Khvarena, and the year the fires went out. They also learn the year the fires were relit – and the ritual that severed the Prince's name from the world.

The Flame‑Keepers of the nine largest temples meet once every nine years at Mount Khvarena, the holiest peak. They do not elect a leader. They sit in a circle, each holding a coal from their own hearth. They speak of weather, of harvests, of bandit raids. They do not speak of the Prince – but each of them, before sleep, places a hand on the ninth coal that sits in a brass brazier at the centre of the circle. If it is warm, they know the Shadow is stretching. If it is hot, they light a signal fire and send riders to every valley: \textit{The name stirs. Renew the seals.}

No outsider has ever attended a meeting of the nine. The Fhara have tried to send spies; the spies never returned, and their names were erased from every ledger. The Tulkani do not try; they leave a knotted cord at the edge of the temple and wait for a reply. The reply is always a single burnt knot.

\subsection*{The Fear of the Fhara and the Tulkani Unease}
\label{subsec:pereshi-fear-and-unease}

\textbf{The Fhara fear the Pereshi.} In the desert oases, water‑clerks whisper that every Pereshi carries a ``name‑knife'' – a coal that can erase a debtor from memory. A merchant whose name is struck from the ledgers is not bankrupt; he is \textit{forgotten}. His wells dry up. His caravans lose their way. His children cannot remember his face. The Fhara pay their debts to the Pereshi with unusual haste, and they never, ever allow a Pereshi fire‑keeper to light a lamp in their counting houses. Some say that the ninth well of Qalat al‑Madina was cursed not by the Kuvani but by a Pereshi exile who wanted to make an oasis vanish.

\textbf{The Tulkani have an uneasy relationship with the Pereshi.} They share a border of silence. Both respect the ninth night, but for different reasons. The Tulkani say the Pereshi stole fire from the Dreamer (Ikasha). The Pereshi say the Tulkani forgot how to build a hearth. They trade stories at the edge of the high passes, but never stay the night. A Tulkani who enters a Pereshi fire temple must leave a knot behind – and the Pereshi burn it. The smoke carries the dream to the sky, but the knot is gone, and the Tulkani cannot retie it. Some Tulkani elders whisper that the Pereshi know how to sever a name because they once tried to sever the Dreamer's name – and failed. They do not say this aloud. The Hollow might be listening.

\textbf{Everyone else underestimates the Pereshi.} Ecktorian clerks call them ``backward highlanders with no coin.'' Vhasian knights sneer at their unpaved roads and unarmoured sentinels. Oshiiran grain factors pity them. The Pereshi prefer it this way. Let the lowlands send their tax collectors; the passes are high, the cold is deep, and every tariff receipt requires a verse that no foreigner can recite correctly. Those who try are turned back at the bridge – or find that their names have been quietly written in ash and left to the wind.

\subsection*{The Forgotten Prince – Heart as Coal, Name as Ash}
\label{subsec:pereshi-forgotten-prince}

Once, the Ashaani Prince – a conqueror from the Veiled Throne, a man whose armies swept from the Red Desert to the eastern hills – turned his gaze to the Pereshi highlands. He burned the great library at Khvarena. He melted the temple bells into arrowheads. He tried to erase the fire‑faith, to replace it with the silence of the Veiled Sisters. For nine years he ruled from a garrison in the lowlands, sending patrols to cut down any who still lit a lamp before dawn.

Then his empire cracked. The Kalrantha fragment that had given him power over fate began to devour him from within. His generals turned on each other. His shadow began to walk independently. He retreated to a fortress in the eastern passes, and the Pereshi followed – not with armies, but with priests.

The ritual took nine nights. They did not kill him. They \textit{unmade} him. They spoke his true name into a lead scroll, sealed it with nine coals, and buried it beneath the largest fire temple. They cut his shadow from his body and pinned it to the Sun King's Throne (now lost in the Red Desert). And they took his heart – still beating, still warm – and placed it in a brass brazier. The ninth coal is not a relic. It is his petrified heart. Once a year, on the anniversary of his fall, it pulses. The Flame‑Keepers feel it. They do not speak of it.

The Prince's shadow still walks because the heart still dreams of conquest. The Pereshi do not fear him – they \textit{remember} him. He is the Debt That Was Paid. But the debt compounds. Every broken name, every forgotten oath, every coal that cracks adds interest. If the ninth coal ever goes cold, the Prince will not return – he will never have been gone.

\subsection*{The Prince's Shadow Clock [8]}
Tracks the Prince's gradual return. Advance the clock when:
\begin{itemize}
  \item The party breaks a Pereshi oath or defiles a fire temple (+1)
  \item A Flame‑Keeper dies without passing on the secret of the ninth coal (+1)
  * The party speaks the Prince's true name aloud (from the First Emperor's scroll) (+2)
  \item The GM spends 2 SB to represent a crack in the archive's seal (+1)
\end{itemize}
When the clock reaches 4: The Prince's shadow begins to appear in mirrors and still water. It asks questions. Answering gives it a thread to pull. When the clock reaches 8: The Prince manifests as a ghostly figure – not a combat boss, but a presence that can possess a willing (or desperate) host. The party must either reseal him (by finding a new name, a new shadow, or a new coal) or bargain with him (he offers great power at the cost of a Pereshi valley). The valleys are not his to give. He offers them anyway. The ninth coal will be warm that night.

\subsection*{Prominent NPCs of the Pereshi Highlands}
\label{sec:pereshi-npcs}

\begin{longtable}{|p{0.2\textwidth}|p{0.25\textwidth}|p{0.3\textwidth}|p{0.15\textwidth}|}
\hline
\textbf{Name} & \textbf{Role/Title} & \textbf{Motivation} & \textbf{Rumoured Quirk} \\
\hline
Flame‑Keeper Zarifa & Priestess of Mount Khvarena & Preserve the old rites and keep the ninth coal from lighting & She has never been seen without a rose – fresh or dried. The rose is from the Prince's own garden, taken the day he fell. She sleeps with one hand on the coal. The coal has begun to pulse in time with her heart. \\
\hline
Lord Rostam & Highland petty king (Adur‑Gah valley) & Unite the valleys against lowland tax collectors – and against the Prince's return & He keeps a cage of doves; each dove carries a different treaty. The ninth dove carries a blank page for the Prince's name. The dove has started to speak. It speaks in the Prince's voice. \\
\hline
Archivist Kaelen & Keeper of the Unfinished Books (Khvarena) & Protect the scrolls that were never completed – especially the Prince's name scroll & He has not left the archive in thirty years. His shadow is missing. The Prince's shadow took it. He does not sleep. He no longer needs to. \\
\hline
Tarkan & Caravan master of the high passes & Secure the passes before the snows close them – and before the Prince's ghost walks & He has three camels, each named after a dead wife. The camels were bred from stock the Prince's army abandoned. They refuse to cross the Pass of the Unforgiven. They kneel when the ninth coal pulses. \\
\hline
The Poet of the Ninth Verse & An old man who stole a missing line of an ancient epic & Recite the verse for a cup of tea made with snowmelt – and keep the Prince from hearing it & His left eye is a piece of amber with a fly trapped inside. The fly is from the Prince's own stable. When the Prince's shadow is near, the fly buzzes. When the coal pulses, the fly is silent. \\
\hline
Fire‑Keeper Zahra (new) & Young priestess of the eastern temple (Adur‑Gah) & Secretly correspond with a Fhara water‑clerk – she is torn between duty and mercy & She writes letters on bark. The bark is from a grove that the Prince burned. She signs with a drop of her own blood. The Fhara clerk has never seen her face. \\
\hline
The Knot‑less (Tulkani exile, new) & A Tulkani who cut her own cord, now living among the Pereshi & Learn the name‑severing ritual – to hide from the Hollow & She carries a single knot in her mouth. She has never swallowed it. She speaks by moving her lips around the knot. The Pereshi have not yet burned her offerings. \\
\hline
The Nameless (Fhara merchant, new) & A spice trader whose name was erased by a Pereshi fire‑keeper & Reclaim his lineage – even if it costs a shard of the Kalrantha amulet & He wears a mask of woven silver. Underneath, his face is smooth – no mouth, no nose, only two eyeholes. He drinks through a straw. He whispers his true name to the ninth coal every night, hoping it will remember. \\
\hline
\end{longtable}

\subsection*{Names of the Pereshi Highlands}
\label{sec:pereshi-names}

Inspired by highland and steppe traditions: names with soft endings, honorific prefixes (Mir, Khan, Beg), and surnames denoting a clan or a valley. Many families carry a second, secret name – the name they used during the Prince's occupation, preserved as a mark of survival. To speak a Pereshi’s secret name is to bind them to a promise; to learn it is to hold a thread of their soul.

\begin{longtable}{|p{0.25\textwidth}|p{0.25\textwidth}|p{0.2\textwidth}|p{0.2\textwidth}|}
\hline
\textbf{First Names (Male)} & \textbf{First Names (Female)} & \textbf{Surnames} & \textbf{Epithets} \\
\hline
Rostam & Zarifa & Kuhistani & \textit{the Fire‑Touched} \\
\hline
Kaelen & Laleh & Qasemi & \textit{the Unfinished} \\
\hline
Tarkhan & Shirin & Salari & \textit{Ninth Verse} \\
\hline
Dariush & Mahsa & Bahrami & \textit{the Bridge‑Mouthed} \\
\hline
Behrouz & Golnar & Khorasani & \textit{the Red Rose} \\
\hline
\end{longtable}

\paragraph*{(Temple/Library/Bridge)} Fire temple with a copper dome and a flame that never goes out; library of unfinished books where scribes copy texts that were burned by the Prince's order; stone bridge over a chasm, each arch carved with a line of poetry – except the ninth, which is deliberately blank. The blank arch hums at midnight.

\section*{Spades — Places}
\label{sec:pereshi-places}
\begin{enumerate}
\setcounter{enumi}{1}
\item \textbf{Fire Temple of the Unquenched} – A stone tower with a copper dome. The fire inside is said to be the first flame. It has never gone out – but it flickers when a king dies.
  \hfill \textit{The flame's colour changes with the season. The ninth coal sits in a brass brazier, untouched since the Prince's fall. Recently, the coal has started to glow at midnight. The glow casts no shadow.} `[RITUAL]` `[MEMORY]`
\item \textbf{Mount Khvarena (The Nine Fires)} – The holiest peak, where the Flame‑Keepers meet once every nine years.
  \hfill \textit{At the summit, a circle of nine braziers. The ninth brazier holds the Prince's heart. The coals never cool. Pilgrims climb the mountain barefoot, carrying a single rose. The rose is burned. The ash is the offering.} `[LEGEND]`
\item \textbf{The Archive of Severed Names} – A cave behind a waterfall, where clay tablets record every name ever erased.
  \hfill \textit{The tablets are stacked to the ceiling. Some are warm. Those names are still trying to return. The archivist has never opened the ninth row.} `[DEATH]` `[SECRET]`
\item \textbf{Adur‑Gah (Fire Temple City)} – A walled town built around a natural gas vent that burns day and night.
  \hfill \textit{The vent is called the Prince's Breath. The locals say it is his sigh, trapped underground. When the wind shifts, the flame roars – and the ninth coal pulses.} `[COMMUNITY]`
\item \textbf{The Pass of the Unforgiven} – A high defile where the Prince's army froze to the rock.
  \hfill \textit{On cold nights, you can still see their shapes in the ice. They are still standing. Their lips move. They are still counting. The Pereshi do not use this pass. They say the dead do not forgive.} `[DEATH]` `[FEAR]`
\item \textbf{The Library of Unfinished Books} – A scriptorium where scribes copy texts that the Prince tried to burn.
  \hfill \textit{The books have no endings. The endings were lost. The scribes write new ones, but the old endings whisper from the margins. The ninth shelf is empty. It is reserved for the book that will contain the Prince's true name.} `[KNOWLEDGE]`
\item \textbf{The Bridge of Nine Echoes} – A rope bridge over a frozen chasm. Each step echoes nine times.
  \hfill \textit{The ninth echo is always a whisper – the name of the next person to fall. Do not answer. The whisper is not for you.} `[DEATH]` `[OMEN]`
\item \textbf{The Rose Garden of Khvarena (Ruined)} – Where the Prince cut every rose bush to build a pyre.
  \hfill \textit{One bush survived. It grows from a crack in the stone. The roses are black. They smell of ash. A single petal can bind a name – or sever it.} `[MAGIC]`
\item \textbf{The Well of Unspoken Names} – A dry well where the Pereshi throw coals that have gone cold.
  \hfill \textit{The coals burn again at the bottom. No one knows how. Some say the Prince's heart keeps them warm.} `[MYSTERY]`
\item[J] \textbf{The Sunken Archive (Fhara border)} – A flooded chamber where water‑tablets and fire‑tablets meet.
  \hfill \textit{The water is warm. The tablets float. The Fhara water‑clerks who come here never leave. The Pereshi do not explain why.} `[THRESHOLD]`
\item[Q] \textbf{The Silent Pass} – A high mountain pass where no bird flies.
  \hfill \textit{The air is too thin for sound. Arguments must be written in the snow. The snow does not lie. The Pereshi leave offerings of silence here.} `[WEATHER]`
\item[K] \textbf{The Ninth Coal’s Vault} – A lead‑lined chamber beneath Mount Khvarena.
  \hfill \textit{The door has nine locks. The keys are held by nine Flame‑Keepers. The ninth key is a memory that no one can recall. The door sweats.} `[SEAL]`
\item[A] \textbf{The First Emperor’s Scroll (Myth)} – A scroll that contains the Prince's true name, written in ash on lead.
  \hfill \textit{It is kept in a box that has no key. The box is opened only when the ninth coal cracks. The last time it cracked, the Prince's shadow walked for nine days.} `[LEGEND]`
\end{enumerate}

\paragraph*{(Priest/Scholar/Envoy)} Flame‑Keeper with a rose and a coal; archivist with ink‑stained fingers; Fhara merchant with a woven silver mask.

\section*{Hearts — People \& Factions}
\label{sec:pereshi-people}
\begin{enumerate}
\setcounter{enumi}{1}
\item \textbf{Flame‑Keeper Zarifa} – Priestess of Mount Khvarena, keeper of the ninth coal.
  \hfill \textit{She has never been seen without a rose – fresh or dried. The rose is from the Prince's own garden. She sleeps with one hand on the coal. The coal has begun to pulse in time with her heart.} `[RITUAL]`
\item \textbf{Archivist Kaelen} – Keeper of the Unfinished Books, guardian of the Prince's name scroll.
  \hfill \textit{He has not left the archive in thirty years. His shadow is missing. The Prince's shadow took it. He does not sleep. He no longer needs to.} `[SECRET]`
\item \textbf{Caravan Master Tarkan} – A highland trader who knows every pass and every ghost.
  \hfill \textit{He has three camels, each named after a dead wife. The camels were bred from stock the Prince's army abandoned. They refuse to cross the Pass of the Unforgiven. They kneel when the ninth coal pulses.} `[TRAVEL]`
\item \textbf{Lord Rostam} – A petty king of Adur‑Gah, trying to unite the valleys.
  \hfill \textit{He keeps a cage of doves; each dove carries a different treaty. The ninth dove carries a blank page for the Prince's name. The dove has started to speak. It speaks in the Prince's voice.} `[POLITICS]`
\item \textbf{The Poet of the Ninth Verse} – An old man who stole a missing line of the epic.
  \hfill \textit{His left eye is a piece of amber with a fly trapped inside. The fly is from the Prince's own stable. When the Prince's shadow is near, the fly buzzes. When the coal pulses, the fly is silent.} `[KNOWLEDGE]`
\item \textbf{Fire‑Keeper Zahra (new)} – A young priestess of the eastern temple, secretly writing to a Fhara water‑clerk.
  \hfill \textit{She writes letters on bark. The bark is from a grove that the Prince burned. She signs with a drop of her own blood. The Fhara clerk has never seen her face.} `[CONFLICT]`
\item \textbf{The Knot‑less (Tulkani exile, new)} – A Tulkani who cut her own cord, now living among the Pereshi.
  \hfill \textit{She carries a single knot in her mouth. She has never swallowed it. She speaks by moving her lips around the knot. The Pereshi have not yet burned her offerings.} `[MYSTERY]`
\item \textbf{The Nameless (Fhara merchant, new)} – A spice trader whose name was erased by a Pereshi fire‑keeper.
  \hfill \textit{He wears a mask of woven silver. Underneath, his face is smooth – no mouth, no nose, only two eyeholes. He drinks through a straw. He whispers his true name to the ninth coal every night, hoping it will remember.} `[TRAGEDY]`
\item \textbf{Covenant Judge} – A grey‑robed notary stationed at the bridge between Pereshi and Vhasian lands.
  \hfill \textit{She carries a water‑clock that runs on melted snow. It never freezes. She has never ruled against a Pereshi – not because she fears them, but because their contracts are written in ash, and ash has no appeal.} `[JUSTICE]`
\item[J] \textbf{The First Censor (Ghost of Khvarena)} – A librarian who burned herself with the scrolls to keep the Prince from reading them.
  \hfill \textit{She appears as a pillar of smoke. She asks: ``Is the name still sealed?'' If you say yes, she weeps ash. If you say no, she screams and the ninth coal cracks.} `[DEATH]`
\item[Q] \textbf{The Prince's Shadow (Manifestation)} – A silhouette without a face, seen in mirrors and still water.
  \hfill \textit{It asks for a name. Any name. Give it a living one, and that person will dream of the Prince every night. Give it a dead one, and the dead will walk for one night. Give it your own, and you lose your shadow – and your name becomes ash.} `[CURSE]`
\item[K] \textbf{The Ninth Flame‑Keeper (Legend)} – The one who sealed the Prince's heart. She never died.
  \hfill \textit{She sits at the bottom of the Well of Unspoken Names, warming the coals with her breath. She has not breathed in nine centuries. The coals breathe for her.} `[MYTH]`
\item[A] \textbf{The Prince's Echo (Forgotten Conqueror)} – A voice from the ninth coal, heard when it pulses.
  \hfill \textit{It speaks in a language no one remembers. It says only one word. The word is a name – yours, or the Prince's. It changes each time.} `[FEAR]`
\end{enumerate}

\paragraph*{(Ash/Name/Coal)} A tablet from the Archive of Severed Names (contains an erased name); a rose from the Ruined Garden (can be burned to ask a yes/no question of the fire); a fragment of the ninth coal (once, reroll a failed negotiation with a Pereshi).

\section*{Clubs — Complications/Threats}
\label{sec:pereshi-complications}
\begin{enumerate}
\setcounter{enumi}{1}
\item \textbf{Name‑Debt Called} – A Pereshi fire‑keeper demands a name you do not wish to give.
  \hfill \textit{The name must be true. It must be someone you know. The fire‑keeper will write it on a clay tablet and place it in the Archive of Severed Names. That person will begin to be forgotten – first by strangers, then by family, then by you.} `[DEBT]` `[SOCIAL]`
\item \textbf{The Ninth Coal Cracks} – The Prince's heart splits. The Shadow stretches.
  \hfill \textit{Advance the Prince's Shadow clock by 2. All mirrors in the scene show a silhouette standing behind each character. The silhouette asks: ``Why do you keep my name?'' Refuse to answer, and it will take a memory.} `[DEATH]` `[FEAR]`
\item \textbf{Fhara Envoy Refused} – A water‑clerk from the oases arrives with gifts. The Pereshi burn them.
  \hfill \textit{The envoy is the Nameless – a merchant whose name was erased. He offers the party a shard of the Kalrantha amulet in exchange for a promise to intercede. The Pereshi will not negotiate. The party must choose between the merchant's gratitude and the fire‑keeper's ire.} `[POLITICS]`
\item \textbf{Tulkani Knot Left at the Threshold} – A Tulkani Caravan Judge has left a knotted cord outside a fire temple.
  \hfill \textit{The knot is a request for parley. The Pereshi will not touch it. They will burn it at dusk. If the Tulkani returns before the fire goes out, they may speak – but only in riddles. The party is asked to act as interpreters. A wrong answer may start a feud.} `[SOCIAL]`
\item \textbf{Ash from the Prince's Pyre} – A storm carries ash from the Red Desert. The ash settles on Pereshi fields.
  \hfill \textit{The ash is warm. It smells of old smoke. The Pereshi say it is the Prince's breath, scattered. Any promise written in it will be forgotten by dawn. The party finds a message written in ash on their tent. It says: ``You carry my name.''} `[OMEN]`
\item \textbf{The First Emperor's Scroll Found} – A copy of the scroll containing the Prince's true name surfaces in a lowland market.
  \hfill \textit{The Pereshi will pay anything – a name, a coal, a memory – to recover it. The party is hired to retrieve it. But other factions want it too: the Veiled Sisters (to free the Prince), the Breath‑less (to destroy it), and a Fhara merchant (to sell it to the highest bidder).} `[QUEST]`
\item \textbf{The Ninth Verse Spoken Aloud} – Someone recites the missing line of the epic.
  \hfill \textit{The line is: ``And the heart that burned, the name that ash – the ninth is warm, the shadow waits.'' The moment it is spoken, the ninth coal pulses. A Oath‑Keeper appears. It asks: ``Who taught you that verse?'' The speaker must answer truthfully or lose their voice.} `[CURSE]`
\item \textbf{Fire Temple Silence} – All the fires in a temple go out at once. The ninth coal is cold.
  \hfill \textit{The Flame‑Keeper is found dead, a rose in her mouth. The ninth coal is not cold – it is gone. Stolen. The party must track the thief before the Prince's shadow claims the coal. The thief is a Tulkani exile who thought the coal would let her cut her last knot.} `[CRIME]`
\item \textbf{The Archive Weeps} – Water seeps from the walls of the Archive of Severed Names.
  \hfill \textit{The water tastes of salt and old ink. It is the tears of the erased names. They are trying to return. A character who touches the water hears a single name whispered – a name they have never heard. That name now belongs to them. The Hollow will expect it.} `[DEBT]`
\item[J] \textbf{The Prince's Shadow Points} – The shadow appears in a mirror and points at a party member.
  \hfill \textit{That character now owes a name to the Prince. They have nine days to find someone else's name to offer, or the Prince will take theirs. The shadow will appear in every reflection until the debt is paid.} `[CURSE]` `[DEBT]`
\item[Q] \textbf{The Rose Garden Burns (Again)} – Someone sets fire to the ruined garden. The black roses scream.
  \hfill \textit{The smoke forms the shape of a crown – the Prince's crown. The fire was lit by a Fhara merchant who thought he could destroy the Prince's memory. Instead, he woke it. The Prince's shadow now walks toward the lowlands. The party must warn the Pereshi – or bargain with the Prince.} `[DISASTER]`
\item[K] \textbf{The Ninth Flame‑Keeper Rises} – The legendary priestess climbs out of the Well of Unspoken Names.
  \hfill \textit{She has not aged. Her eyes are coals. She says: ``The seal is failing. I need a new name to bind.'' She looks at the party. She asks: ``Who volunteers?'' Refusal will not be accepted. The party must offer a false name, a dead name, or a memory of a name. The flame‑keeper will know if it is false.} `[LEGEND]` `[FEAR]`
\item[A] \textbf{The Prince Walks} – The ninth coal cracks completely. The Prince's shadow becomes solid.
  \hfill \textit{He does not fight. He negotiates. He offers: ``Give me a kingdom, and I will give you the world. Give me a name, and I will forget yours. Give me a coal, and I will warm your hearth forever.'' The Pereshi will not help. They are hiding. The party must decide alone.} `[APOCALYPSE]`
\end{enumerate}

\paragraph*{(Ash/Name/Coal)} Fragment of the ninth coal (once, reroll a failed negotiation with a Pereshi); a rose from the Ruined Garden (burn it to ask a yes/no question of the fire); a tablet from the Archive of Severed Names (contains an erased name – can be used to blackmail someone who has forgotten that name).

\section*{Diamonds — Rewards/Leverage}
\label{sec:pereshi-rewards}
\begin{enumerate}
\setcounter{enumi}{1}
\item \textbf{Fire‑Temple Token (Copper)} – A small disc stamped with a flame. Gain +1 Position in any Pereshi court or negotiation.
  \hfill \textit{The token is warm. It cools when you lie. The Pereshi will notice.} `[SOCIAL]`
\item \textbf{Rose from the Ruined Garden} – A black rose, dried, from the bush that survived the Prince's pyre.
  \hfill \textit{Burn it to ask a yes/no question of the fire. The fire will answer with a colour: red for yes, blue for no, white for the question is forbidden. The rose crumbles after use.} `[MAGIC]`
\item \textbf{Tablet of a Severed Name} – A clay tablet from the Archive, inscribed with a name that was erased.
  \hfill \textit{The name belongs to a person who has been forgotten – even by themselves. Speak the name aloud, and that person will remember who they are. Use it to blackmail someone who once owed that person a debt. The tablet crumbles after one use.} `[SECRET]` `[DEBT]`
\item \textbf{Fragment of the Ninth Coal} – A piece of the Prince's heart, still warm.
  \hfill \textit{Once, reroll a failed negotiation with a Pereshi. The fragment then turns to ash. But the Prince will remember your face. The Hollow will too.} `[LEGEND]`
\item \textbf{The Ninth Verse (Scroll)} – The missing line of the epic, written on a scrap of leather.
  \hfill \textit{Recite it aloud to summon a Oath‑Keeper. The Oath‑Keeper will answer one question about the Prince's shadow – but it will then ask you a question in return. You must answer truthfully or lose your voice.} `[KNOWLEDGE]` `[DEBT]`
\item \textbf{A Coal from Mount Khvarena} – A regular coal, blessed by the ninth flame.
  \hfill \textit{Carry it to be immune to cold weather for one leg of a journey. The coal glows faintly when the Prince's shadow is near.} `[TRAVEL]`
\item \textbf{Archivist's Seal} – A clay seal that opens any archive door in the highlands.
  \hfill \textit{The seal is shaped like a rose. It must be pressed to the lock. The lock will open – but the archivist will know who opened it. The seal works once.} `[ACCESS]`
\item \textbf{The Prince's True Name (in a lead scroll)} – A copy of the name, written in ash.
  \hfill \textit{Speak it aloud, and the Prince's shadow will manifest and obey one command. Then the scroll burns, and the Prince will hunt you. Use once. The GM gains 5 SB.} `[LEGEND]` `[CURSE]`
\item \textbf{Fhara Water‑Clerk's Ledger (Stolen)} – A ledger that records every name the Pereshi have erased for Fhara merchants.
  \hfill \textit{The ledger is a weapon. Present it to a Fhara well‑judge, and they will cancel a debt. Present it to a Pereshi fire‑keeper, and they will grant you a favour – but they will also ask whose name you want to erase in return.} `[TRADE]` `[DEBT]`
\item[J] \textbf{The Knot‑less's Last Knot} – The knot that the Tulkani exile carries in her mouth.
  \hfill \textit{If you untie it, the exile will remember her name – but she will also owe you a story. If you cut it, she will die, and you will gain the ability to hide from the Hollow for one night (the Hollow cannot see you). The knot is warm. It smells of sleep.} `[MYSTERY]`
\item[Q] \textbf{The Nameless's True Name (Whispered)} – The name that the Fhara merchant lost.
  \hfill \textit{Whisper it to him, and he will grant you a boon: a shard of the Kalrantha amulet, or a map to the archive of severed names. But the Pereshi will know you spoke his name. They will mark you.} `[TRADE]` `[DEBT]`
\item[K] \textbf{The Ninth Flame‑Keeper's Rose} – A rose that never wilts, from the legendary priestess.
  \hfill \textit{Once, present it to any Pereshi to demand a single act of mercy – the release of a prisoner, the return of a name, a night's shelter during a storm. The rose will then turn to ash. The priestess will ask why you used it. Answer poorly, and she will take a memory.} `[LEGEND]`
\item[A] \textbf{The Prince's Shadow (Bound)} – A piece of the Prince's shadow, captured in a lead‑lined mirror.
  \hfill \textit{Break the mirror, and the shadow will fight for you for one scene (Cap 3, Harm 3). Then it will turn on you. The mirror is heavy. It whispers. It knows your name.} `[MAGIC]` `[CURSE]`
\end{enumerate}

\section*{Quick use notes}
\label{sec:pereshi-quick-use}
\begin{itemize}
\item Draw until you have all four suits: Spade = place, Heart = actor, Club = pressure, Diamond = leverage. Highest rank sets the main Clock (2–5 → 4, 6–10 → 6, J/Q/K → 8, A → 10).
\item Diamonds are codified outcomes (tokens, scrolls, coals) that change position rather than call for a roll.
\item If any A appears, echo \textbf{fire‑omens} – flames that burn green, scrolls that ignite without heat, a bridge that recites a verse in a dead language, a coal that glows in the dark. Also tick the Prince's Shadow clock by 1 (the curse advances). The coal is not a metaphor. It is warm.
\end{itemize}

\subsection*{Additional Features (Cultural Mechanics)}
\label{sec:pereshi-mechanics}

\begin{itemize}
\item \textbf{Fire Temples and the Ninth Coal:} Every fire temple keeps a brazier with nine coals. The ninth coal is never lit – it is a promise to the Hollow and a ward against the Prince. Lighting it summons the Oath‑Keeper or the Prince's ghost. Flame‑Keepers inherit the secret of which coal is the ninth; it changes with each keeper. The coal is warm to the touch. It should not be. Recently, all the ninth coals across the highlands have grown warm at the same time. They pulse in the same rhythm. The rhythm is not a heartbeat.
\item \textbf{Verse Law:} In Pereshi courts, a well‑recited verse carries more weight than a witness. The old epics contain every legal precedent; a litigant who can quote the correct verse wins automatically. The problem is that the epics are incomplete – and the missing verses are hidden in the tombs of forgotten kings. The Prince tried to finish the epic with his own verses. They are still there, in the sealed archive, waiting to be burned – or to be sung. The singers who try to learn them lose their voices for a year.
\item \textbf{Caravan Tariffs as Currency:} A paid tariff receipt can be traded, sold, or used as proof of credit. The Pereshi have no central bank; the receipt is the coin. Forged receipts are a capital crime – the forger is forced to copy the entire epic by hand, blindfolded. The Prince's tax collectors were the first to be executed this way. Their ghosts still copy, still blindfolded. The ink they use is rust‑coloured.
\item \textbf{The Prince's Seal (Forgotten):} In the deepest vault of the archive, sealed in a lead box, is the Prince's own signet ring. It is said to grant authority over any Pereshi – because it proves the Prince's claim. No one has opened the box. The key was buried with the last emperor. The ring is still warm. Recently, the box has begun to hum. The hum is the same pitch as the ninth coal's pulse.
\item \textbf{Name‑Debt Mechanics:} Instead of tracking coin, a Pereshi may demand a “name” – the true name of an enemy, a forgotten ancestor, a secret that has never been spoken. A character who cannot pay loses one Bond or one Asset (the name is erased from records and memory). The GM may offer a choice: pay a name you know, or lose one of yours. This is not a debt. It is an erasure.
\item \textbf{Fire‑Temple Hospitality:} The three‑day feast rule is sacred. Breaking it marks the offender as \textit{anathema} – all Pereshi refuse to speak to them, and the ninth coal in every hearth spits ash in their direction. Removing the mark requires a pilgrimage to Mount Khvarena and the offering of a true name to the ninth coal. The coal may accept or refuse. It usually refuses.
\item \textbf{Local Patrons (Syncretic Names):}
  \begin{itemize}
    \item \textbf{The Undying Flame (Oath of Flame and Light):} The fire that witnessed the Prince's defeat – and still burns in witness and in warning. The warning is the coal's warmth.
    \item \textbf{The Keeper of Sealed Names (The Witness):} The archivist who records what must never be spoken – including the Prince's true name. She has started to forget. The forgetting is a mercy.
    \item \textbf{The Ninth Coal (The Ninth):} The coal that was never lit; the promise that the Prince will not return. It is warm. It is warm. It is warm.
    \item \textbf{The Silent Smith's Patron (Mykkiel):} The law folded into steel; the covenant that cannot be broken – unless the Prince's shadow crosses it. The shadow is crossing it now.
  \end{itemize}
\end{itemize}

\subsection*{Lore Echoes (Cross‑Expansion Hooks)}
\label{sec:pereshi-lore}

\begin{itemize}
  \item \textbf{Fire Temples and the Covenant (Mykkiel):} Covenant judges recognise the authority of flame‑keepers. A fire‑temple token can be used as a witness in a Covenant court – the flame itself is considered an incorruptible witness. The Prince once stood before a Covenant judge. He lied. The flame flickered, but did not burn. The judge was paid. The judge's ghost now tends a temple in the high passes, still waiting for the truth. The flame flickers when the judge's ghost speaks. It flickers in the same rhythm as the coal.
  \item \textbf{The Ninth Verse and the Ukwe (Sekogo):} The missing verse of the Pereshi epic is said to be carved on the Ukwe's oldest root. The Lethai‑ar know the verse but will not recite it – doing so would call the Hollow. The verse is the Prince's own confession. The Ukwe's root has begun to bleed sap the colour of rust. The sap smells of the Prince's garden.
  \item \textbf{Silent Smith's Swords and the Aeler (Ngomebe):} The silent smiths are said to be descendants of Aeler vent‑priors who stayed in the highlands. Their folding technique is a lost art; the Ngomebe masons would pay a fortune for a working forge. The smiths will only trade for a Prince's relic – and the Prince's shadow is now offering them his own chains. The chains are warm. They pulse.
  \item \textbf{The Buried Archive and the Kalrantha Amulet (Dhahara):} The First Emperor's scrolls contain a description of the amulet's creation. The Veiled Sisters have been searching for the archive for a century. The party may find it first – or find it already looted. The Prince's ghost is also looking. He believes the amulet can restore his name. He is wrong. The amulet can only give him a new one. He does not want a new one.
  \item \textbf{The Prince's Ghost (Red Desert Connection):} Some say the Prince did not die. He walked into the Red Desert and became part of it. His ghost haunts the high passes on stormy nights, asking for a cup of water. Do not give it to him. He will ask for your name next. Lately, he has started asking for a shadow to borrow. He returns them. They are never quite the same.
\end{itemize}

\subsection*{Pereshi Superstitions (burn for +1 die once)}
\begin{itemize}
  \item \textbf{Bread to the Fire:} Crumble a crust into the temple brazier. Ignore the first \emph{Name‑Debt Called} penalty. \hfill `[SUPERSTITION]`
  \item \textbf{Knot the Scroll:} Tie a single ribbon around a book you are returning. Cancel \emph{Forged Scroll} or take –1 die on your next legal roll (your choice). \hfill `[SUPERSTITION]`
  \item \textbf{Name to the Bridge:} Whisper a dead poet's name into the water below. Gain +1 Effect when reciting a verse this scene, but mark Obligation to the Bridge's spirits. \hfill `[SUPERSTITION]`
  \item \textbf{Coal to the Shadow:} Drop a glowing coal into a basin of water while speaking the Prince's forbidden name backwards. The next time the Prince's shadow appears, it will hesitate for one exchange – but it will remember your face. It will remember your face in your dreams. \hfill `[SUPERSTITION]`
\end{itemize}

\subsection*{Thematic SB Spend Table}
\label{sec:pereshi-sb}

\paragraph{Minor Complications (1 SB)}
\begin{itemize}
\item \textbf{Exposure:} Your actions draw unwanted attention from \textbf{fire‑keepers or archivists}.
\item \textbf{Noise:} Sounds of your actions alert nearby \textbf{caravan guards or Tulkani traders}.
\item \textbf{Trace:} Evidence of your passage marks your route for \textbf{Prince's shadow or name‑hunters}.
\item \textbf{Delay:} A brief but meaningful setback costs you \textbf{time or safe passage through a pass}.
\item \textbf{Supply Strain:} Mark +1 segment on a relevant \textbf{resource clock}.
\end{itemize}

\paragraph{Moderate Setbacks (2 SB)}
\begin{itemize}
\item \textbf{Alarm Raised:} \textbf{Local Flame‑Keeper or lord's herald} becomes aware and begins responding.
\item \textbf{Position Lost:} You lose advantageous ground/cover/stealth due to \textbf{ashfall or a bridge recitation failure}.
\item \textbf{Foe Appears:} A \textbf{name‑hungry spirit or a Fhara debt collector} arrives on scene.
\item \textbf{Gear Trouble:} A piece of equipment becomes \textbf{Compromised/Neglected}.
\item \textbf{Lock/Barrier:} A simple obstacle now requires a test to overcome.
\end{itemize}

\paragraph{Serious Trouble (3 SB)}
\begin{itemize}
\item \textbf{Reinforcements:} Additional \textbf{fire‑temple guards, Pereshi riders, or Prince's shadow fragments} arrive.
\item \textbf{Key Gear Breaks:} A crucial tool/weapon becomes temporarily unusable.
\item \textbf{Major Twist:} The situation fundamentally changes – \textbf{name debt called / ninth coal cracks / archive floods}.
\item \textbf{Rail Tick:} Advance a relevant campaign/front clock by 1 segment.
\item \textbf{Condition Applied:} Mark \textbf{Fatigue 1/Harm 1/Condition} appropriate to fiction.
\end{itemize}

\paragraph{Major Turns (4+ SB)}
\begin{itemize}
\item \textbf{Trap Springs:} A prepared danger activates with full effect.
\item \textbf{Authority Arrival:} \textbf{Ninth Flame‑Keeper, Archivist Kaelen, or the Prince's shadow} intervenes.
\item \textbf{Scene Shift:} The environment changes dramatically – \textbf{temple burns / archive collapses / ninth coal shatters}.
\item \textbf{Patron Omen:} Divine/arcane forces take notice – \textbf{omen appears / blessing lost / curse manifests}.
\item \textbf{Narrative Pivot:} The story takes an unexpected turn that reframes objectives.
\end{itemize}

\paragraph{Region-Specific SB Options}
\begin{itemize}
\item \textbf{Pereshi (Fire Law):} Flames change colour, scrolls ignite without heat, a coal pulses out of rhythm.
\item \textbf{Pereshi (Name Law):} A character forgets their own name for a moment, a ghost asks who you are, a written word turns to ash.
\item \textbf{Pereshi (Bridge Law):} A bridge recites a verse from a dead language, a plank snaps at the ninth step, the wind carries a warning.
\end{itemize}

\subsection*{Pursuit Ladder (Near / Mid / Far)}
\begin{itemize}
\item \textbf{Advance/Withdraw (Wits + Survival):} On a hit, shift one band. On a strong hit, also impose \emph{Ash Cloud}.
\item \textbf{Cut the Name (Presence + Sway):} On a hit, freeze range for one exchange; if a rose from the Ruined Garden is in your possession, +1 die.
\item \textbf{Weather the Shadow (Resolve + Spirit):} On a hit in \emph{Prince's Shadow Points}, you pick who is targeted; on a miss, the GM does.
\end{itemize}

\subsection*{Boss Seeds (Ash, Name \& Heart)}
\begin{itemize}
  \item \textbf{The Name‑Eater (Nine Erasures)} – \emph{Phases}: Name Debt Called $\rightarrow$ Orchestrated Forgetting $\rightarrow$ Archive Coup. \emph{Cracks}: return a severed name to its owner, expose a false erasure, survive a night in the Archive of Severed Names.
  \item \textbf{The Ninth Flame‑Keeper (Legend)} – \emph{Phases}: Coal Pulses $\rightarrow$ Shadow Stretches $\rightarrow$ The Heart Cracks. \emph{Cracks}: offer a new name to bind (a willing sacrifice), speak the Prince's true name in the Well of Unspoken Names, or convince the Pereshi that the Prince's debt is finally paid.
  \item \textbf{The Prince's Shadow (Manifestation)} – \emph{Phases}: Whisper $\rightarrow$ Demand $\rightarrow$ Possession. \emph{Cracks}: burn a rose from the Ruined Garden in his presence, speak his true name while holding the ninth coal, or offer him a shadow in exchange for his own – but the shadow must be willing.
\end{itemize}

\begin{tcolorbox}[colback=black!3,colframe=black!40!white,title={Patronage \& Power}]
Among the Pereshi, power flows through fire, ink, and the weight of ancient witness. A lord's authority rests on the age of his library; a flame‑keeper's influence on the purity of her fire. The petty lords are constantly forging genealogies, and the archivists are constantly exposing them. Power is not inherited – it is cited, and it is witnessed by the flame. But now a new power stirs: the Prince's shadow, offering forgotten knowledge, impossible bargains, and the chance to rewrite history. Some Pereshi have begun to listen. The ninth coal is warm. The shadow is patient.

\textbf{For the GM:} Use the \textbf{Prince's Shadow Clock} as the campaign's ticking heart. Every session, describe how the ninth coal feels warmer, how a mirror in a lord's hall has begun to show a second figure, how a scribe's shadow now moves independently. Let the party decide whether to reseal the Prince, bargain with him, or – most dangerously – try to seize his power for themselves. The Pereshi are not a people waiting for a saviour. They are a people waiting for the fire to decide. And the fire is listening. The coal is warm. The shadow is stretching.

Remember the relationships: the Fhara send envoys who cannot meet Pereshi eyes; the Tulkani leave knots at thresholds and wait for the smoke. The lowlands do not see the Pereshi at all. That is how the Pereshi like it. But when the coal cracks, everyone will see – and it will be too late.
\end{tcolorbox}

% Index entries
\index{Pereshi|see{Roof of the Way}}
\index{Theme generator}
\index{Roof of the Way generator}
\index{Places!Pereshi}
\index{People!Pereshi}
\index{Complications!Pereshi}
\index{Rewards!Pereshi}
\index{Flame-Keeper}
\index{Forgotten Prince}
\index{Ninth coal}
\index{Name erasure}
\index{Fhara (fear)}
\index{Tulkani (unease)}
\index{Mount Khvarena}
\index{Archive of Severed Names}
\index{Prince's Shadow clock}
\index{Local Patrons!Pereshi}
\newpage

% -------------------- KUVANI --------------------
\chapter{Kuvani}
\emph{The Remount Steppe.}

\noindent
Sarnai says: \textit{“The Kuvani do not ride horses. They borrow them. A horse’s name is its lineage. Its lineage is its story. I told a remount master a story about a wolf who forgot her pack. He gave me a grey mare. The mare knew the story already.”}

\noindent
The Kuvani ride the eastern steppe beneath the Eternal Blue Sky. Kuva, the sky father, sees all: every oath sworn in the open, every horse born, every blade of grass bent by the wind. They have no temples – the sky is their shrine. Their shamans, called Wind‑Singers, read the future in the flight of hawks and the whisper of grass. Their dead are not buried; they are given to the sky on funeral platforms, that Kuva may carry their souls to the ancestral herd that rides the northern lights. A horse is not a mount; it is a relative. Its name is a lineage stretching back to the first mare that Kuva breathed life into.

\vspace{0.5em}
\noindent\textit{“Never ride a Kuvani horse without asking its name. The horse will answer by throwing you. If it throws you, you have insulted its lineage.”} \\
— Remount Master Alp Arslan

\vspace{1em}
\begin{almanac}
\textbf{What do people eat for breakfast?} Fermented mare’s milk and dried meat – the milk from the herd, the meat from the hunt. The hunt is a ceremony. The ceremony is a debt to the grass.

\textbf{What does a child learn before they can walk?} To ride. Placed on a horse before they can stand. The horse is patient. It remembers the child’s grandfather.

\textbf{What is the first thing you notice about a stranger?} Their saddle. Worn means many miles. New means many miles ahead. No saddle means a ghost. Ghosts do not need saddles.

\textbf{What is the most common crime?} Horse theft – a declaration of feud. A stolen horse is a stolen relative. The feud lasts until the horse is returned or the thief’s line is ended.

\textbf{What does no one talk about but everyone knows?} The ghost herd is real – spectral horses that ride the salt pan on moonless nights. They are the souls of riders who died without honour. They ride forever, searching for a battle that will never come.

\textbf{Where do people go when they die?} To the sky. Placed on funeral platforms, wrapped in felt. The birds take the flesh. The wind takes the bones. The sky takes the soul.

\textbf{How do you apologize for a serious wrong?} You offer a horse – the best in your herd. If the wronged party accepts, the debt is paid. If they kill the horse, the debt is blood.

\textbf{What is the most important festival?} The Ride of the Ancestors (autumn equinox). The clans ride to the salt pan. The ghost herd appears at midnight. The living ride with the dead.

\textbf{What is the secret that could destroy a powerful person?} The Khagan’s horse is a ghost. It was killed in battle twenty years ago. The Khagan does not know. He rides a skeleton wrapped in hide. The horse does not tire, does not eat, does not forgive.
\end{almanac}

\subsection*{Etiquette Hooks}
\begin{itemize}
  \item \textbf{Guest Tent Rope}: Leave a rope coiled outside to signal neutrality → +1 Position on first Parley in camp.
  \item \textbf{Bowl Before Blade}: Accept a simple bowl before speaking of arms → Audience: Cool → Warm with elders.
  \item \textbf{Honour the Horse}: Praise mounts; tend tack before sleep → DV –1 on next Chase related to horse teams.
  \item \textbf{Faux Pas (SB)}: Spurned Mount – neglecting a horse lets the GM impose Position –1 among riders.
\end{itemize}

\subsection*{Strings (sample)}
\begin{itemize}
  \item \textbf{Remount Chain Mark} – a horse’s tooth; exchange for a fresh mount at any Kuvani station.
  \item \textbf{Grazing Compact} – sheepskin document; safe passage on clan land.
  \item \textbf{Sky‑Road Song} – a melody on a bone flute; navigate any storm.
\end{itemize}

\subsection*{Venue Tags}
\begin{itemize}
  \item \textbf{Remount Post}: swap fresh mounts; once/leg treat Chase (horse) DV –1.
  \item \textbf{Sky‑Road Cairns}: waystones; spend a song token to convert Navigate Position +1 across open grass.
\end{itemize}

\newpage
\section{Kuvani --- ``The Remount Steppe''}
\label{chap:kuvani}

\subsection*{Quick Reference}
\textbf{Tagline:} \textit{The sky hears every oath. The grass remembers every step.}

\begin{quote}
  \textbf{Wind-Singer (Elite):}  
  ``Kuva watches from the blue. He does not speak. He listens. A promise made under open sky is carved into the wind itself. The grass will bend around it. The horses will remember it. Break it, and the sky will turn its face away. You will ride forever in shadow.''
\end{quote}

\begin{quote}
  \textbf{Remount Master (Commoner):}  
  ``A horse's name is its lineage. Its lineage is its story. You do not ride a Kuvani horse. You borrow it. And when you return it, you will tell me a story that the horse has not already heard. If you lie, the horse will throw you next time — and Kuva will laugh.''
\end{quote}

\textbf{Genre / Mood:} Sky‑worshipping nomads, ancestral horse‑magic, wind‑carried oaths.

\textbf{Starting Location:} A felt tent circle at a remount station, where a corral of horses stands behind a rope of braided hair. A Remount Master sits on a saddle, a bell in his hand: ``You need fresh horses. I have twelve. Each horse has a name. Each name has a price. The price is not silver — it is a story I have not heard. Tell me one. If it makes me laugh, you ride for free. If it makes me weep, you owe me a foal. If Kuva sends a hawk to watch, you owe a prayer.''

\vspace{2mm}
\begin{tcolorbox}[colback=black!3,colframe=blue!60!black,title={Theme \& Atmosphere}]
The Kuvani ride the eastern steppe beneath the Eternal Blue Sky. Kuva, the sky father, sees all: every oath sworn in the open, every horse born, every blade of grass bent by the wind. The Kuvani have no temples — the sky is their shrine. Their shamans, called Wind-Singers, read the future in the flight of hawks and the whisper of grass. Their dead are not buried; they are given to the sky on funeral platforms, that Kuva may carry their souls to the ancestral herd that rides the northern lights.

Unlike their Ykrul neighbours, who worship geometry and oath‑strings, the Kuvani worship movement itself. A horse is not a mount — it is a relative. Its name is a lineage stretching back to the first mare that Kuva breathed life into. A Kuvani who cannot name his horse's ancestors is a Kuvani without a shadow. The steppe is not empty; it is full of ghosts: the ghost herd of the salt pan, the sky‑road markers that hum at dusk, the Iron Corral where blades whisper the names of the dead. A Kuvani does not fight for land — he fights for grazing rights, for the memory of grass, for the honour of his horse's line.

\vspace{2mm}
\textbf{What you notice first:} The smell of horse sweat and dry grass. The way every rider touches a small felt bundle at their belt — a piece of a dead horse's mane, a prayer to Kuva. The constant sound of bells on harnesses, each tone a different ancestor's name.

\textbf{The one rule that kills newcomers:} Never ride a Kuvani horse without asking its name. The horse will answer — by throwing you. If it throws you, you have insulted its lineage. The remount master will demand a story as compensation. If the story is not good enough, you walk — and Kuva will turn the wind against you.
\end{tcolorbox}

\subsection*{Regional Mood: The Weight of the Herd}

Power here is measured in remount chains, grazing compacts, and the favour of Kuva. A clan's status depends on the number of horses it can offer to caravans; a remount master's reputation on the speed and health of his string. The Kuvani have no kings — only Wind-Singers, who speak for the sky, and Remount Masters, who hold the earthly ledger of horse‑flesh. Every caravan that crosses the steppe must pay the remount toll: not in coin, but in stories. The Kuvani believe that a story shared is a prayer heard, and that a traveller who has no story has no soul — only an empty shape that the wind will soon erase.

\subsection*{Prominent NPCs of the Kuvani}
\label{sec:kuvani-npcs}

\begin{longtable}{|p{0.2\textwidth}|p{0.25\textwidth}|p{0.3\textwidth}|p{0.15\textwidth}|}
\hline
\textbf{Name} & \textbf{Role/Title} & \textbf{Motivation} & \textbf{Rumoured Quirk} \\
\hline
Alp Arslan & Remount master of the central route & Expand grazing rights into Ykrul territory & He has never been seen sleeping; his eyes are always on the horizon, watching for Kuva's hawks. \\
\hline
Borte & Wind-Singer of the White Mane Clan & Keep the songs of the clans from being forgotten & She has a horse with a white blaze; the blaze is a map of the steppe, and it shifts each season. \\
\hline
Temur & Caravan guide & Find the fastest routes through the shifting grass & He names his dogs after dead enemies. The dogs remember. They growl when the enemy's ghost is near. \\
\hline
Seljuk & Ykrul envoy (negotiator) & Negotiate shared pastures without war & He carries a hostage string with nine knots; each knot is a season of peace. He is running out of knots. \\
\hline
The Grass Reader & A hermit who lives in a felt hut & Read the future in the way the grass bends & He never speaks above a whisper. The grass speaks for him. His tongue is green. \\
\hline
\end{longtable}

\subsection*{Names of the Kuvani}
\label{sec:kuvani-names}

Inspired by Turkic and Mongolian traditions: short, hard consonants, vowel clusters, patronymic suffixes (-oglu, -ova). Surnames denote a clan, a horse-line, or a wind-name.

\begin{longtable}{|p{0.25\textwidth}|p{0.25\textwidth}|p{0.2\textwidth}|p{0.2\textwidth}|}
\hline
\textbf{First Names (Male)} & \textbf{First Names (Female)} & \textbf{Surnames} & \textbf{Epithets} \\
\hline
Alp & Borte & Arslanoglu & \textit{the Iron Rein} \\
\hline
Temur & Selin & Kuvani-bek & \textit{the White Mane} \\
\hline
Seljuk & Chagatai & Borjigin & \textit{Ninth Wind} \\
\hline
Batu & Nergui & Oirat & \textit{the Grass-Tongue} \\
\hline
Toghrul & Otgon & Kipchak & \textit{the Remount-Keeper} \\
\hline
\end{longtable}

\paragraph*{(Station/Circle/Kurgan)} Remount station with a felt tent and a corral; Wind-Singer's circle of saddles; a kurgan (burial mound) with a stone horse that is missing a leg.

\section*{Spades --- Places}
\label{sec:kuvani-places}
\begin{enumerate}
\setcounter{enumi}{1}
\item \textbf{Remount Station} --- A felt tent with a corral; horses for hire. The horses have no names — only brands.
  \hfill \textit{The brands are the names of lost children. The horses remember them. Kuva's hawks circle when a brand is new.} `[TRADE]`
\item \textbf{Wind-Singer's Circle} --- A ring of saddles; the singer sits in the centre. The songs are not sung — they are hummed.
  \hfill \textit{The hum is the wind. The wind carries the song to the next valley. The next valley answers with a hawk's cry.} `[RITUAL]`
\item \textbf{Kurgan of the Broken Oath} --- A burial mound with a stone horse missing a leg.
  \hfill \textit{The leg is a treaty. The treaty was broken. The horse is waiting for it to be returned. Its eye follows you.} `[DEATH]`
\item \textbf{Grass Maze} --- A section of the steppe where the grass grows taller than a rider on horseback.
  \hfill \textit{The paths shift with the wind. Only a Wind-Singer can find the way out. The grass whispers warnings.} `[MYSTERY]`
\item \textbf{Salt Pan of the Ghost Herd} --- A dry lake; the salt is crusted with hoofprints.
  \hfill \textit{The prints are from a herd that died a century ago. They appear every spring. The herd is still running. Kuva's hawks do not fly over it.} `[DEATH]`
\item \textbf{Felt Village (Mobile)} --- A temporary camp of round tents. The tents face the wind.
  \hfill \textit{The wind is the clan's direction. If the wind changes, the village moves. The door always faces the rising sun.} `[COMMUNITY]`
\item \textbf{Horse Cemetery} --- A field of skulls, each on a wooden stake. The stakes are carved with names.
  \hfill \textit{The names are the riders' names. The riders are buried elsewhere. The horses wait. Their skulls whicker at dawn.} `[DEATH]`
\item \textbf{Bridge of the Ninth Hardship} --- A log over a ravine. The log is replaced each year. The old log is burned.
  \hfill \textit{The ash is the clan's memory. The ash is scattered on the spring grass. Kuva drinks the smoke.} `[RITUAL]`
\item \textbf{Stone of the Hostage String} --- A standing stone with notches. Each notch is a hostage taken.
  \hfill \textit{The hostages are carved into the stone. The stone weeps when a hostage dies. The tears taste of salt and iron.} `[DEBT]`
\item[J] \textbf{The Iron Corral} --- A pen made of salvaged sword blades. The blades are from a war that never ended.
  \hfill \textit{The horses are afraid of them. The blades whisper the names of the dead. The whisper is the wind.} `[MAGIC]`
\item[Q] \textbf{Sky-Road Marker} --- A cairn painted blue. The blue is from a stone that fell from the sky.
  \hfill \textit{The stone is still warm. It marks the path of the summer stars. Touch it, and you will dream of falling.} `[TRAVEL]`
\item[K] \textbf{Khagan's Felt Throne} --- A tent within a tent; the seat is a saddle. The saddle has no stirrups.
  \hfill \textit{The Khagan does not need them. He rides the wind. The wind has no stirrups. The saddle was his grandfather's.} `[AUTHORITY]`
\item[A] \textbf{The Steppe's Edge} --- Where the grass ends and the desert begins. The line is a scar.
  \hfill \textit{The scar is a treaty. The treaty was signed in blood. The blood was Kuvani. Kuva turned his face away that day.} `[LEGEND]`
\end{enumerate}

\paragraph*{(Master/Singer/Guide)} Remount master with calloused hands and a horsehair whip; Wind-Singer with grey hair and a bone flute; caravan guide with a string of dog teeth.

\section*{Hearts --- People \& Factions}
\label{sec:kuvani-people}
\begin{enumerate}
\setcounter{enumi}{1}
\item \textbf{Remount Master} --- Calloused hands, squinting eyes, a whip of braided horsehair.
  \hfill \textit{He knows the price of a horse in salt, in silk, or in stories. The price changes with the wind. He has never been cheated.} `[TRADE]`
\item \textbf{Wind-Singer (Shaman)} --- Grey-haired, blind, but never lost. She hums the direction to the nearest water.
  \hfill \textit{The hum changes with the wind. If she stops humming, the wind stops too. Her left eye sees the living; her right sees the dead.} `[RITUAL]`
\item \textbf{Caravan Guide (Temur)} --- Young, fast, a scar across his lip. He names his dogs after dead enemies.
  \hfill \textit{The dogs remember. They growl when they smell the enemy's kin. The dogs also see ghosts.} `[TRAVEL]`
\item \textbf{Ykrul Envoy (Seljuk)} --- Braided hair, a hostage string on his wrist. The string has nine knots.
  \hfill \textit{Each knot is a season of peace. He is running out of knots. The last knot is tied with a hair from a Khagan's daughter.} `[POLITICS]`
\item \textbf{Grass Reader (Hermit)} --- A hermit in a felt hut, never seen without a handful of grass.
  \hfill \textit{He reads the future by how the grass bends. The future is always wind. He has never been wrong.} `[DIVINATION]`
\item \textbf{Horse Thief} --- Young, missing two fingers. The fingers were a bet. He lost the bet — but not the horses.
  \hfill \textit{He can steal a horse from a corral without waking the dogs. The dogs are his. He paid for them with a story.} `[STEALTH]`
\item \textbf{Felt Village Elder} --- A woman with a pipe made of bone. She settles disputes by listening to the wind.
  \hfill \textit{``The wind says you are lying,'' she says. The wind has never been wrong. The pipe belonged to her grandmother.} `[JUSTICE]`
\item \textbf{Khagan's Shadow} --- A silent figure in black felt. He never speaks; he writes commands in the dust.
  \hfill \textit{The dust blows away. The commands are remembered anyway. His face is never seen.} `[AUTHORITY]`
\item \textbf{Horse-Whisperer (Child)} --- A child who can calm any stallion. The child speaks in neighs.
  \hfill \textit{The horses answer in riddles. The child is learning to understand. One day, a horse will answer with a name.} `[MYSTERY]`
\item[J] \textbf{Ghost Herd's Rider} --- A translucent figure on a skeletal horse, seen at the salt pan at dawn.
  \hfill \textit{He asks: ``Where is my herd?'' If you answer, you become part of it. If you refuse, he follows you for nine days.} `[DEATH]`
\item[Q] \textbf{The Steppe's Voice (Oracle)} --- A rock with a crack. The wind speaks through it.
  \hfill \textit{The wind answers questions. The answers are always poems. The poems are always warnings. The rock is warm.} `[DIVINATION]`
\item[K] \textbf{The Last Khagan (Legend)} --- A king who united the clans, then rode into a storm and never returned.
  \hfill \textit{He is said to appear during the worst blizzards, offering shelter. The shelter is a trap. His horse has no shadow.} `[LEGEND]`
\item[A] \textbf{The First Remount (Myth)} --- A horse that never tired. Its bones are scattered across the steppe.
  \hfill \textit{Find all nine bones, and you can summon it. The summoning costs a year of your life. The horse will remember your name.} `[MYTH]`
\end{enumerate}

\paragraph*{(Sickness/Fire/Stampede)} A horse sickness spreads through the remounts — all mounts lose 1 Speed; grass fire — the smoke spells a name; a wild horse stampede — the horses are fleeing something worse.

\section*{Clubs --- Complications/Threats}
\label{sec:kuvani-complications}
\begin{enumerate}
\setcounter{enumi}{1}
\item \textbf{Horse Sickness (Curse)} --- A cough spreads through the remounts. The cough is a curse from a Ykrul bone-singer.
  \hfill \textit{The cure is a story that the bone-singer has not heard. The story must be true. Kuva's hawks will judge.} `[CURSE]`
\item \textbf{Grass Fire} --- The steppe burns; the smoke is black. The smoke spells a name.
  \hfill \textit{The name is the one who lit the fire. The wind will carry the name to the clans. They will hunt. The name also appears in your dreams.} `[DISASTER]`
\item \textbf{Stampede (Wild Herd)} --- A herd of wild horses crashes through camp. The horses are not wild.
  \hfill \textit{They are fleeing something worse. Something that hunts horses. The party is in its path. The ground shakes.} `[COMBAT]`
\item \textbf{Hostage String Cut} --- A Ykrul envoy's string is broken. The clan demands a new hostage.
  \hfill \textit{The party must choose a member to offer. The hostage will be held for nine seasons. Kuva will witness the oath.} `[DEBT]`
\item \textbf{Rasputitsa (Mud)} --- The spring thaw turns the steppe to mud. Wagons sink. Progress slows.
  \hfill \textit{The mud is warm. Something is buried underneath. It is waking up. The mud smells of old blood.} `[WEATHER]`
\item \textbf{Wind-Singer's Silence} --- The Wind-Singer stops humming. The wind stops too. The silence is a warning.
  \hfill \textit{The warning is: ``A storm is coming. Not of rain — of ghosts.'' The ghosts are the unburied.} `[DEATH]`
\item \textbf{Water Poisoned (Salt Pan)} --- A well is found with dead animals nearby. The poison is salt from the cursed pan.
  \hfill \textit{The salt pan wants a sacrifice. The sacrifice is a horse. The horse must be willing. The horse will neigh its acceptance.} `[CURSE]`
\item \textbf{Felt Village Gone} --- A camp is empty; the tents are still standing. The felt is warm.
  \hfill \textit{The people vanished at noon. The horses are still in the corral. The horses are trembling. Their eyes are white.} `[MYSTERY]`
\item \textbf{Khagan's Demand} --- The Khagan (or his shadow) asks for a tribute. The tribute is a story.
  \hfill \textit{The story must be about a broken promise. If it is not good enough, the Khagan takes a horse. The horse will never be seen again.} `[DEBT]`
\item[J] \textbf{Iron Corral Bleeding} --- The sword-blade fence begins to weep rust. The rust is blood.
  \hfill \textit{The blood is from the war that never ended. The war is about to resume. The blades are warm.} `[OMEN]`
\item[Q] \textbf{Sky-Road Marker Falls} --- The blue cairn collapses. The stone inside is a meteor. It whispers.
  \hfill \textit{The whisper is a prophecy. The prophecy is about the end of the steppe. The stone is cold.} `[OMEN]`
\item[K] \textbf{Ghost Herd Charges} --- Spectral horses trample the camp. They are looking for their riders.
  \hfill \textit{The riders are dead. The ghosts do not know. One rider's ghost is in your party. It does not know it is dead.} `[DEATH]`
\item[A] \textbf{Steppe's Edge Recedes} --- The grass grows into the desert. The desert is retreating.
  \hfill \textit{Something is driving the desert back. Something that wants the steppe. Kuva is silent.} `[APOCALYPSE]`
\end{enumerate}

\paragraph*{(Mark/Compact/Song)} Remount chain mark (a horse's tooth, exchange for a fresh mount); grazing compact (a sheepskin document, safe passage on clan land); sky-road song (a melody on a bone flute, navigate any storm).

\section*{Diamonds --- Rewards/Leverage}
\label{sec:kuvani-rewards}
\begin{enumerate}
\setcounter{enumi}{1}
\item \textbf{Remount Chain Mark} --- A horse's tooth on a leather thong. Exchange it for a fresh mount at any Kuvani station.
  \hfill \textit{The tooth is from a horse that died in battle. Its spirit runs with the new mount. The new mount will be braver.} `[TRAVEL]`
\item \textbf{Grazing Compact (Sheepskin)} --- A written agreement. Allows your caravan to pasture on clan land for one season.
  \hfill \textit{The compact is sealed with horsehair. The hair is from the Khagan's own mare. Break the seal, and the mare will know.} `[TRADE]`
\item \textbf{Sky-Road Song (Bone Flute)} --- A melody recorded on a flute. Play it to navigate any storm.
  \hfill \textit{The flute was carved from a rib. The rib was the Wind-Singer's. She gave it willingly. The flute is warm.} `[TRAVEL]`
\item \textbf{Hostage String (Borrowed)} --- A braided cord with knots. Present it to demand a truce from a Ykrul clan.
  \hfill \textit{The string is hair. The hair is from a Khagan's daughter. The Ykrul will recognise it. They will honour it. Once.} `[SOCIAL]`
\item \textbf{Felt Tent Pass} --- A felt circle with a hole. Safe lodging in any Kuvani camp.
  \hfill \textit{The hole is for smoke. The smoke is your story. The clan will listen. If they laugh, you stay. If they weep, you owe a foal.} `[SOCIAL]`
\item \textbf{Kurgan Key (Stone Tablet)} --- A flat stone. Opens the Horse Cemetery gate.
  \hfill \textit{The gate is a skull. The key fits the eye socket. The skull will scream. The scream is a name — the name of the horse you must bury.} `[MAGIC]`
\item \textbf{Iron Corral Nail} --- One sword blade from the fence. Use it to seal an oath.
  \hfill \textit{The blade is rusty. The rust is from broken promises. It will rust your oath too. Speak true, or bleed.} `[OATH]`
\item \textbf{Sky-Road Stone (Chip)} --- A piece of the blue cairn. Carry it to avoid getting lost on the steppe.
  \hfill \textit{The stone is warm. It points to the north. The north is where the grass is greenest. The stone hums at dawn.} `[TRAVEL]`
\item \textbf{Khagan's Whistle (Horse Phalange)} --- A brass whistle. Blow it to summon a remount from any station.
  \hfill \textit{The whistle is a horse's toe bone. The horse was the Khagan's favourite. It will come — but only once.} `[MAGIC]`
\item[J] \textbf{Ghost Herd's Bridle (Leather)} --- A strap of old leather. Use it to command one spirit horse.
  \hfill \textit{The bridle is from a horse that drowned in the salt pan. The spirit is still wet. It will carry you across water.} `[DEATH]`
\item[Q] \textbf{Wind-Singer's Staff (Tent Pole)} --- A wooden rod. Plant it to stop the wind for one hour.
  \hfill \textit{The staff was a tent pole. The tent was the singer's home. The home is gone. The wind will listen — for one hour.} `[WEATHER]`
\item[K] \textbf{Khagan's Saddle (Felt Pad)} --- Ride in it to gain +1 Position on all mounted rolls.
  \hfill \textit{The saddle has no stirrups. You must learn to ride without them. The Khagan did. His ghost rides behind you.} `[COMBAT]`
\item[A] \textbf{Steppe's Edge Stone} --- A piece of the boundary where grass meets desert. Touch it to make the grass grow.
  \hfill \textit{The stone is a tooth. The tooth is from the Steppe itself. The Steppe is hungry. The grass will grow — but it will be thirsty.} `[LEGEND]`
\end{enumerate}

\section*{Quick use notes}
\label{sec:kuvani-quick-use}
\begin{itemize}
\item Draw until you have all four suits: Spade = place, Heart = actor, Club = pressure, Diamond = leverage. Highest rank sets the main Clock (2–5 → 4, 6–10 → 6, J/Q/K → 8, A → 10).
\item Diamonds are codified outcomes (marks, compacts, flutes) that change position rather than call for a roll.
\item If any A appears, echo \textbf{sky-omens}—hawks that circle counter‑clockwise, grass that bends against the wind, a bell that rings without being struck, a shadow that falls where no sun shines.
\end{itemize}

\subsection*{Additional Features (Cultural Mechanics)}
\label{sec:kuvani-mechanics}

\begin{itemize}
\item \textbf{The Story Toll:} Kuvani remount stations do not accept coin. To take a fresh horse, you must tell a story that the remount master has not heard. The story is judged by the horse — if the horse stamps its foot, the story is false. If it whinnies, the story is true. If it remains still, the story is boring — you walk. Kuva's hawks watch; a hawk's cry means the story is worthy of the sky.
\item \textbf{Wind-Songs and Sky-Roads:} The Kuvani navigate by the stars and the wind. A Wind-Singer can hum a melody that encodes the route to water, to grass, or to safety. The song changes with the season. Only the singer knows the current melody. Those who hear the song without permission are cursed — they will hear it in their dreams until they pay a horse to the singer.
\item \textbf{Grazing Compacts as Currency:} A grazing compact (a sheepskin document sealed with horsehair) can be traded, sold, or used as proof of passage. A compact from a powerful clan is worth more than a chest of silver. Forging a compact is punishable by exile — the forger is left on the steppe without a horse. Kuva will not watch over them.
\item \textbf{Kuva's Witness (Oath Mechanic):} Any oath sworn under open sky (no roof, no canopy) is witnessed by Kuva. Breaking such an oath imposes −1 die on all rolls until the oath is fulfilled or atoned for by a sacrifice: a horse, a memory, or a year of service to the wronged party. The GM may also tick a \textbf{Kuva's Attention [6]} clock; when full, a sky-omen (drought, storm, ghost herd) strikes the oath‑breaker.
\item \textbf{Local Patrons (Syncretic Names):}
  \begin{itemize}
    \item \textbf{Kuva} — The Eternal Blue Sky. *The witness of all oaths, the father of horses, the wind that carries prayers.*
    \item \textbf{The White Mare} — The Traveler. *The first horse, who showed the Kuvani the sky‑road. Her hoofprints are the stars.*
    \item \textbf{The Iron Hoof} — Khemesh. *The pressure beneath the steppe; the hunger that pulls horses down into the salt pan.*
    \item \textbf{The Shadow Khagan} — Malachai. *A cursed king who sold his clan's horses for immortality; he rides alone, and his horses have no shadows.*
  \end{itemize}
\end{itemize}

\subsection*{Lore Echoes (Cross-Expansion Hooks)}
\begin{itemize}
  \item \textbf{Remount Chains and the Caravan Trade (Way of Silk):} Every major caravan on the eastern route uses Kuvani remount stations. A disruption in the remount chain — a sickness, a feud, a stolen horse — can halt trade for an entire season. The Tulkani have a saying: ``A Kuvani horse is worth nine stories. A dead Kuvani horse is worth nine feuds.''
  \item \textbf{Grazing Compacts and the Ykrul:} The Ykrul are the Kuvani's rivals and sometimes allies. A grazing compact with a Kuvani clan may require a Ykrul hostage string as surety. Breaking the compact means the hostage is executed — and Kuva will turn his face from the oath‑breaker. The Ykrul do not fear Kuva, but they respect his silence.
  \item \textbf{Wind-Singers and the Sky (Kuva):} The Kuvani worship Kuva as the watcher of oaths. Their wind-songs are prayers to the sky. A Wind-Singer who loses her voice is considered cursed by Kuva; she must walk the steppe alone until she hears the sky again. Some never return; their bones become sky‑road markers.
  \item \textbf{The Ghost Herd and the Hollow (Eastern Realms):} The ghost herd at the salt pan is said to be the Hollow's own horses — ridden by the Oath-Keepers. Aelaerem leave offerings of salt to keep the herd from riding through their villages. The Kuvani avoid the salt pan entirely. They say the ghost herd is Kuva's punishment for a Khagan who broke an oath to the sky.
  \item \textbf{The Iron Corral and the Black Banners:} The sword blades of the Iron Corral are from a war that never ended — a condotta that was paid twice, first in coin, then in blood. The Black Banner captains know the story; they leave offerings of rust at the corral when they pass. A captain who refuses to pay finds his horses spooked and his supply lines cut.
\end{itemize}

\subsection*{Kuvani Superstitions (burn for +1 die once)}
\begin{itemize}
  \item \textbf{Bread to the Horses:} Crumble a crust into the corral. Ignore the first \emph{Horse Sickness} penalty. Kuva's hawks will watch but not strike.
  \item \textbf{Knot the Hostage String:} Tie a new knot in a borrowed string. Cancel \emph{Hostage String Cut} or take –1 die on your next negotiation with Ykrul (your choice). The knot must be tied while facing the rising sun.
  \item \textbf{Name to the Wind:} Whisper a dead horse's name into the grass. Gain +1 Effect when riding this scene, but mark \emph{Obligation} to the Wind-Singer. The wind will carry the name to Kuva. He may answer with a storm.
\end{itemize}

\subsection*{Thematic SB Spend Table}
\label{sec:kuvani-sb}

\paragraph{Minor Complications (1 SB)}
\begin{itemize}
\item \textbf{Exposure:} Your actions draw unwanted attention from \textbf{Remount Masters or Wind-Singers}.
\item \textbf{Noise:} Sounds of your actions alert nearby \textbf{herdsmen or felt village guards}.
\item \textbf{Trace:} Evidence of your passage marks your route for \textbf{rival clans or ghost herd scouts}.
\item \textbf{Delay:} A brief but meaningful setback costs you \textbf{time or favourable grazing access}.
\item \textbf{Supply Strain:} Mark +1 segment on a relevant \textbf{resource clock}.
\end{itemize}

\paragraph{Moderate Setbacks (2 SB)}
\begin{itemize}
\item \textbf{Alarm Raised:} \textbf{Local Remount Master or Wind-Singer} becomes aware and begins responding.
\item \textbf{Position Lost:} You lose advantageous ground/cover/stealth due to \textbf{grass fire or sudden fog}.
\item \textbf{Foe Appears:} A \textbf{rival clan rider or Ykrul envoy} arrives on scene.
\item \textbf{Gear Trouble:} A piece of equipment becomes \textbf{Compromised/Neglected}.
\item \textbf{Lock/Barrier:} A simple obstacle now requires a test to overcome.
\end{itemize}

\paragraph{Serious Trouble (3 SB)}
\begin{itemize}
\item \textbf{Reinforcements:} Additional \textbf{Kuvani riders, Ykrul warriors, or ghost herd} arrive.
\item \textbf{Key Gear Breaks:} A crucial tool/weapon becomes temporarily unusable.
\item \textbf{Major Twist:} The situation fundamentally changes—\textbf{water poisoned/hostage string cut/khagan's demand issued}.
\item \textbf{Rail Tick:} Advance a relevant campaign/front clock by 1 segment.
\item \textbf{Condition Applied:} Mark \textbf{Fatigue 1/Harm 1/Condition} appropriate to fiction.
\end{itemize}

\paragraph{Major Turns (4+ SB)}
\begin{itemize}
\item \textbf{Trap Springs:} A prepared danger activates with full effect.
\item \textbf{Authority Arrival:} \textbf{Khagan, Wind-Singer of the White Mane, or Ghost Herd's Rider} intervenes.
\item \textbf{Scene Shift:} The environment changes dramatically—\textbf{steppe catches fire/salt pan floods/sky-road marker falls}.
\item \textbf{Patron Omen:} Divine/arcane forces take notice—\textbf{omen appears/blessing lost/curse manifests}.
\item \textbf{Narrative Pivot:} The story takes an unexpected turn that reframes objectives.
\end{itemize}

\paragraph{Region-Specific SB Options}
\begin{itemize}
\item \textbf{Kuvani (Remount Law):} Horse bells ring out of sequence, saddle leather creaks a warning, a horse refuses to move.
\item \textbf{Kuvani (Wind-Song):} Grass whispers a name, the wind stops mid-sentence, a bone flute plays itself.
\item \textbf{Kuvani (Ghost Herd):} Hoofprints appear where there are no horses, salt tastes of blood, a skeletal horse watches from the horizon.
\end{itemize}

\subsection*{Pursuit Ladder (Near / Mid / Far)}
\begin{itemize}
\item \textbf{Advance/Withdraw (Wits + Ride):} On a hit, shift one band. On a strong hit, also impose \emph{Wind at Back}.
\item \textbf{Cut the Herd (Presence + Sway):} On a hit, freeze range for one exchange; if a remount chain mark is in your possession, +1 die.
\item \textbf{Weather the Curse (Resolve + Spirit):} On a hit in \emph{Horse Sickness}, you pick which mount survives; on a miss, the GM does.
\end{itemize}

\subsection*{Boss Seeds (Hoof \& Wind)}
\begin{itemize}
\item \textbf{The Bone-Singer (Nine Plagues)} — \emph{Phases}: Horse Sickness Spread $\rightarrow$ Orchestrated Famine $\rightarrow$ Herd Coup. \emph{Cracks}: find a true story the bone-singer has not heard, expose a false cure, survive a stampede. \emph{Kuva's Judgment}: if the party swears an oath to end the sickness under open sky, the bone-singer loses one phase.
\item \textbf{The Storm Khagan (Rider of the White Blizzard)} — \emph{Phases}: Grass Fires Called Mercy $\rightarrow$ Procession of Empty Saddles $\rightarrow$ Wind-Baptism of a Remount Master. \emph{Cracks}: relic wind-song, survivor testimony, the sky-road turning against them. \emph{Kuva's Favour}: if a party member offers their own horse to the sky (releasing it onto the steppe), the Storm Khagan must pause for one exchange.
\end{itemize}

\begin{tcolorbox}[colback=black!3,colframe=black!40!white,title={Patronage \& Power}]
Among the Kuvani, power flows through remount chains, wind-songs, and the truth of a story. A clan's authority rests on the speed of its horses and the skill of its singers. The Khagan is a figurehead; the real power lies with the Remount Masters who control the stations, and the Wind-Singers who read the grass — and Kuva's will.

\textbf{For the GM:}  
Power in Kuvani society is measured in horse teeth, sheepskin compacts, and the memory of a good story. Patronage often takes the form of remount marks, grazing compacts, or bone flutes — each tied to a specific station, clan, or singer. To emphasize the itinerant and sky‑worshipping culture:
\begin{itemize}
\item Require players to \textbf{tell a story} to access remounts, shelter, or grazing rights. The GM can judge the story's quality (or have the player roll Presence + Sway vs. DV based on the remount master's mood). A failed story means the player must offer a truth — a secret, a confession, or a memory — to Kuva directly (speak it aloud under open sky).
\item Use \textbf{Wind-Singers} as quest-givers and as intermediaries with Kuva. They hear songs on the wind that only the party can resolve. The song is always a warning or a debt call, and often a prophecy.
\item Let \textbf{grazing compacts} be a form of currency that can be stolen, forged, or voided. A party without a compact is trespassing — the clans can demand a story as payment, or seize a horse as collateral. A compact sealed with a hair from the Khagan's mare is unbreakable without divine intervention.
\item Emphasize that \textbf{the steppe is alive and Kuva watches}: the grass remembers, the wind listens, and a broken promise is written in the sky-road. The Hollow's attention ticks when the party lies to a Kuvani. Kuva's attention ticks when they break an oath sworn under open sky.
\item Use \textbf{sky-omens} as scene-setting tools: a hawk's cry, a sudden stillness, a cloud shaped like a horse. These can foreshadow complications or reward truthful behaviour.
\end{itemize}
In Kuvani society, the grass forgets nothing — and Kuva's sky remembers every word spoken in the open.
\end{tcolorbox}

% Index entries
\index{Kuvani|see{Steppe Hosts}}
\index{Theme generator}
\index{Steppe Hosts generator}
\index{Places!Kuvani}
\index{People!Kuvani}
\index{Complications!Kuvani}
\index{Rewards!Kuvani}
\index{Remount stations}
\index{Wind-Singers}
\index{Grazing compacts}
\index{Horse sickness}
\index{Ghost herd}
\index{Sky-road}
\index{Iron corral}
\index{Khagan}
\index{Ykrul}
\index{Story toll}
\index{Kuva}
\index{Sky deity}
\index{Local Patrons!Kuvani}
\newpage

% -------------------- TULKANI --------------------
\chapter{Tulkani}
\emph{The Itinerant Courts.}

\noindent
Sarnai smiles: \textit{“My people have no home but the road. Our wagons are our hearths. Our knots are our memories. I tie a knot for every story I collect. When the cord is full, I will return to my grandmother’s circle and untie them, one by one, until the debt is paid.”}

\noindent
The Tulkani have no fixed homeland; they are wanderers who follow the seasons and the trade routes. Their wagons are their hearths, their circles are their courts, and their knots are their memories. They are welcome everywhere and trusted nowhere – a position they wear with wry humour. A Tulkani caravan can appear at a lord’s border, camp for seven nights, and leave before the eighth – because on the eighth night, the lord can demand rent. On the ninth night, something else comes. Their neutrality is sacred: a Tulkani camp is sanctuary for any traveller, regardless of allegiance, as long as they respect the circle and offer a story when asked.

\vspace{0.5em}
\noindent\textit{“Never sleep outside the circle at night. The wagon ring is sanctuary. Outside it, the Hollow walks. A Tulkani will not rescue you – they will only mourn you at dawn.”} \\
— Caravan Judge Sarnai (my grandmother)

\vspace{1em}
\begin{almanac}
\textbf{What do people eat for breakfast?} Tea and hardtack – the tea brewed from leaves gathered along the road, the hardtack baked at the start of the season. The hardtack is always stale. The tea is always fresh.

\textbf{What does a child learn before they can walk?} To tie a knot. The eight journey‑knots: Departure, Arrival, Promise, Debt, Payment, Witness, Apology, and the ninth – the knot that is never tied. The ninth knot is for the Hollow.

\textbf{What is the first thing you notice about a stranger?} Their cord. Every Tulkani carries a knotted cord – their life’s ledger. Many knots mean travelled far. Few knots mean just begun. No cord means a ghost.

\textbf{What is the most common crime?} Story‑theft – but story‑theft is a compliment. A stolen story is a story worth telling. The original teller is honoured. The thief is welcomed.

\textbf{What does no one talk about but everyone knows?} The Tulkani were exiled from their homeland nine centuries ago, for a crime that no one remembers. The crime is written in the ninth knot of the Elder’s cord. The cord is too long to read.

\textbf{Where do people go when they die?} To the road. Bodies left at a crossroads. The road takes them to the next valley, the next telling. The dead become stories. Stories are immortal.

\textbf{How do you apologize for a serious wrong?} You untie a knot – one knot from your cord, chosen by the wronged party. The knot is burned. The ash is scattered on the wind. If the wind blows toward the wronged party, you are forgiven.

\textbf{What is the most important festival?} The Tying of the New Cord (new year). Every Tulkani receives a new cord, unknotted. The old cord is burned. The ash is mixed with water and drunk. The drink tastes of memory.

\textbf{What is the secret that could destroy a powerful person?} The Tulkani are not wanderers. They are pilgrims – walking toward a destination that no longer exists. Their homeland was swallowed by the desert. The desert is still hungry. They will never arrive.
\end{almanac}

\subsection*{Etiquette Hooks}
\begin{itemize}
  \item \textbf{Story as Kinship}: Tell a true story (in character) to a Tulkani elder → gain a temporary String (Memory‑Bead).
  \item \textbf{Circle Courtesy}: Do not step over the cord that marks the wagon circle → first social roll in the circle gains +1 Position.
  \item \textbf{Ninth Night Offering}: Leave a knotted cord at the central post → the Hollow’s attention is delayed by one night.
  \item \textbf{Faux Pas (SB)}: Counting to nine aloud – the GM may start \textit{Hollow’s Patience} [4] immediately.
\end{itemize}

\subsection*{Strings (sample)}
\begin{itemize}
  \item \textbf{Memory‑Bead} – a single knot; present it to ask for a story from any Tulkani.
  \item \textbf{Safe‑Haven Oath} – a knotted cord; guarantees sanctuary in any Tulkani circle.
  \item \textbf{Dream‑Cord} – a silk cord with nine knots (the ninth loose); sleep on it to ask Ikasha one question in a dream.
\end{itemize}

\subsection*{Venue Tags}
\begin{itemize}
  \item \textbf{Wagon Circle}: neutral ground; no violence permitted; breaking the circle starts \textit{Exile Clock} [4].
  \item \textbf{Caravan Court}: disputes settled by storytelling; a good story (Performance test DV 3) grants +1 Effect on the verdict.
\end{itemize}

\newpage
\section{Tulkani --- ``The Itinerant Courts''}
\label{chap:tulkani}

\subsection*{Quick Reference}
\textbf{Tagline:} \textit{The knot remembers. The story pays. The circle holds.}

\begin{quote}
  \textbf{Caravan Judge (Elite):}  
  ``The road does not forgive a debt. It only reschedules it. A promise spoken under the open sky is witnessed by the wind, the grass, and the wheels of every cart that passes. I do not write judgments. I tie knots. The knots remember. When you are ready to pay, untie one. If you are not ready, tie another.''
\end{quote}

\begin{quote}
  \textbf{Wagon-Keeper (Commoner):}  
  ``My wagon is my hearth, my court, my grave. The Hollow cannot cross the circle of wheels unless invited. The tax collector cannot enter without a story. The lord cannot demand rent unless we camp on his land for nine nights. On the eighth night, we leave. That is the law.''
\end{quote}

\textbf{Genre / Mood:} Itinerant community, story-driven justice, folk horror at the edge of the Hollow.

\textbf{Starting Location:} A wagon circle at a crossroads, where nine painted wagons form a ring. A Caravan Judge sits on an overturned barrel, a knotted cord in her hand. ``You are standing on ground that belongs to the lord of the eastern valley. He will be here at dawn to demand rent. But you are in the circle. You are safe until the ninth hour. Tell me a story — a true one — and I will tell you how to leave before the lord arrives.''

\vspace{2mm}
\begin{tcolorbox}[colback=black!3,colframe=blue!60!black,title={Theme \& Atmosphere}]
The Tulkani have no fixed homeland; they are wanderers who follow the seasons and the trade routes. Their wagons are their hearths, their circles are their courts, and their knots are their ledgers. They are welcome everywhere and trusted nowhere — a position they wear with wry humour. A Tulkani caravan can appear at a lord's border, camp for seven nights, and leave before the eighth — because on the eighth night, the lord can demand rent. On the ninth night, the Hollow comes. The Tulkani never stay nine nights in one place. Their courts are circles of wagons; their judges are elders who carry braided cords. Each knot is a vow, a story, or a favour owed — but never merely a debt. They spread the cult of Ikasha — She Who Sleeps — not through temples, but through dream-cords tied in the dark. Their neutrality is sacred: a Tulkani camp is sanctuary for any traveller, regardless of allegiance, as long as they respect the circle and offer a story when asked.

\vspace{2mm}
\textbf{What you notice first:} The smell of woodsmoke and cooking oil. The painted wheels of wagons arranged in a circle, each wheel a different colour. The way everyone touches a knotted cord before speaking of an obligation.

\textbf{The one rule that kills newcomers:} Never sleep outside the circle at night. The wagon ring is sanctuary. Outside it, the Hollow walks. A Tulkani will not rescue you — they will only mourn you at dawn.
\end{tcolorbox}

\subsection*{Regional Mood: The Weight of the Knot}

Power here is measured in vow-cords, story‑credits, and the number of nights camped on a lord's land. The Tulkani have no king, no judges, no written law — only the Caravan Circle, where disputes are settled by storytelling. A liar is exposed not by evidence but by the knots on his cord: if a knot is tied in false witness, it will tighten until blood is drawn. Ikasha's cult is not a religion; it is a survival tactic — the dream inside the wagon, the whisper in the creaking wheel. The Tulkani pay no permanent taxes, but they pay tolls in stories. A lord who demands coin will be told a story so boring that his crops fail. A lord who offers hospitality will be told a story so beautiful that his harvest doubles. The Tulkani do not track debts; they track \textbf{reciprocity}. What you give, you get back — maybe tomorrow, maybe in the next valley. But the circle remembers.

\subsection*{Prominent NPCs of the Tulkani}
\label{sec:tulkani-npcs}

\begin{longtable}{|p{0.2\textwidth}|p{0.25\textwidth}|p{0.3\textwidth}|p{0.15\textwidth}|}
\hline
\textbf{Name} & \textbf{Role/Title} & \textbf{Motivation} & \textbf{Rumoured Quirk} \\
\hline
Elder Sarnai & Caravan Judge of the Nine Knots & Keep the circle whole & She ties a knot for every judgment; her cord drags on the ground. \\
\hline
Temur Wagon-Keeper & Master of the painted wheel & Find new grazing ground before the ninth night & He paints a new colour on his wagon each season; the colours are memories. \\
\hline
Gulshat & Story-Seller & Spread the dream of Ikasha through tales & She carries a dream-cord with nine knots; the ninth knot is loose. \\
\hline
The Knife (Unnamed) & Ikasha's agent & Free a spirit bound by the Veiled Sisters & He wears a mask made of knotted cord; he has no face underneath. \\
\hline
Vow-Broken & An exile who cut his own cord & Earn back his place in the circle & He carries a single knot tied in his own hair; the knot is his promise to return. \\
\hline
\end{longtable}

\subsection*{Names of the Tulkani}
\label{sec:tulkani-names}

Inspired by the nomadic traditions of the steppe and desert: short, guttural, with honorifics and clan prefixes. Surnames are often replaced by epithets describing a deed or a cord.

\begin{longtable}{|p{0.25\textwidth}|p{0.25\textwidth}|p{0.2\textwidth}|p{0.2\textwidth}|}
\hline
\textbf{First Names (Male)} & \textbf{First Names (Female)} & \textbf{Clan/Epithet} & \textbf{Knot Name} \\
\hline
Temur & Sarnai & \textit{the Wagon-Keeper} & \textit{Ninth Cord} \\
\hline
Batu & Gulshat & \textit{the Story-Teller} & \textit{Dream-Knot} \\
\hline
Orhan & Nergui & \textit{the Knife} & \textit{Unbound} \\
\hline
Cengiz & Altan & \textit{Vow-Broken} & \textit{Cut Cord} \\
\hline
Toghrul & Odval & \textit{Circle-Keeper} & \textit{Neutral Knot} \\
\hline
\end{longtable}

\paragraph*{(Circle/Wagon/Ground)} A wagon circle at dusk, lanterns lit; a painted wagon with a dream-cord hanging from the wheel; a crossroads where nine roads meet.

\section*{Spades --- Places}
\label{sec:tulkani-places}
\begin{enumerate}
\setcounter{enumi}{1}
\item \textbf{Wagon Circle (Night)} --- Nine wagons arranged in a ring, lanterns burning, wheels touching.
  \hfill \textit{Inside the circle, the Hollow cannot enter. Outside, it waits. The lanterns must never go out.} `[SANCTUARY]`
\item \textbf{Caravan Court (Day)} --- A carpet spread on grass; a judge on a barrel; a cord on the ground.
  \hfill \textit{The cord is the law. Step over it, and you accept the judge's verdict.} `[JUSTICE]`
\item \textbf{Dream-Wagon (Ikasha's Shrine)} --- A painted wagon with no wheels, sitting on stones. The door is a curtain of knots.
  \hfill \textit{Inside, a single candle and a bed of straw. Ikasha dreams there. Do not wake her.} `[MYSTIC]`
\item \textbf{Cord-Maker's Stall} --- A canvas awning where a woman braids horsehair and silk.
  \hfill \textit{Each cord is a promise. The colour tells the weight: white for truth, red for blood, black for a secret.} `[CRAFT]`
\item \textbf{Crossroads of the Ninth Road} --- A meeting of nine paths, only eight of which are visible.
  \hfill \textit{The ninth road is for the dead. The Tulkani leave an offering there each full moon.} `[THRESHOLD]`
\item \textbf{Abandoned Circle} --- A ring of burnt grass and wheel ruts. No wagons. No people.
  \hfill \textit{The circle was broken on the ninth night. The Hollow came. The grass will not grow back.} `[OMEN]`
\item \textbf{Vow-Cord Tree} --- An old oak with hundreds of knotted cords hanging from its branches.
  \hfill \textit{Each cord is a promise kept. The tree remembers. The wind unties one each season.} `[COMMUNITY]`
\item \textbf{Hollow's Gate (Seasonal)} --- A place where the boundary is thin, marked by a single painted stone.
  \hfill \textit{The Tulkani camp here only on the new moon. On the full moon, they are a day's ride away.} `[THRESHOLD]`
\item \textbf{Tinkanar (Itinerant Caravanserai)} --- A mobile market of wagons, open for one day each season.
  \hfill \textit{The location changes. The Tulkani announce it by dream — the dream of a sleeping merchant.} `[TRADE]`
\item[J] \textbf{Story-Well} --- A dry well where the Tulkani throw written stories on parchment.
  \hfill \textit{The stories are offerings to Ikasha. The well is full. The parchment never rots.} `[RITUAL]`
\item[Q] \textbf{The Ninth Cord's Post} --- A wooden post at the centre of every circle, with nine notches.
  \hfill \textit{The ninth notch is never carved. It is left for the Hollow. The Hollow never takes it.} `[SYMBOL]`
\item[K] \textbf{The Wandering Hearth} --- A wagon that is never hitched, always moving. It carries the council fire.
  \hfill \textit{The fire has been burning for three hundred years. It has never touched the ground.} `[LEGEND]`
\item[A] \textbf{Ikasha's Dreaming Place (Myth)} --- A cave beneath a hill, where She Who Sleeps lies on a bed of knotted cords.
  \hfill \textit{If you untie a cord, she wakes. If she wakes, the world ends.} `[MYTH]`
\end{enumerate}

\paragraph*{(Judge/Keeper/Storyteller)} Caravan Judge with a dragging cord; Wagon-Keeper painting a wheel; Story-Seller with a dream-cord and a cup of water.

\section*{Hearts --- People \& Factions}
\label{sec:tulkani-people}
\begin{enumerate}
\setcounter{enumi}{1}
\item \textbf{Caravan Judge} --- Old, grey-braided, a cord so long it drags on the ground.
  \hfill \textit{She ties a knot for every judgment. Untying a knot voids the judgment. She has never untied one.} `[JUSTICE]`
\item \textbf{Wagon-Keeper} --- A woman with paint on her hands, a brush behind her ear.
  \hfill \textit{She paints a new colour on her wagon each season. The colours are memories. She is running out of colours.} `[CRAFT]`
\item \textbf{Story-Seller (Gulshat)} --- A young woman with a dream-cord around her wrist. She speaks in rhymes.
  \hfill \textit{Her stories are true — but only if you listen with your eyes closed.} `[SOCIAL]` `[MYSTIC]`
\item \textbf{Ikasha's Knife} --- A hooded figure with a mask of knotted cord. No face underneath.
  \hfill \textit{He cuts the cords of oath-breakers. The cut cord becomes a ghost. The ghost follows the breaker.} `[JUSTICE]` `[DEATH]`
\item \textbf{Vow-Broken (Exile)} --- A man with a single knot tied in his own hair. He never smiles.
  \hfill \textit{He is saving stories to buy back his name. He needs nine. He has eight.} `[COMMUNITY]`
\item \textbf{Cord-Maker} --- A woman with braids of horsehair and silk between her fingers.
  \hfill \textit{She can tell the weight of a promise by the feel of the hair. Horsehair is heavy. Silk is light.} `[CRAFT]`
\item \textbf{Hollow-Watcher} --- A boy who sits at the edge of the circle at dusk, watching the dark.
  \hfill \textit{He can see the Hollow when it walks. He points. The adults throw salt.} `[SCOUT]`
\item \textbf{The Dreamer} --- A woman who sleeps in the Dream-Wagon. She never wakes during the day.
  \hfill \textit{She dreams the caravan's path. The wagon follows her dreams.} `[MYSTIC]`
\item \textbf{Tinkanar Merchant} --- A fat man with nine bells on his hat. He sells knots and stories.
  \hfill \textit{He accepts only promises as payment. His ledger is a cord. It has nine thousand knots.} `[TRADE]`
\item[J] \textbf{The Ghost of the Ninth Night} --- A translucent figure seen at abandoned circles.
  \hfill \textit{It was a Tulkani who stayed too long. It asks: ``Is it the ninth hour yet?''} `[DEATH]`
\item[Q] \textbf{The First Wagon-Keeper (Legend)} --- A skeleton driving a wagon of bones. The wheels are skulls.
  \hfill \textit{It appears when a circle is broken. It offers a ride. The ride leads to the Hollow.} `[LEGEND]`
\item[K] \textbf{Ikasha's Voice} --- A whisper heard in the Dream-Wagon at midnight. It speaks in riddles.
  \hfill \textit{``The ninth cord is not a cord,'' it says. ``It is the space between the knots.''} `[MYSTIC]`
\item[A] \textbf{The Unwitnessed} --- A Tulkani who never tied a single knot. No one remembers her name.
  \hfill \textit{She walks through circles unnoticed. She can break any vow because no one saw her make it.} `[MYTH]`
\end{enumerate}

\paragraph*{(Toll/Story/Hollow)} A lord demands rent for camping on his land — negotiate nights or tell a story; the Hollow walks on the ninth night — a character hears their own name whispered; a stranger recognises your vow-cord and asks for a tale.

\section*{Clubs --- Complications/Threats}
\label{sec:tulkani-complications}
\begin{enumerate}
\setcounter{enumi}{1}
\item \textbf{Recognition (Vow-Cord Seen)} --- Someone from your past sees the knots you carry. They ask for a story you owe. `[SOCIAL]` `[RECIPROCITY]`
  \hfill \textit{The story must be about a promise you kept — or broke. If you cannot tell one, they tie a new knot.}
\item \textbf{Lord's Rent Demand} --- The local lord arrives at the circle at dawn. ``You have been here seven nights. Pay the eighth.'' `[AUTHORITY]` `[NEGOTIATION]`
  \hfill \textit{The price is a story that flatters him. If the story fails, he seizes a wagon. The circle holds its breath.}
\item \textbf{Hollow Walks (Ninth Night)} --- Someone in the circle has miscounted. The ninth hour approaches. `[DEATH]` `[FEAR]`
  \hfill \textit{The lanterns flicker. A shadow stands at the edge. It whispers a name. The name is the last person who broke an oath.}
\item \textbf{Vow-Cord Tangled} --- Two promises conflict. The knots tighten. Blood is drawn. `[CONFLICT]` `[RECIPROCITY]`
  \hfill \textit{The judge must cut one cord. The severed cord becomes a ghost. The ghost accuses the cutter of favouritism.}
\item \textbf{Wagon Wheel Broken} --- A wheel cracks. The circle is incomplete. The Hollow can enter. `[DISASTER]`
  \hfill \textit{The wheel must be repaired before dusk. The repair requires a story — a true one — told to the wood. The wood listens.}
\item \textbf{Tulkani Schism} --- Two judges dispute the meaning of a knot. The circle splits. `[POLITICS]`
  \hfill \textit{Each judge claims authority. The party must choose which cord to follow. The other cord will remember — and so will the Hollow.}
\item \textbf{Ikasha's Dream Disturbed} --- A stranger enters the Dream-Wagon. The Dreamer wakes. `[MYSTIC]` `[CRISIS]`
  \hfill \textit{Ikasha stirs. The ground trembles. The stranger must tell a story so good that Ikasha falls back asleep. The stranger is you.}
\item \textbf{Story Toll (Bridge/Pass)} --- A guard demands a story to cross. Your story is not good enough. `[TRAVEL]` `[SOCIAL]`
  \hfill \textit{The guard demands a truth instead. The truth must be a secret. The secret must hurt. The guard will remember.}
\item \textbf{Vow-Broken's Return} --- An exile tries to rejoin the circle. The judge refuses. `[COMMUNITY]` `[DRAMA]`
  \hfill \textit{The exile demands the party speak for him. If they refuse, he cuts his own throat. If they speak, they owe a knot — and a story.}
\item[J] \textbf{The Ninth Cord Cut} --- Someone cuts the ninth knot on the central post. The Hollow floods in. `[DEATH]` `[TRAP]`
  \hfill \textit{The lanterns go out. The circle becomes a cage. The party must re-tie the knot in the dark. The cord is warm.}
\item[Q] \textbf{Ikasha's Knife Betrayed} --- The Knife is offered a bribe to cut a false cord. He refuses. The briber attacks. `[COMBAT]` `[JUSTICE]`
  \hfill \textit{The Knife is wounded. He asks the party to carry his mask. The mask weighs nothing — until you put it on. Then it weighs everything.}
\item[K] \textbf{Tinkanar Closed (No Stories)} --- The mobile market is empty. The merchant has run out of knots. `[TRADE]` `[CRISIS]`
  \hfill \textit{He is selling his own name. The price is a story that makes him weep. He has not wept in forty years. The last story he heard was a lie.}
\item[A] \textbf{Ikasha Wakes} --- The Dreamer dies. Ikasha opens her eyes. The sky turns white. `[APOCALYPSE]`
  \hfill \textit{She asks one question: ``Who cut the ninth cord?'' The answer must be a name. If no name is given, she takes all names.}
\end{enumerate}

\paragraph*{(Cord/Bead/Charm)} Vow-bead (a single knot, can be used to call in a favour once); safe-haven oath (a knotted cord that guarantees sanctuary in any Tulkani circle); dream-cord (gift from Ikasha, one question answered in a dream).

\section*{Diamonds --- Rewards/Leverage}
\label{sec:tulkani-rewards}
\begin{enumerate}
\setcounter{enumi}{1}
\item \textbf{Vow-Bead (Single Knot)} --- A small knot on a leather thong. Present it to call in a favour from any Tulkani. `[SOCIAL]`
  \hfill \textit{The knot will loosen after use. If it is not retied within a year, the favour becomes a story owed — and stories are harder to repay.}
\item \textbf{Safe-Haven Oath (Knotted Cord)} --- A braided cord with three knots. Hang it on your wagon to guarantee sanctuary in any Tulkani circle. `[TRAVEL]`
  \hfill \textit{The cord must be renewed each season. A new knot is added. After nine knots, it becomes a dream-cord — and the Hollow will notice you.}
\item \textbf{Dream-Cord (Ikasha's Gift)} --- A silk cord with nine knots, the ninth loose. Once, sleep on it to ask Ikasha one question. She will answer in a dream. `[MYSTIC]`
  \hfill \textit{The answer is always a riddle. The ninth knot tightens after use. When it tightens fully, Ikasha wakes — and the world changes.}
\item \textbf{Caravan Judge's Cord (Borrowed)} --- A length of braided horsehair. Tie it around your wrist to gain +1 Position in any negotiation with Tulkani. `[SOCIAL]`
  \hfill \textit{The cord is the judge's own. She will want it back. The return is a negotiation — and a story.}
\item \textbf{Wagon-Keeper's Paint (Pot)} --- A clay pot of blue pigment. Paint your own wagon wheel, and no lord will demand rent for nine days. `[TRAVEL]`
  \hfill \textit{The paint is made from crushed sky-stone. It glows under moonlight. The Hollow avoids it — but the lord will send a spy.}
\item \textbf{Story-Seller's Verse} --- A single verse of a true story. Recite it to calm a spirit or end a feud. `[SOCIAL]` `[MYSTIC]`
  \hfill \textit{The verse works once. After that, the story belongs to the listener. You cannot tell it again — but you can tell a different truth.}
\item \textbf{Ikasha's Mask (Knotted)} --- A mask of knotted cord. Wear it to become invisible to the Hollow for one night. `[STEALTH]` `[DEATH]`
  \hfill \textit{The mask has no eye holes. You must walk by memory. The memory may fail. If it fails, the Hollow sees you — and smiles.}
\item \textbf{Hollow-Watcher's Salt Pouch} --- A leather pouch of salt blessed by the boy who watches the circle. Throw it to banish a Debt-Walker. `[COMBAT]` `[DEATH]`
  \hfill \textit{The salt is bitter. The Debt-Walker will taste it. It will remember your face — and your smell.}
\item \textbf{Tinkanar Merchant's Bell} --- One of nine bells from the merchant's hat. Ring it to summon him anywhere. `[TRADE]`
  \hfill \textit{He will come. He will demand payment in stories. He will not leave until paid. He has a long memory.}
\item[J] \textbf{Vow-Broken's Hair Knot} --- A single knot tied in human hair. It contains a broken promise. Use it to void a contract. `[SOCIAL]` `[MAGIC]`
  \hfill \textit{The knot is from an exile. He wants it back. He will pay a story for it — the story of how he broke his vow.}
\item[Q] \textbf{The Ninth Cord's Post Fragment} --- A splinter of wood from the central post. Touch it to any knot to learn who tied it. `[MYSTIC]`
  \hfill \textit{The splinter is warm. It will burn your hand if the knot was tied in false witness. The truth leaves no scar.}
\item[K] \textbf{First Wagon-Keeper's Skull} --- A human skull painted with wheel patterns. Place it at the centre of a circle to keep the Hollow away for nine nights. `[DEATH]` `[SANCTUARY]`
  \hfill \textit{The skull whispers on the ninth night. It asks: ``Is it time?'' Say no. If you say yes, it will take your place.}
\item[A] \textbf{Ikasha's Dream (Bottle)} --- A glass bottle containing a sleeping breath. Break it to send a question to Ikasha. She will answer in a vision. `[MYSTIC]`
  \hfill \textit{The vision lasts one minute. The cost is one memory. The memory is forgotten forever. The memory is always a happy one.}
\end{enumerate}

\section*{Quick use notes}
\label{sec:tulkani-quick-use}
\begin{itemize}
\item Draw until you have all four suits: Spade = place, Heart = actor, Club = pressure, Diamond = leverage. Highest rank sets the main Clock (2–5 → 4, 6–10 → 6, J/Q/K → 8, A → 10).
\item Diamonds are codified outcomes (knots, cords, beads) that change position rather than call for a roll.
\item If any A appears, echo \textbf{knot-omens}—a cord that tightens without being pulled, a knot that unties itself, a wheel that stops turning at the ninth spoke.
\end{itemize}

\subsection*{Additional Features: Itinerant Mechanics}
\label{sec:tulkani-itinerant}

\subsubsection*{The Ninth Night Rule}
A Tulkani caravan may camp on a lord's land for seven nights without payment. On the eighth night, the lord may demand rent — typically a story, a promise, or a small service. On the ninth night, the Hollow comes. No Tulkani stays nine nights in one place. A party camping with Tulkani must leave by the eighth dusk or face the Hollow's attention.

\textbf{Mechanic:} Each night the party camps in a Tulkani circle, mark a segment on a \textbf{Ninth Night Clock [8]}. At 7 segments, the lord (or local authority) arrives to demand rent. At 8 segments, the Hollow manifests. The party can leave before 8 to reset the clock.

\subsubsection*{Story as Currency}
Tulkani do not use coin. They trade in stories, promises, and knots. A good story can buy a meal, a horse, or safe passage. A bad story can start a feud. The GM should encourage players to tell brief in-character stories; the quality (or a Presence + Sway roll) determines the value.

\textbf{Mechanic:} When a player tells a story, they roll Presence + Sway (DV set by the listener's mood). On a success, the story is worth 1 Measure of value. On a critical success, it is worth 2. On a failure, the listener is offended — the party gains -1 Position on their next social roll in that circle. A player may also offer a \textbf{truth} (a secret, a confession, a memory) instead of a story; a truth is always worth 2 Measures but marks 1 Obligation to Ikasha.

\subsubsection*{Land Rent and Neutrality}
Tulkani circles are neutral ground. No lord may collect taxes or enforce laws inside a circle, regardless of whose territory it sits on. However, the lord may demand rent at the circle's edge, and the Tulkani will negotiate. A party inside a Tulkani circle is under the Caravan Judge's protection — but if the party commits a crime, the judge may expel them to the Hollow.

\textbf{Mechanic:} When a lord demands rent, the Caravan Judge mediates. The party may offer a story (one Measure), a promise (a future favour, represented by a Reciprocity Clock [4]), or a service (a quest). Refusing to pay means the lord may attack the circle — which breaks the neutrality and summons the Hollow.

\subsubsection*{Vow-Cords and Reciprocity}
Every Tulkani carries a vow-cord — a braided strand of horsehair and silk, with knots tied for every promise made and kept. A Tulkani who breaks a promise must cut the corresponding knot; the cut cord becomes a ghost that follows the oath-breaker until the promise is fulfilled. Outsiders may borrow a Tulkani vow-cord; the knots then record their promises. But the Tulkani see this not as debt but as \textbf{reciprocity}: what goes around comes around, and the circle always balances.

\textbf{Mechanic:} A character who accepts a vow-cord gains +1 die to any roll made to keep a promise related to the cord. However, if they break that promise, they must cut a knot — mark 1 Corruption and gain a Condition: \emph{Followed by a Ghost} (-1 die to stealth until the promise is fulfilled). A character may tie a new knot to make a new promise; each knot beyond the first increases the penalty for breaking. But a character who fulfills nine knots may ask the Caravan Judge for a \textbf{Circle Boon} — a permanent +1 die on all social rolls with Tulkani.

\subsubsection*{Local Patrons (Syncretic Names)}
\begin{itemize}
  \item \textbf{The Unwitnessed Weaver} — Inaea. *The knot that binds gently, the cord that remembers mercy.*
  \item \textbf{The Cord-Cutter} — Isoka. *The clean break, the shedding of false promises.*
  \item \textbf{The Dreamer Below} — Ikasha. *She who sleeps in the wagon, she who dreams the road.*
  \item \textbf{The Ninth Traveller} — The Traveler. *The road that has no end, the threshold that is always one step away.*
\end{itemize}

\subsection*{Lore Echoes (Cross-Expansion Hooks)}
\begin{itemize}
  \item \textbf{Tulkani Vow-Cords and the Sidhi Freedman's Seal:} Sidhi merchants on the Brass Coast accept Tulkani vow-cords as proof of credit. A well-knotted cord can buy a ship. A cut cord can sink one. The Sidhi know that a Tulkani knot is worth more than a signature.
  \item \textbf{Ikasha and the Veiled Sisters (Ashaan):} The Tulkani are the primary spreaders of Ikasha's cult outside Ashaan. A Breath-less agent may be disguised as a Tulkani story-seller; a Tulkani caravan may be a mobile safe house. The Sisters have tried to infiltrate the circles three times. Each time, the spy was found with a cord around her throat — tied in a knot she could not untie.
  \item \textbf{The Ninth Night and the Hollow (Eastern Realms):} The Tulkani's ninth-night rule is the only reason the Hollow has not overrun the eastern trade routes. Aelaerem hedge-keepers sometimes hire Tulkani guides to reset the Hollow's attention clocks. The payment is always a story — and the story is always about a door that was left open.
  \item \textbf{Story Debts and the Ukwe (Sekogo):} The Tulkani believe that the Ukwe is the first vow-cord tree. They leave offerings of knotted grass at its roots. A Tulkani judge may consult the Ukwe to verify a story's truth. The Ukwe speaks in rustles. A single rustle means true. Two mean false. Three mean the story is not finished.
\end{itemize}

\subsection*{Tulkani Superstitions (burn for +1 die once)}
\begin{itemize}
  \item \textbf{Bread to the Circle:} Crumble a crust at the centre of the wagon ring. Ignore the first \emph{Recognition} penalty. `[SUPERSTITION]`
  \item \textbf{Knot the Lantern:} Tie a single knot in the rope of a lantern. Cancel \emph{Lord's Rent Demand} or take –1 die on your next negotiation (your choice). `[SUPERSTITION]`
  \item \textbf{Name to the Ninth Post:} Whisper a dead judge's name into the central post. Gain +1 Effect when telling a story this scene, but mark \emph{Obligation} to Ikasha. `[SUPERSTITION]`
\end{itemize}

\subsection*{Thematic SB Spend Table}
\label{sec:tulkani-sb}

\paragraph{Minor Complications (1 SB)}
\begin{itemize}
\item \textbf{Exposure:} Your actions draw unwanted attention from \textbf{Tulkani judges or local lords}.
\item \textbf{Noise:} Sounds of your actions alert nearby \textbf{circle-watchers or Hollow scouts}.
\item \textbf{Trace:} Evidence of your passage marks your route for \textbf{vow-breakers or story-hunters}.
\item \textbf{Delay:} A brief but meaningful setback costs you \textbf{time or favourable camping ground}.
\item \textbf{Supply Strain:} Mark +1 segment on a relevant \textbf{resource clock}.
\end{itemize}

\paragraph{Moderate Setbacks (2 SB)}
\begin{itemize}
\item \textbf{Alarm Raised:} \textbf{Local Caravan Judge or lord's herald} becomes aware and begins responding.
\item \textbf{Position Lost:} You lose advantageous ground/cover/stealth due to \textbf{lantern failure or broken wheel}.
\item \textbf{Foe Appears:} A \textbf{debt collector or Hollow manifestation} arrives on scene.
\item \textbf{Gear Trouble:} A piece of equipment becomes \textbf{Compromised/Neglected}.
\item \textbf{Lock/Barrier:} A simple obstacle now requires a test to overcome.
\end{itemize}

\paragraph{Serious Trouble (3 SB)}
\begin{itemize}
\item \textbf{Reinforcements:} Additional \textbf{Tulkani riders, lord's men, or Hollow spirits} arrive.
\item \textbf{Key Gear Breaks:} A crucial tool/weapon becomes temporarily unusable.
\item \textbf{Major Twist:} The situation fundamentally changes—\textbf{knot cut/wagon broken/ninth night triggered}.
\item \textbf{Rail Tick:} Advance a relevant campaign/front clock by 1 segment.
\item \textbf{Condition Applied:} Mark \textbf{Fatigue 1/Harm 1/Condition} appropriate to fiction.
\end{itemize}

\paragraph{Major Turns (4+ SB)}
\begin{itemize}
\item \textbf{Trap Springs:} A prepared danger activates with full effect.
\item \textbf{Authority Arrival:} \textbf{High Caravan Judge, Ikasha's Knife, or Hollow's Debt-Walker} intervenes.
\item \textbf{Scene Shift:} The environment changes dramatically—\textbf{circle breaks/wagon burns/lanterns go out}.
\item \textbf{Patron Omen:} Divine/arcane forces take notice—\textbf{omen appears/blessing lost/curse manifests}.
\item \textbf{Narrative Pivot:} The story takes an unexpected turn that reframes objectives.
\end{itemize}

\paragraph{Region-Specific SB Options}
\begin{itemize}
\item \textbf{Tulkani (Knot Law):} Cords tighten without warning, a knot unties itself, a wheel stops at the ninth spoke.
\item \textbf{Tulkani (Story Debt):} A story you told returns to you changed, a listener laughs at a tragedy, the wind carries your words to someone who should not hear.
\item \textbf{Tulkani (Ninth Night):} A lantern flickers, a shadow stands at the circle's edge, a whisper says your name from the dark.
\end{itemize}

\subsection*{Pursuit Ladder (Near / Mid / Far)}
\begin{itemize}
\item \textbf{Advance/Withdraw (Wits + Stealth):} On a hit, shift one band. On a strong hit, also impose \emph{False Knot}.
\item \textbf{Cut the Cord (Presence + Sway):} On a hit, freeze range for one exchange; if a vow-bead is in your possession, +1 die.
\item \textbf{Weather the Hollow (Resolve + Spirit):} On a hit in \emph{Hollow Walks}, you pick who hears the whisper; on a miss, the GM does.
\end{itemize}

\subsection*{Boss Seeds (Knot \& Dream)}
\begin{itemize}
\item \textbf{The Cord-Cutter (Nine Broken Promises)} — \emph{Phases}: Vow-Cords Severed $\rightarrow$ Orchestrated Exile $\rightarrow$ Circle Coup. \emph{Cracks}: re-tie a cut cord, expose a false judge, survive a night outside the circle.
\item \textbf{The Dream-Eater (Ikasha's Shadow)} — \emph{Phases}: Dreamers Waking $\rightarrow$ Procession of Empty Wagons $\rightarrow$ Knot-Baptism of a Storyteller. \emph{Cracks}: relic dream-cord, survivor testimony, the ninth cord refusing to tighten.
\end{itemize}

\begin{tcolorbox}[colback=black!3,colframe=black!40!white,title={Patronage \& Power}]
Among the Tulkani, power flows through knots, stories, and the memory of the circle. A Caravan Judge's authority is recognised by every lord from the steppe to the coast — not because of armies, but because the Hollow respects the judge's cord. The Tulkani have no temples, no kings, no borders. They have the ninth night, and that is enough.

\textbf{For the GM:}  
Power in Tulkani society is measured in knots tied, stories told, and nights counted. Patronage often takes the form of vow-beads, dream-cords, or safe-haven oaths — each carrying the weight of a promise. To emphasize the itinerant mechanics:
\begin{itemize}
\item Require players to \textbf{negotiate land rent} with local lords, using stories or promises as currency. A failed negotiation may lead to the Hollow.
\item Use the \textbf{Ninth Night Clock} as a tension mechanic: every night in a Tulkani camp advances the clock. At 8, the Hollow manifests unless the party leaves.
\item Let \textbf{vow-cords} be physical props — actual knotted strings that players can hold, tie, or cut. A cut cord is a dramatic moment.
\item Emphasise that \textbf{storytelling is a skill}: reward creative players with mechanical bonuses, and let failed stories have consequences — a bored lord may raise rent, an offended spirit may curse.
\item Treat Tulkani society as a \textbf{reciprocity economy}, not a debt economy. What you give, you get back — but the circle decides when.
\end{itemize}
In Tulkani society, the knot remembers everything — and the ninth night forgets nothing.
\end{tcolorbox}

% Index entries
\index{Tulkani|see{Itinerant Courts}}
\index{Theme generator}
\index{Itinerant Courts generator}
\index{Places!Tulkani}
\index{People!Tulkani}
\index{Complications!Tulkani}
\index{Rewards!Tulkani}
\index{Ninth night}
\index{Vow-cords}
\index{Story as currency}
\index{Ikasha}
\index{Dream-cord}
\index{Caravan Judge}
\index{Wagon circle}
\index{Hollow}
\index{Tinkanar}
\index{Local Patrons!Tulkani}
\newpage

% ----------------------------------------------------------------------
% APPENDIX: CARAVAN & TRIBE RULES
% ----------------------------------------------------------------------
\appendix
\chapter{Caravan \& Tribe Rules}

This appendix distils the travel‑and‑trade engine from the \textit{Sands of Moon and Brass} expansion. Use these procedures when running campaigns focused on the Way of Silk.

\section*{Caravan / Tribe Sheet}
\begin{longtable}{|p{0.3\textwidth}|p{0.55\textwidth}|}
\hline
\textbf{Caravan/Tribe Name} & \\
\hline
\textbf{Seat (home oasis/camp)} & \\
\hline
\textbf{Council Affiliation} & (Fhara, Sidhi, Pereshi, Kuvani, Tulkani – or independent) \\
\hline
\textbf{Strings (2–3)} & water‑share tablet ∙ lamp‑oath seal ∙ trade charter ∙ feud settlement ∙ way‑song ∙ guide‑stone \\
\hline
\textbf{Tags (2–4)} & Water‑Wise ∙ Star‑Pathed ∙ Heat‑Trained ∙ Merchant‑Known ∙ Hospitable ∙ Law‑Strict ∙ Poet‑Loud ∙ Falcon‑Banner \\
\hline
\textbf{Tracks} & Standing [6] (reputation), Feud [4] (active quarrel), Exposure [6] (to foreign powers) \\
\hline
\textbf{Oath Ledger} & favours owed/held (name + what; acts like Favour/Leverage) \\
\hline
\textbf{Notables} & guide ∙ war‑leader ∙ speaker ∙ poet ∙ water‑keeper \\
\hline
\end{longtable}

\subsection*{Standing [6]}
Tick up for kept oaths, fair trade, valour. Tick down for oath‑breaking, cowardice, sacrilege. High Standing grants \textit{Audience: Respectful}; low invites \textit{Audience: Skeptical}.

\subsection*{Feud [4]}
Name the counterparty. Ticks on insult, theft, or harm. Clear by blood‑silver, ordeal, or deed. When filled, blood‑price owed or open war.

\subsection*{Exposure [6]}
Ticks when the party attracts attention from foreign powers (Admiralty, Veiled Sisters, etc.). When filled, a hostile faction acts directly.

\section*{Mount / Camel Sheet}
Pick a mount and 3–4 tags.

\paragraph{Dromedary (caravan)}: water‑efficient, steady pace, cargo capacity.
\paragraph{Charger (war)}: fast sprint, agile turns, combat‑trained.
\paragraph{Sand Runner (scout)}: light, narrow feet, collapsible saddle for dunes.

\subsection*{Mount Tags (choose 3–4)}
\begin{itemize}
  \item \textbf{Sand‑Sure}: Desert sand Drive/Handle DV –1; dune ridges, quicksand.
  \item \textbf{Water‑Efficient}: Half water consumption; ignore first Fatigue from heat.
  \item \textbf{Night‑Runner}: In darkness, Navigate DV –1; Ambush at Dawn starts one Position higher.
  \item \textbf{Heat‑Ready}: In sun/fire scenes, first Condition +1 is ignored.
  \item \textbf{Qanat‑Swimmer}: Underground water passages; Disengage gains Position +1 in tunnels.
  \item \textbf{Falcon‑Eye}: +1 die on Lookout/Watches; spot distant smoke/sand columns.
  \item \textbf{Brass‑Bearer}: Carries heavy loads; ignore encumbrance penalties for trade goods.
  \item \textbf{Poet’s Drum}: Once/score gain \textit{Audience: Fierce} aboard; can rally morale.
  \item \textbf{Guide‑Stones}: Hidden cairn‑marks; once/score Navigate Position +1 in trackless sand.
\end{itemize}

\section*{Roles (choose at table)}
Guide (navigation), War‑Leader (combat/assault), Speaker (parley), Poet (morale/saga), Water‑Keeper (supplies/repairs), Lookout (weather/sandstorm).

\section*{Season Wheel}
\begin{itemize}
  \item \textbf{Spring}: Caravan season – trade routes open.
  \item \textbf{Summer}: Heat – survival focus; oases contested.
  \item \textbf{Autumn}: Trade/Settle – markets busy; councils meet.
  \item \textbf{Winter}: Council – disputes settled; oaths renewed.
\end{itemize}
At each transition: advance Politics (Mandate/Crisis) for regional powers; roll Weather/Heat [4–6] for the current theatre; offer Oath Opportunities (escort, feud settlement, shrine warding).

\section*{Score Types}
Pick approach (Deceit ∙ Speed ∙ Shock ∙ Parley) → set Position/DV from venue/tags → roll.

\subsection*{Desert Caravan (Profit / Passage)}
Clocks: Alarm [4–6], Trade [6–8], Mount Damage [4], Blood‑Price [4].  
Entry: sand gate, hidden wadi, qanat shortcut, market ruse.  
On 1s: GM spends SB → sandstorm warning, hidden quicksand, rival caravan ambush, guide betrayal.  
Resolution: When Trade fills, choose coin or Strings. If Blood‑Price fills, mark Feud +1 or pay blood‑silver.

\subsection*{Oasis Siege (Water / Tribute)}
Clocks: Thirst [4], Defense [6], Sentries [4], Terms [6].  
Position tweaks: Water‑Efficient/Heat‑Ready help; sandstorms/hail hurt.

\subsection*{Council Moot (Law / Settlement)}
Venue: circle of stones, winter pavilion, qanat meeting hall.  
Moves: Oath‑Swear (commit under penalty), Wager Wyrd (ordeal by feat), Blood‑Silver (compensation roll), Witness the Verse (Poet stakes a truth).  
Outcomes: resolve Feud, write Oath to ledger (durable String), assign seasonal water rights or trade route.

\section*{Weather / Heat Matrix}
Advance on 1s or fiction:
\begin{itemize}
  \item \textbf{Khamsin/Sandstorm}: sight Position –1; Sand‑Sure cancels. On 1, Pursuit +1 (lost bearings).
  \item \textbf{Flash Flood}: ranged actions –1 die; on 1, Condition +1 (saddle/gear).
  \item \textbf{Ward Storm}: Navigate DV +1; on 1, choose Delay (Distance stalls) or Sand‑Blast (Condition +1).
\end{itemize}

\section*{Generators (Quick Reference)}
\subsection*{Desert Targets (d66)}
11 star‑temple ∙ 12 beacon tower ∙ 13 cliff monastery ∙ 14 salt works ∙ 15 oasis fleet ∙ 16 lord’s pavilion ∙ 21 dune storehouse ∙ … (full table in \textit{Sands of Moon and Brass})

\subsection*{Desert Hazards (d12)}
1 dune on a bend ∙ 2 sudden flash flood ∙ 3 quicksand pool ∙ 4 deadwood snag ∙ 5 scorpion nest ∙ 6 heat shimmer ∙ 7 hidden side‑cut ∙ 8 rock echo ∙ 9 toll chain half‑raised ∙ 10 reeds conceal archers ∙ 11 rain‑swollen ford ∙ 12 sand‑suck bank.

\subsection*{Council Cases (d12)}
1 insult in verse ∙ 2 stolen guide stone ∙ 3 broken oath on winter grain ∙ 4 blood‑price disputed ∙ 5 marriage claim ∙ 6 way‑song stolen ∙ 7 hostage pledge lapsed ∙ 8 salvage rights ∙ 9 border cairn moved ∙ 10 feud cooling terms ∙ 11 verse witness contest ∙ 12 mercenary pay withheld.

\subsection*{Verses \& Boons (d12)}
1 drum‑song that steadies hands ∙ 2 falcon omen at dawn ∙ 3 amber find ∙ 4 guide’s ghost shows a cut ∙ 5 sand‑storm spares you ∙ 6 omen of red sails ∙ 7 poet’s verse spreads ∙ 8 winter pavilion adopts you ∙ 9 qanat seal renewed ∙ 10 cliff bell silent ∙ 11 stars open path ∙ 12 oath‑ring warms (truth told).

% ----------------------------------------------------------------------
% EPILOGUE
% ----------------------------------------------------------------------
\chapter*{Epilogue: The Ninth Knot}
\addcontentsline{toc}{chapter}{Epilogue}

My grandmother’s cord had nine hundred knots. I have added nine – one for each story in this book. The cord is still heavy. The debt is not paid. But I have tied the ninth knot loosely, as the Tulkani do. The Hollow is patient. So am I.

On the road to the Crimson Valley, with a fresh notebook and a borrowed lamp. The ink is squid. The wind is north. The next tale is already waiting.

\vspace{1em}
\noindent\textit{— Sarnai, Knot‑Tyer’s Hand}

\vspace{2em}
\begin{quote}
  \small
  \textit{“The road is a ledger. Every oasis is a credit, every broken wheel a debt. The Tulkani tie knots to remember who owes whom. The Fhara weigh water before they weigh silver. The Pereshi recite verses to keep the bridges open. And the Sidhi keep lighthouse lamps burning for the dead who still have not reached the shore.”}
  \\ — Dusana of the Raven Road (from whom I learned the telling)
\end{quote}

% ----------------------------------------------------------------------
% COLOPHON
% ----------------------------------------------------------------------
\chapter*{Colophon}
\addcontentsline{toc}{chapter}{Colophon}

This book was written on paper that had been a cargo manifest, a love letter, and a bandage. The ink is lampblack cut with salt water and the ash of a burnt confession. The binding is grey cord tied with nine knots. The ninth knot is not cut.

Sarnai composed the preface in a Tulkani wagon circle, while a Oath‑Keeper counted her breaths. The almanacs were gathered from traders, spies, and the occasional honest farmer. The generator tables were tested in the field – meaning she ran them at gunpoint in a caravan court while a Fhara well‑judge weighed her words.

If you find yourself on the Way of Silk and a Tulkani offers you a cup of tea, accept it. If they offer you a knotted cord, think carefully before you untie it.

The road is a cord. Every knot is a promise. Untie it, and you pay the debt.

\end{document}
